% =========================================================================
% DOCUMENT CLASS & SETUP
% =========================================================================

% \documentclass[sigconf]{acmart}       % For final camera-ready peer-reviewed submissions
\documentclass[sigconf,nonacm]{acmart} % For arXiv/preprint: removes copyright blocks & ACM footers

% =========================================================================
% GLOBAL PARAGRAPH AND PROSE RULES
% =========================================================================
% 1. Target Length: Aim for 3–7 sentences per paragraph. Avoid single-sentence 
%    paragraphs and dense walls of text.
% 2. Functional Focus: Each paragraph must serve a singular rhetorical function 
%    (e.g., context, problem, method, finding, interpretation, implication, 
%    limitation, or future work).
% 3. Explicit Signposting: Maintain visible logical flow across paragraph 
%    boundaries using structural transitions (e.g., "However", "As a result", 
%    "Consequently", "In contrast").
% 4. Micro-Sections: Use \myparagraph{...} blocks to signify named mini-paragraphs 
%    with specialized micro-rhetorical goals (e.g., "Contributions", "Significance").

% =========================================================================
% CONDITIONAL DEBUGGING FLAGS
% =========================================================================

\newif\ifDEBUG  % Declares a boolean switch called "DEBUG" for conditional compilation
\DEBUGtrue      % Set DEBUG to true (activates draft/todo notes)
\DEBUGfalse     % Set DEBUG to false (silences draft notes for clean compilation)
% NET EFFECT: \DEBUGfalse runs last, so debug is OFF by default. While OFF,
% \TODO{}, \CITE{}, and \NMS{} (co-author notes) render as nothing.
% To see draft markers, comment out the \DEBUGfalse line above.

% =========================================================================
% MODULAR PREAMBLE IMPORTS
% =========================================================================

% =========================================================================
% misc/data.tex
% Global Data Macro Definitions
% =========================================================================
        % Imports global stats, variables, and data macros
% =========================================================================
% core typography & layout engine
% =========================================================================

% \usepackage{amssymb,amsmath,amsfonts} % Optional math packages (currently commented)
\usepackage[T1]{fontenc}               % Uses modern 8-bit font encoding for crisp glyph rendering
\usepackage{balance}                   % Equalizes the length of columns on the final page of a 2-column layout

% =========================================================================
% code listing & custom graphical callout containers
% =========================================================================

\usepackage[most]{tcolorbox}           % Advanced custom framed boxes with advanced skins/listings hooks
\usepackage{listings}                  % Native source code block layout engine
\usepackage{xcolor}                    % Comprehensive driver-independent color naming & mixing support

% =========================================================================
% standard acm-approved helper packages
% =========================================================================

% --- structural & page geometry control
\usepackage{afterpage}                 % Postpones execution of commands until the next page break
\usepackage{float}                     % Enhances figure place control; introduces the rigid [H] specifier
\usepackage{lscape}                    % Places specific content pages in landscape orientation
\usepackage{rotating}                  % Rotates text, tables, and figures along arbitrary angle limits

% --- mathematics & computational logic
\usepackage{algorithmic}               % Structural pseudocode layout environment builder

% --- advanced tabular structures
\usepackage{array}                     % Extended configuration controls for table column definitions
\usepackage{booktabs}                  % Renders publication-grade professional horizontal table rules
\usepackage{makecell}                  % Multiline table cell formatting with manual line breaks
\usepackage{multirow}                  % Enables cells spanning multiple horizontal rows
\usepackage{tabularx}                  % Auto-calculates column widths to fit specified horizontal constraints
\usepackage{longtable}                 % Handles multi-page tables natively within the standard stream

% --- captions & bibliographic citations
\usepackage[skip=4pt]{caption}         % Uniform control over caption spacing across figures and tables
\usepackage[compress]{cite}            % Automatically sorts and collapses numeric sequence reference tags
\usepackage{natbib}[numbers]           % Explicit configuration for structured numeric reference citation types

% --- document hyperlinking & reference automation
\usepackage{hyperref}                  % Injects interactive hyperlinks and global document indexes
\usepackage{cleveref}                  % Smart inline references (adds contextual prefixes like "Section" or "Fig.")

% --- list layout control
\usepackage{enumitem}                  % Fine-grained layout control over itemize, enumerate, and description lists

% --- graphic assets, icons, & utility packages
% \usepackage{filecontents}            % Pass-through payload file management (currently commented)
\usepackage{fontawesome5}              % Vector icon engine for system design and repository markers
\usepackage{graphicx}                  % Core graphics scaling, cropping, and rendering extension
\usepackage{lipsum}                    % Generates dummy text placeholders for layout prototyping
\usepackage{siunitx}                   % Uniform typesetting formatting engine for technical units and data measurements
\usepackage{soul}                      % Provides tracking, kerning, and strike-through/text-highlighting features
\usepackage{textcomp}                  % Provides specialized miscellaneous symbols and text glyphs
\usepackage[normalem]{ulem}            % Smooth underlining control; [normalem] prevents breaking \emph highlights
\usepackage{xspace}                    % Intelligent trailing space macro insertion engine

% =========================================================================
% page design & heading spacing tweaks
% =========================================================================

% \usepackage[expectsections, unnumberedsections, fencedCode]{markdown} % Direct Markdown parsing engine (commented)
% \usepackage[vskip=1em,font=itshape,leftmargin=2em,rightmargin=2em]{quoting} % Advanced blockquote layout engine (commented)

% Compacted section structural layout overrides (currently commented)
% \titlespacing*{\section}{0pt}{6pt}{4pt}
% \titlespacing*{\subsection}{0pt}{5pt}{3pt}
% \titlespacing*{\subsubsection}{0pt}{4pt}{2pt}

% =========================================================================
% CUSTOM UTILITY MACROS
% =========================================================================

% --- standard abbreviated latin syntax styling
\newcommand{\ie}{\textit{i.e.},\xspace}  % Uniformly formats "id est" with correct italicization and spacing
\newcommand{\eg}{\textit{e.g.},\xspace}  % Uniformly formats "exempli gratia" with correct italicization and spacing

% --- layout constructs
\newcommand{\myparagraph}[1]{\vspace{0.20cm}\noindent\textbf{#1} \noindent{}} % Spaced bold headers for inline subsections

% --- generative framework / prompts callout containers
% =========================================================================
% LLM PROMPT MACRO TEMPLATE
% Purpose: Formats Large Language Model system instructions, context windows,
%          or user prompts within a distinct, clean visual container.
%
% Usage: \myprompt{Argument #1: Box Heading}{Argument #2: Prompt Text Payload}
% Example: \myprompt{System Prompt}{Identify software reuse patterns...}
% =========================================================================

\newcommand{\myprompt}[2]{%
  \begin{tcolorbox}[
    enhanced,              % Activates advanced skin engine hooks (requires tcolorbox library 'skins')
    colback=gray!10,       % Light gray canvas fill background (10% opacity black over white canvas)
    colframe=black!20,     % Muted gray border line frame color (20% opacity black)
    boxrule=0.4pt,         % Ultra-thin, crisp frame line boundary thickness
    arc=1mm,               % Sets a subtle, modern 1mm corner radius curve
    % Tight geometric padding rules to prevent large text blocks from forcing column overflows:
    left=2mm,              % Internal text padding buffer from the left border
    right=2mm,             % Internal text padding buffer from the right border
    top=1mm,               % Internal text padding buffer beneath the title separator rule
    bottom=1mm,            % Internal text padding buffer above the lower frame floor
    fonttitle=\bfseries,   % Sets the box title header typography font to boldface
    coltitle=black,        % Keeps the title string text plain black for contrast over gray background
    title={#1},            % Evaluates Positional Argument #1 as the structural header text
    % Elastic vertical paragraph spacing rules that stretch or shrink to prevent layout underfull breaks:
    before skip=3pt plus 2pt minus 1pt,
    after skip=3pt plus 2pt minus 1pt
  ]
  \small #2                % Reduces text size inside the box to \small for better scannability of dense text payloads
  \end{tcolorbox}%         % The trailing '%' character strips accidental trailing whitespaces from compilation output
}

% --- structured empirical finding blocks
\newcounter{finding}                   % Declares custom structural metric serialization register
\newcommand{\newfinding}{%             % Instantiates inline empirical milestone callouts
    \refstepcounter{finding}           % Increments register metrics sequentially and binds anchor links for \label
    \textbf{Finding \thefinding.}   % Prints stylized numeric marker index into current structural block stream
}

% =========================================================================
% REVIEWER INTERACTION & COLLABORATIVE DEBUG SWITCHES
% =========================================================================

\ifDEBUG
    % --- active draft editing marker functions
    \newcommand{\TODO}[1]{\hl{TODO: #1}} % Visual yellow structural highlight marker
    \newcommand{\CITE}[1]{\hl{CITE: #1}} % Visual yellow missing source tracking marker
    
    % --- colored co-author editorial discussion nodes
    \newcommand{\NMS}[1]{\textcolor{orange}{[NMS: #1]}} % Nicholas M. Synovic
\else
    % --- production compilation overrides (silences markers entirely)
    \newcommand{\TODO}[1]{}
    \newcommand{\CITE}[1]{}
    
    \newcommand{\NMS}[1]{}
\fi % Imports formatting overrides, extra packages, and layout tweaks

% =========================================================================
% CUSTOM LOGO MACROS
% =========================================================================

% Delays definition until document start to prevent package conflicts
\AtBeginDocument{%
    % Safely defines \BibTeX logo if not already defined by a loaded package
    \providecommand\BibTeX{{%
        \normalfont B\kern-0.5em{%   % Backspaces slightly after 'B'
            \scshape                 % Switches to small caps for 'ib'
            i\kern-0.25em b}\kern-0.8em\TeX}}} 
% Micro-spaces 'ib' and pulls 'TeX' close

% =========================================================================
% RIGHTS & COPYRIGHT MANAGEMENT (Suppressed while 'nonacm' is active)
% =========================================================================

\setcopyright{acmlicensed}          % Specifies the ACM licensing agreement type
\copyrightyear{2026}                % Sets the explicit copyright year
\acmYear{2026}                      % Sets the publishing year index
\acmDOI{XXXXXXX.XXXXXXX}            % Placeholder for the Digital Object Identifier string
\acmISBN{978-1-4503-XXXX-X/18/06}   % Placeholder for the standard publication ISBN

% Defines conference acronym, full metadata description, dates, and geographic location
\acmConference[Conference acronym 'XX]{Make sure to enter the correct
  conference title from your rights confirmation email}{June 03--05,
  2018}{Woodstock, NY}

% =========================================================================
% DOCUMENT CONTENT START
% =========================================================================

\begin{document}

\title{\textbf{INSERT TITLE}} % Formats the core paper title in bold face

% -------------------------------------------------------------------------
% AUTHOR METADATA
% -------------------------------------------------------------------------

\author{Nicholas M. Synovic}
\email{nsynovic@luc.edu}
\orcid{0000-0003-0413-4594} % Unique digital author identifier
\affiliation{
  \institution{Loyola University Chicago}
  \streetaddress{1032 W Sheridan Rd}
  \city{Chicago}
  \state{Illinois}
  \country{USA}
  \postcode{60660}
}

% Replaces the full author list in page running headers with a concise version to save space
\renewcommand{\shortauthors}{Synovic}

% -------------------------------------------------------------------------
% ABSTRACT SECTION
% -------------------------------------------------------------------------

% ABSTRACT: one self-contained, three-part narrative block (6–8 sentences total).
% HARD RULES: no citations and no raw statistical tables/numbers in the abstract.
\begin{abstract}
% Part 1 — Context & Problem (Sentences 1–3):
% Establish the broad SE domain and why it matters, then pivot to the core
% friction point or consequence. Function: context → problem.
Lorem ipsum dolor sit amet consectetur adipiscing elit. 
Quisque faucibus ex sapien vitae pellentesque sem placerat. 
In id cursus mi pretium tellus duis convallis. 

% Part 2 — Approach & Scope (Sentences 4–5):
% State "In this paper, we…" and summarize the design/method at one-sentence
% granularity, including evaluation scope (data size, system types, participants).
Tempus leo eu aenean sed diam urna tempor. 
Pulvinar vivamus fringilla lacus nec metus bibendum egestas. 

% Part 3 — Impact & Insights (Sentences 6–8):
% Deliver the headline qualitative/quantitative outcome, the supporting insight,
% and one concrete implication for practice or future research.
% Keep it high-level: no detailed numbers or statistics.
Iaculis massa nisl malesuada lacinia integer nunc posuere. 
Ut hendrerit semper vel class aptent taciti sociosqu. 
Ad litora torquent per conubia nostra inceptos himenaeos.
\end{abstract}

% =========================================================================
% ACM COMPUTING CLASSIFICATION SYSTEM (https://dl.acm.org/ccs)
% =========================================================================
\begin{CCSXML}
\end{CCSXML}

% The \ccsdesc macros render XML classifications into the document footer
\ccsdesc[500]{Software and its engineering}

% =========================================================================
% KEYWORDS INDEXING
% =========================================================================

\keywords{Replace me}

\maketitle

% =========================================================================
% SECTION: INTRODUCTION
% Function: Sell the paper. Six-paragraph arc; each block has ONE rhetorical job.
% Arc (one block each): broad context/hook → narrowed gap → goal & core insight
%   → contributions (bulleted) → significance/scope → roadmap.
% Rules: 3–7 sentences per paragraph; explicit transitions between blocks
%   ("However", "As a result"); no unquantified hype ("novel", "significant").
% =========================================================================
% RULE: every figure must be cited in-text via \Cref{fig:...} (e.g. \Cref{fig:template_figure}).
% =========================================================================
% TEMPLATE FIGURE PLACEHOLDER
% Uses a standard boxed paragraph container to stand in for missing visual assets.
% =========================================================================

\begin{figure}[h] % [h] requests placement "here" inline with the surrounding text stream
    \centering    % Horizontally centers the figure artifact within the column/page margins
    
    % \fbox draws a visible frame border around its inner content block
    \fbox{
        % \parbox structure creates an explicit, fixed-size paragraph boundary container
        % Arguments format: [inner-alignment][height][content-alignment]{width}
        \parbox[c][2in][c]{3in}{
            \centering Missing Figure % Centered placeholder text string inside the bounding box
        }
    }
    
    % Core descriptive label rendered immediately beneath the figure box frame
    \caption{Template figure placeholder}
    
    % Unique cross-referencing anchor tag; reference this asset elsewhere via \Cref{fig:template_figure}
    \label{fig:template_figure}
\end{figure}
\section{Introduction}
\label{sec:introduction}

% --- PARAGRAPH 1: CONTEXT AND HOOK ---
% Function: Intro paragraph 1.
% Target: 3–6 sentences.
% Content: Establish the broad software engineering domain, why it matters (reliability, cost, UX, scientific impact), and a hook that shows the domain is non‑trivial. No detailed technicalities here.
Lorem ipsum dolor sit amet consectetur adipiscing elit.
Quisque faucibus ex sapien vitae pellentesque sem placerat. 
In id cursus mi pretium tellus duis convallis.
Tempus leo eu aenean sed diam urna tempor. 
Pulvinar vivamus fringilla lacus nec metus bibendum egestas.
Iaculis massa nisl malesuada lacinia integer nunc posuere. 

% --- PARAGRAPH 2: NARROWED PROBLEM AND GAP ---
% Function: Intro paragraph 2.
% Target: 4–7 sentences.
% Content: Summarize current practice/prior work, then pivot with a clear ‘However/Yet/Despite this’ sentence that states the precise gap or friction point. End by emphasizing why this gap is consequential.
Lorem ipsum dolor sit amet consectetur adipiscing elit. 
Quisque faucibus ex sapien vitae pellentesque sem placerat.
In id cursus mi pretium tellus duis convallis. 
Tempus leo eu aenean sed diam urna tempor.
Pulvinar vivamus fringilla lacus nec metus bibendum egestas. 
Iaculis massa nisl malesuada lacinia integer nunc posuere.
Ut hendrerit semper vel class aptent taciti sociosqu. 

% --- PARAGRAPH 3: GOAL AND APPROACH (THE CORE INSIGHT) ---
% Function: Intro paragraph 3.
% Target: 3–6 sentences.
% Content: Present the overall research goal and high‑level approach (system, study, framework). Include plain‑language intuition about why your approach should address the stated gap.
Lorem ipsum dolor sit amet consectetur adipiscing elit. 
Quisque faucibus ex sapien vitae pellentesque sem placerat.
In id cursus mi pretium tellus duis convallis. 
Tempus leo eu aenean sed diam urna tempor.
Pulvinar vivamus fringilla lacus nec metus bibendum egestas. 
Iaculis massa nisl malesuada lacinia integer nunc posuere.

% --- PARAGRAPH 4: CONTRIBUTIONS ---
% Structure: Use 3–4 bullets, each describing one type of contribution (artifact, empirical study, insight, open science). 
% Formatting Rule: Keep each item one sentence plus a short clarifying clause, avoiding vague phrases (‘novel’, ‘significant’) without specifics.
\myparagraph{Contributions}
This paper makes the following core contributions:
\begin{itemize}[leftmargin=16pt, rightmargin=5pt, topsep=0pt]
    % Contribution 1: The Artifact (Tool / Framework / Dataset)
    \item Lorem ipsum dolor sit amet consectetur adipiscing elit.
    % Contribution 2: The Empirical Scale & Matrix Analysis
    \item Lorem ipsum dolor sit amet consectetur adipiscing elit.
    % Contribution 3: The Findings, Open Science & Replication Package
    \item Lorem ipsum dolor sit amet consectetur adipiscing elit.
\end{itemize}

% --- PARAGRAPH 5: SIGNIFICANCE & SCOPE ---
% Function: Intro paragraph 5 — "headline results + scope" preview (required).
% Target: 3–5 sentences.
% Content: Provide 1–2 headline quantitative or qualitative outcomes, plus a brief statement of scope/assumptions (where it works, where it might not). Do not repeat details from the Results; this is a teaser, not the findings.
\myparagraph{Significance}
Lorem ipsum dolor sit amet consectetur adipiscing elit. 
Quisque faucibus ex sapien vitae pellentesque sem placerat.
In id cursus mi pretium tellus duis convallis. 
Tempus leo eu aenean sed diam urna tempor.
Pulvinar vivamus fringilla lacus nec metus bibendum egestas. 

% --- PARAGRAPH 6: PAPER ROADMAP ---
% Function: Manuscript outline (Intro paragraph 6).
% Target: 2–4 sentences.
% Content: Walk the reader through the paper in document order so the roadmap
%   doubles as a structural map. The actual order is:
%   Background (sec:background) → Research Questions (sec:research_questions)
%   → RQ1 (sec:rq1_evaluation) → RQ2 (sec:rq2_evaluation)
%   → RQ3 (sec:rq3_evaluation) → Discussion (sec:discussion)
%   → Threats (sec:threats) → Conclusion (sec:conclusion).
%   NOTE: there is no standalone "Methodology" section — methodology is embedded
%   inside each RQ section. Use lowercase section names + \Cref macros; never
%   hardcode numbers.
\myparagraph{Paper Outline}
This paper is organized as follows: 
\Cref{sec:background} provides essential background and reviews related work, and 
\Cref{sec:research_questions} formalizes the research questions.
\Cref{sec:rq1_evaluation}, \Cref{sec:rq2_evaluation}, and \Cref{sec:rq3_evaluation} each present the methodology and empirical findings for one research question.
The paper concludes in \Cref{sec:discussion} with a discussion of implications for the software engineering community and future work, followed by an examination of potential threats to validity in \Cref{sec:threats} and closing remarks in \Cref{sec:conclusion}.

% =========================================================================
% SECTION: BACKGROUND AND RELATED WORK
% Narrative Arc: "What world are we in?" -> "What fails here?" -> "How do we fix it?"
% =========================================================================
\section{Background and Related Work}
\label{sec:background}

% --- MAJOR SIGNPOSTING & RE-CONTEXTUALIZATION ---
% Entry paragraph (3–5 sentences): re‑state the problem space with more technical depth than the Introduction and briefly signpost the three subsections (domain overview, existing approaches, conceptual foundation).
Lorem ipsum dolor sit amet consectetur adipiscing elit. 
Quisque faucibus ex sapien vitae pellentesque sem placerat.
In id cursus mi pretium tellus duis convallis. 
Tempus leo eu aenean sed diam urna tempor.
Pulvinar vivamus fringilla lacus nec metus bibendum egestas. 

% =========================================================================
% SUBSECTION 1: THE DOMAIN OVERVIEW (Domain overview)
% =========================================================================
\subsection{Subsection 1}
\label{subsec:subsection-1}

% Paragraph 1 – Domain definitions and workflows: 
% Technical definitions and standard workflows. Clarify entities (projects, pipelines, registries), typical processes, and key terminology.
Lorem ipsum dolor sit amet consectetur adipiscing elit. 
Quisque faucibus ex sapien vitae pellentesque sem placerat.
In id cursus mi pretium tellus duis convallis. 
Tempus leo eu aenean sed diam urna tempor.
Pulvinar vivamus fringilla lacus nec metus bibendum egestas. 
Iaculis massa nisl malesuada lacinia integer nunc posuere.

% Paragraph 2 – Cause‑and‑effect linkages: 
% Explain how practices and constraints in this domain lead to specific challenges (e.g., fragility, reproducibility issues). Use ‘As a result/Consequently’ to make causal links explicit.
Lorem ipsum dolor sit amet consectetur adipiscing elit.
Quisque faucibus ex sapien vitae pellentesque sem placerat. 
In id cursus mi pretium tellus duis convallis.
Tempus leo eu aenean sed diam urna tempor. 
Pulvinar vivamus fringilla lacus nec metus bibendum egestas.
Iaculis massa nisl malesuada lacinia integer nunc posuere. 

% Paragraph 3 – Transition to existing approaches: 
% Close the subsection by hinting at what practitioners or researchers currently do to cope with these challenges, setting up Subsection 2.
Lorem ipsum dolor sit amet consectetur adipiscing elit. 
Quisque faucibus ex sapien vitae pellentesque sem placerat.
In id cursus mi pretium tellus duis convallis. 
Tempus leo eu aenean sed diam urna tempor.
Pulvinar vivamus fringilla lacus nec metus bibendum egestas. 
Iaculis massa nisl malesuada lacinia integer nunc posuere.

% =========================================================================
% SUBSECTION 2: EXISTING APPROACHES & PRIOR WORK
% =========================================================================
\subsection{Subsection 2}
\label{subsec:subsection-2}

% Paragraph 1 – Subsection entry and literature grouping: 
% Target: 3–5 sentences. 
% Content: Explicitly tie back to Subsection 1 and describe high‑level categories of existing approaches (e.g., tooling, guidelines, empirical studies). Group literature by theme, not chronologically.
Lorem ipsum dolor sit amet consectetur adipiscing elit. 
Quisque faucibus ex sapien vitae pellentesque sem placerat.
In id cursus mi pretium tellus duis convallis. 
Tempus leo eu aenean sed diam urna tempor.
Pulvinar vivamus fringilla lacus nec metus bibendum egestas. 
Iaculis massa nisl malesuada lacinia integer nunc posuere.

% Paragraph 2 – Contrastive analysis and limitations: 
% Content: Describe what these approaches achieve and where they fall short. Use contrast markers (‘However’, ‘In contrast’) and connect limitations to your gap.
Lorem ipsum dolor sit amet consectetur adipiscing elit.
Quisque faucibus ex sapien vitae pellentesque sem placerat. 
In id cursus mi pretium tellus duis convallis.
Tempus leo eu aenean sed diam urna tempor. 
Pulvinar vivamus fringilla lacus nec metus bibendum egestas.
Iaculis massa nisl malesuada lacinia integer nunc posuere. 

% Paragraph 3 – Exit transition to conceptual lens: 
% Content: Summarize the overall shortcomings of existing work and point forward to the conceptual foundation you will use to reason about them (Subsection 3).
Lorem ipsum dolor sit amet consectetur adipiscing elit. 
Quisque faucibus ex sapien vitae pellentesque sem placerat.
In id cursus mi pretium tellus duis convallis. 
Tempus leo eu aenean sed diam urna tempor.
Pulvinar vivamus fringilla lacus nec metus bibendum egestas. 
Iaculis massa nisl malesuada lacinia integer nunc posuere.

% =========================================================================
% SUBSECTION 3: CONCEPTUAL FOUNDATION
% =========================================================================
\subsection{Subsection 3}
\label{subsec:subsection-3}

% Paragraph 1 – Subsection entry and theoretical framing: 
% Content: Introduce the conceptual/theoretical lens (e.g., socio‑technical, diffusion, reliability) that underpins your study. Motivate why this lens is appropriate given the domain and prior work.
Lorem ipsum dolor sit amet consectetur adipiscing elit. 
Quisque faucibus ex sapien vitae pellentesque sem placerat.
In id cursus mi pretium tellus duis convallis. 
Tempus leo eu aenean sed diam urna tempor.
Pulvinar vivamus fringilla lacus nec metus bibendum egestas. 
Iaculis massa nisl malesuada lacinia integer nunc posuere.

% Paragraph 2 – Technical hook into methodology: 
% Content: Tie the chosen lens directly to your upcoming methodology, hinting at key constructs, relationships, or hypotheses that will be operationalized in the Methods and RQ sections.
Lorem ipsum dolor sit amet consectetur adipiscing elit. 
Quisque faucibus ex sapien vitae pellentesque sem placerat.
In id cursus mi pretium tellus duis convallis. 
Tempus leo eu aenean sed diam urna tempor.
Pulvinar vivamus fringilla lacus nec metus bibendum egestas. 
Iaculis massa nisl malesuada lacinia integer nunc posuere.

% =========================================================================
% SECTION: RESEARCH QUESTIONS (The Structural Bridge)
% =========================================================================
\section{Research Questions}
\label{sec:research_questions}

% --- INTRODUCTORY SIGNPOSTING & GAP ANCHORING ---
% Target: 1–3 sentences. 
% Content: Remind the reader of the key gap from Background and state that you now formalize the evaluation goals as research questions.
Lorem ipsum dolor sit amet consectetur adipiscing elit. 
Quisque faucibus ex sapien vitae pellentesque sem placerat.
In id cursus mi pretium tellus duis convallis. 

% --- THE FORMAL RQ LIST ---
% Function: The formal, enumerated research questions (the structural bridge).
% Content: List 3–4 concise RQs, each isolating ONE dimension (e.g., prevalence,
%   patterns, impact). Keep each RQ to a single sentence prefaced by a bolded
%   dimensional anchor (e.g., [Prevalence]).
% Format: per the style guide, the colon sits OUTSIDE the bold: \textbf{RQ\arabic*}:
\noindent We address the following research questions:
\begin{enumerate}[label=\textbf{RQ\arabic*}:, leftmargin=*]
    \item \textbf{[Dimension 1, e.g., Prevalence].} Lorem ipsum dolor sit amet consectetur adipiscing elit?
    \item \textbf{[Dimension 2, e.g., Patterns].} Lorem ipsum dolor sit amet consectetur adipiscing elit?
    \item \textbf{[Dimension 3, e.g., Impact].} Lorem ipsum dolor sit amet consectetur adipiscing elit?
\end{enumerate}

% --- MAPPING PARAGRAPH ---
% Target: 2–3 sentences. 
% Content: Map each RQ to its evaluation section (RQ1, RQ2, RQ3) and briefly mention which data and methods answer which question. This proves the downstream structure aligns with stated goals.
Lorem ipsum dolor sit amet consectetur adipiscing elit. 
Quisque faucibus ex sapien vitae pellentesque sem placerat.
In id cursus mi pretium tellus duis convallis. 

% =========================================================================
% SECTION: EVALUATION OF RESEARCH QUESTION 1 (RQ1)
% RQ1 block arc: overview & motivation → RQ1‑specific methodology → RQ1 findings → formal answer box → pointer to Discussion.
% =========================================================================
\section{Evaluation of RQ1: [Dimension Title, e.g., Prevalence]}
\label{sec:rq1_evaluation}

% --- RQ1 OVERVIEW PARAGRAPH ---
% Target: 6–8 sentences. 
% Content: Restate RQ1 verbatim, explain its role relative to other RQs (baseline, patterns, impact), and describe why this dimension matters. End with a sentence that introduces the upcoming methodology subsection.
Lorem ipsum dolor sit amet consectetur adipiscing elit. 
Quisque faucibus ex sapien vitae pellentesque sem placerat. 
In id cursus mi pretium tellus duis convallis.
Tempus leo eu aenean sed diam urna tempor. 
Pulvinar vivamus fringilla lacus nec metus bibendum egestas.
Iaculis massa nisl malesuada lacinia integer nunc posuere. 
Ut hendrerit semper vel class aptent taciti sociosqu. 
Ad litora torquent per conubia nostra inceptos himenaeos.

% --- SUBSECTION: METHODOLOGY FOR RQ1 ---
\subsection{RQ1 Methodology and Operationalization}
\label{subsec:rq1_methodology}
% RULE: cite every figure in-text via \Cref{fig:...}; never rely on placement alone.
% =========================================================================
% TEMPLATE FIGURE PLACEHOLDER
% Uses a standard boxed paragraph container to stand in for missing visual assets.
% =========================================================================

\begin{figure}[h] % [h] requests placement "here" inline with the surrounding text stream
    \centering    % Horizontally centers the figure artifact within the column/page margins
    
    % \fbox draws a visible frame border around its inner content block
    \fbox{
        % \parbox structure creates an explicit, fixed-size paragraph boundary container
        % Arguments format: [inner-alignment][height][content-alignment]{width}
        \parbox[c][2in][c]{3in}{
            \centering Missing Figure % Centered placeholder text string inside the bounding box
        }
    }
    
    % Core descriptive label rendered immediately beneath the figure box frame
    \caption{Template figure placeholder}
    
    % Unique cross-referencing anchor tag; reference this asset elsewhere via \Cref{fig:template_figure}
    \label{fig:template_figure}
\end{figure}

% Paragraph 1 – Design & rationale: 
% Explain the study design (experiment, case study, mining, survey) and why it is appropriate for RQ1.
Lorem ipsum dolor sit amet consectetur adipiscing elit. 
Quisque faucibus ex sapien vitae pellentesque sem placerat. 
In id cursus mi pretium tellus duis convallis. 
Tempus leo eu aenean sed diam urna tempor.

% Paragraph 2 – Data sources & sampling: 
% Describe datasets, projects, participants, and inclusion/exclusion rules. Provide enough detail for replication.
Lorem ipsum dolor sit amet consectetur adipiscing elit. 
Quisque faucibus ex sapien vitae pellentesque sem placerat. 
In id cursus mi pretium tellus duis convallis. 
Tempus leo eu aenean sed diam urna tempor.

% Paragraph 3 – Procedure: 
% Outline the step‑by‑step execution (pipelines run, tasks assigned, scripts used), emphasizing chronology and reproducibility.
Lorem ipsum dolor sit amet consectetur adipiscing elit. 
Quisque faucibus ex sapien vitae pellentesque sem placerat. 
In id cursus mi pretium tellus duis convallis. 
Tempus leo eu aenean sed diam urna tempor.

% Paragraph 4 – Measures & analysis plan: 
% Enumerate metrics and constructs, describe how they are calculated, and specify statistical/analytical methods used to answer RQ1.
Lorem ipsum dolor sit amet consectetur adipiscing elit. 
Quisque faucibus ex sapien vitae pellentesque sem placerat. 
In id cursus mi pretium tellus duis convallis. 
Tempus leo eu aenean sed diam urna tempor.

% Threats to Validity for RQ1:
% Short paragraph. Note key RQ1‑specific threats (e.g., confounding, sampling, measurement) and point to the global Threats section for deeper analysis.
\myparagraph{Threats to Validity for RQ1}
Lorem ipsum dolor sit amet consectetur adipiscing elit. 
Quisque faucibus ex sapien vitae pellentesque sem placerat. 
In id cursus mi pretium tellus duis convallis. 

% --- SUBSECTION: EMPIRICAL FINDINGS FOR RQ1 ---
\subsection{RQ1 Empirical Findings}
\label{subsec:rq1_findings}

% Paragraph 1 – Overview and preprocessing: 
% Describe the subset of data used for RQ1, any preprocessing, exclusions, and basic descriptive stats (e.g., counts, ranges).
Lorem ipsum dolor sit amet consectetur adipiscing elit. 
Quisque faucibus ex sapien vitae pellentesque sem placerat. 
In id cursus mi pretium tellus duis convallis. 
Tempus leo eu aenean sed diam urna tempor.

% Paragraphs 2–3 – Main findings with figures/tables: 
% Present the core results in logical order (e.g., primary metrics, distributions, comparative stats). Every figure/table referenced must be explicitly cited in the text (‘As shown in Figure X…’). No interpretation beyond factual patterns.
Lorem ipsum dolor sit amet consectetur adipiscing elit. 
Quisque faucibus ex sapien vitae pellentesque sem placerat.
In id cursus mi pretium tellus duis convallis. 
Tempus leo eu aenean sed diam urna tempor.
Pulvinar vivamus fringilla lacus nec metus bibendum egestas. 
Iaculis massa nisl malesuada lacinia integer nunc posuere.
Ut hendrerit semper vel class aptent taciti sociosqu. 

% Paragraph 4 – Trends, patterns, and null results: 
% Summarize observed trends, highlight interesting edge cases, and explicitly mention negative or null results. Stay descriptive; save interpretation for the Discussion.
Lorem ipsum dolor sit amet consectetur adipiscing elit. 
Quisque faucibus ex sapien vitae pellentesque sem placerat. 
In id cursus mi pretium tellus duis convallis. 
Tempus leo eu aenean sed diam urna tempor.

% Formal ‘Answer to RQ1’ box: 
% 2–3 sentences of neutral synthesis directly answering the RQ text, with no speculation.
% =========================================================================
% EMPIRICAL FINDING DISPLAY BOX
% DEPENDENCY WARNING: This structural container relies directly on the 
% \newfinding command (defined in typeset/misc.tex). 
% Calling \newfinding inside this block:
%   1. Increments the central sequential 'finding' counter register.
%   2. Globally registers a target link for downstream \label and \ref calls.
%   3. Injects the bold string "Finding X." directly into the box canvas stream.
% =========================================================================

\begin{tcolorbox}[
    colback=yellow!30,    % Warm yellow alert canvas interior shading tint
    colframe=black,       % High-contrast solid black border outline
    sharp corners,        % Disables rounded bounding radii to mimic a standard \fcolorbox
    boxrule=0.5pt,        % Ultra-clean frame line width boundary scale
    % Comprehensive layout padding overrides to prevent text edge crowding:
    left=2mm, 
    right=2mm, 
    top=2mm, 
    bottom=5pt, 
    width=0.98\linewidth  % Consistently centers and bounds the box safely inside column margins
]
    % Note: Typically, you will invoke only ONE unique \newfinding call per box container.
    % Multiple consecutive calls here are used to demonstrate sequence step increments.
    \newfinding In id cursus mi pretium tellus duis convallis.
    
    \newfinding Pulvinar vivamus fringilla lacus nec metus bibendum.
    
    \newfinding Iaculis massa nisl malesuada lacinia integer.
    \label{finding: template-finding}
\end{tcolorbox}

% Pointer sentence to Discussion: 
% Single sentence that forwards the reader to the Discussion for qualitative interpretation and implications.
Lorem ipsum dolor sit amet consectetur adipiscing elit.

% =========================================================================
% SECTION: EVALUATION OF RESEARCH QUESTION 2 (RQ2)
% RQ2 block arc: overview & motivation → RQ2‑specific methodology → RQ2 findings → formal answer box → pointer to Discussion.
% =========================================================================
\section{Evaluation of RQ2: [Dimension Title, e.g., Patterns]}
\label{sec:rq2_evaluation}

% --- RQ2 OVERVIEW PARAGRAPH ---
% Target: 6–8 sentences.
% Content: Restate RQ2 verbatim, explain its role relative to the other RQs (it
%   builds on RQ1's baseline by characterizing patterns/mechanics), and explain
%   why this dimension matters. End with a sentence introducing the methodology.
Lorem ipsum dolor sit amet consectetur adipiscing elit.
Quisque faucibus ex sapien vitae pellentesque sem placerat. 
In id cursus mi pretium tellus duis convallis.
Tempus leo eu aenean sed diam urna tempor. 
Pulvinar vivamus fringilla lacus nec metus bibendum egestas.
Iaculis massa nisl malesuada lacinia integer nunc posuere. 
Ut hendrerit semper vel class aptent taciti sociosqu. 
Ad litora torquent per conubia nostra inceptos himenaeos.

% --- SUBSECTION: METHODOLOGY FOR RQ2 ---
\subsection{RQ2 Methodology and Categorization Pipeline}
\label{subsec:rq2_methodology}
% RULE: cite every figure in-text via \Cref{fig:...}; never rely on placement alone.
% =========================================================================
% TEMPLATE FIGURE PLACEHOLDER
% Uses a standard boxed paragraph container to stand in for missing visual assets.
% =========================================================================

\begin{figure}[h] % [h] requests placement "here" inline with the surrounding text stream
    \centering    % Horizontally centers the figure artifact within the column/page margins
    
    % \fbox draws a visible frame border around its inner content block
    \fbox{
        % \parbox structure creates an explicit, fixed-size paragraph boundary container
        % Arguments format: [inner-alignment][height][content-alignment]{width}
        \parbox[c][2in][c]{3in}{
            \centering Missing Figure % Centered placeholder text string inside the bounding box
        }
    }
    
    % Core descriptive label rendered immediately beneath the figure box frame
    \caption{Template figure placeholder}
    
    % Unique cross-referencing anchor tag; reference this asset elsewhere via \Cref{fig:template_figure}
    \label{fig:template_figure}
\end{figure}

% Paragraph 1 – Design & rationale:
% Explain the study design (experiment, case study, mining, survey) and why it
% is appropriate for RQ2.
Lorem ipsum dolor sit amet consectetur adipiscing elit.
Quisque faucibus ex sapien vitae pellentesque sem placerat. 
In id cursus mi pretium tellus duis convallis. 
Tempus leo eu aenean sed diam urna tempor.

% Paragraph 2 – Data sources & sampling:
% Describe datasets, projects, participants, and inclusion/exclusion rules.
% Provide enough detail for replication.
Lorem ipsum dolor sit amet consectetur adipiscing elit.
Quisque faucibus ex sapien vitae pellentesque sem placerat. 
In id cursus mi pretium tellus duis convallis. 
Tempus leo eu aenean sed diam urna tempor.

% Paragraph 3 – Procedure:
% Outline the step-by-step execution (pipelines run, tasks assigned, scripts
% used), emphasizing chronology and reproducibility.
Lorem ipsum dolor sit amet consectetur adipiscing elit.
Quisque faucibus ex sapien vitae pellentesque sem placerat. 
In id cursus mi pretium tellus duis convallis. 
Tempus leo eu aenean sed diam urna tempor.

% Paragraph 4 – Measures & analysis plan:
% Enumerate metrics and constructs, describe how they are calculated, and specify
% the statistical/analytical methods used to answer RQ2.
Lorem ipsum dolor sit amet consectetur adipiscing elit.
Quisque faucibus ex sapien vitae pellentesque sem placerat. 
In id cursus mi pretium tellus duis convallis. 
Tempus leo eu aenean sed diam urna tempor.

% Threats to Validity for RQ2:
% Short paragraph (2–3 sentences). Note key RQ2-specific threats (e.g.,
% confounding, sampling, measurement) and point to the global Threats section.
\myparagraph{Threats to Validity for RQ2}
Lorem ipsum dolor sit amet consectetur adipiscing elit.
Quisque faucibus ex sapien vitae pellentesque sem placerat. 
In id cursus mi pretium tellus duis convallis. 

% --- SUBSECTION: EMPIRICAL FINDINGS FOR RQ2 ---
\subsection{RQ2 Empirical Findings}
\label{subsec:rq2_findings}

% Paragraph 1 – Overview and preprocessing:
% Describe the subset of data used for RQ2, any preprocessing, exclusions, and
% basic descriptive stats (e.g., counts, ranges).
Lorem ipsum dolor sit amet consectetur adipiscing elit.
Quisque faucibus ex sapien vitae pellentesque sem placerat. 
In id cursus mi pretium tellus duis convallis. 
Tempus leo eu aenean sed diam urna tempor.

% Paragraphs 2–3 – Main findings with figures/tables:
% Present the core results in logical order. Every figure/table referenced must
% be explicitly cited in the text ("As shown in \Cref{fig:...}"). Stay strictly
% descriptive; no interpretation beyond factual patterns.
Lorem ipsum dolor sit amet consectetur adipiscing elit. 
Quisque faucibus ex sapien vitae pellentesque sem placerat.
In id cursus mi pretium tellus duis convallis. 
Tempus leo eu aenean sed diam urna tempor.
Pulvinar vivamus fringilla lacus nec metus bibendum egestas. 
Iaculis massa nisl malesuada lacinia integer nunc posuere.
Ut hendrerit semper vel class aptent taciti sociosqu. 

% Paragraph 4 – Trends, patterns, and null results:
% Summarize observed trends, highlight interesting edge cases, and explicitly
% mention negative or null results. Stay descriptive; save interpretation for the
% Discussion.
Lorem ipsum dolor sit amet consectetur adipiscing elit.
Quisque faucibus ex sapien vitae pellentesque sem placerat. 
In id cursus mi pretium tellus duis convallis. 
Tempus leo eu aenean sed diam urna tempor. 

% Formal 'Answer to RQ2' box:
% 2–3 sentences of neutral synthesis directly answering the RQ text, no speculation.
% =========================================================================
% EMPIRICAL FINDING DISPLAY BOX
% DEPENDENCY WARNING: This structural container relies directly on the 
% \newfinding command (defined in typeset/misc.tex). 
% Calling \newfinding inside this block:
%   1. Increments the central sequential 'finding' counter register.
%   2. Globally registers a target link for downstream \label and \ref calls.
%   3. Injects the bold string "Finding X." directly into the box canvas stream.
% =========================================================================

\begin{tcolorbox}[
    colback=yellow!30,    % Warm yellow alert canvas interior shading tint
    colframe=black,       % High-contrast solid black border outline
    sharp corners,        % Disables rounded bounding radii to mimic a standard \fcolorbox
    boxrule=0.5pt,        % Ultra-clean frame line width boundary scale
    % Comprehensive layout padding overrides to prevent text edge crowding:
    left=2mm, 
    right=2mm, 
    top=2mm, 
    bottom=5pt, 
    width=0.98\linewidth  % Consistently centers and bounds the box safely inside column margins
]
    % Note: Typically, you will invoke only ONE unique \newfinding call per box container.
    % Multiple consecutive calls here are used to demonstrate sequence step increments.
    \newfinding In id cursus mi pretium tellus duis convallis.
    
    \newfinding Pulvinar vivamus fringilla lacus nec metus bibendum.
    
    \newfinding Iaculis massa nisl malesuada lacinia integer.
    \label{finding: template-finding}
\end{tcolorbox}

% Pointer sentence to Discussion:
% Single sentence forwarding the reader to the Discussion for qualitative
% interpretation and implications.
Lorem ipsum dolor sit amet consectetur adipiscing elit.

% =========================================================================
% SECTION: EVALUATION OF RESEARCH QUESTION 3 (RQ3)
% RQ3 block arc: overview & motivation → RQ3‑specific methodology → RQ3 findings → formal answer box → pointer to Discussion.
% =========================================================================
\section{Evaluation of RQ3: [Dimension Title, e.g., Impact]}
\label{sec:rq3_evaluation}

% --- RQ3 OVERVIEW PARAGRAPH ---
% Target: 6–8 sentences.
% Content: Restate RQ3 verbatim, explain its role relative to the other RQs (it
%   typically caps the arc by assessing impact/consequence), and explain why this
%   dimension matters. End with a sentence introducing the methodology subsection.
Lorem ipsum dolor sit amet consectetur adipiscing elit.
Quisque faucibus ex sapien vitae pellentesque sem placerat. 
In id cursus mi pretium tellus duis convallis.
Tempus leo eu aenean sed diam urna tempor. 
Pulvinar vivamus fringilla lacus nec metus bibendum egestas.
Iaculis massa nisl malesuada lacinia integer nunc posuere. 
Ut hendrerit semper vel class aptent taciti sociosqu. 
Ad litora torquent per conubia nostra inceptos himenaeos.

% --- SUBSECTION: METHODOLOGY FOR RQ3 ---
\subsection{RQ3 Methodology and Analytical Matrix}
\label{subsec:rq3_methodology}
% RULE: cite every figure in-text via \Cref{fig:...}; never rely on placement alone.
% =========================================================================
% TEMPLATE FIGURE PLACEHOLDER
% Uses a standard boxed paragraph container to stand in for missing visual assets.
% =========================================================================

\begin{figure}[h] % [h] requests placement "here" inline with the surrounding text stream
    \centering    % Horizontally centers the figure artifact within the column/page margins
    
    % \fbox draws a visible frame border around its inner content block
    \fbox{
        % \parbox structure creates an explicit, fixed-size paragraph boundary container
        % Arguments format: [inner-alignment][height][content-alignment]{width}
        \parbox[c][2in][c]{3in}{
            \centering Missing Figure % Centered placeholder text string inside the bounding box
        }
    }
    
    % Core descriptive label rendered immediately beneath the figure box frame
    \caption{Template figure placeholder}
    
    % Unique cross-referencing anchor tag; reference this asset elsewhere via \Cref{fig:template_figure}
    \label{fig:template_figure}
\end{figure}

% Paragraph 1 – Design & rationale:
% Explain the study design (experiment, case study, mining, survey) and why it
% is appropriate for RQ3.
Lorem ipsum dolor sit amet consectetur adipiscing elit.
Quisque faucibus ex sapien vitae pellentesque sem placerat. 
In id cursus mi pretium tellus duis convallis. 
Tempus leo eu aenean sed diam urna tempor.

% Paragraph 2 – Data sources & sampling:
% Describe datasets, projects, participants, and inclusion/exclusion rules.
% Provide enough detail for replication.
Lorem ipsum dolor sit amet consectetur adipiscing elit.
Quisque faucibus ex sapien vitae pellentesque sem placerat. 
In id cursus mi pretium tellus duis convallis. 
Tempus leo eu aenean sed diam urna tempor. 

% Paragraph 3 – Procedure:
% Outline the step-by-step execution (pipelines run, tasks assigned, scripts
% used), emphasizing chronology and reproducibility.
Lorem ipsum dolor sit amet consectetur adipiscing elit.
Quisque faucibus ex sapien vitae pellentesque sem placerat. 
In id cursus mi pretium tellus duis convallis. 
Tempus leo eu aenean sed diam urna tempor.

% Paragraph 4 – Measures & analysis plan:
% Enumerate metrics and constructs, describe how they are calculated, and specify
% the statistical/analytical methods used to answer RQ3.
Lorem ipsum dolor sit amet consectetur adipiscing elit.
Quisque faucibus ex sapien vitae pellentesque sem placerat. 
In id cursus mi pretium tellus duis convallis. 
Tempus leo eu aenean sed diam urna tempor.

% Threats to Validity for RQ3:
% Short paragraph (2–3 sentences). Note key RQ3-specific threats (e.g.,
% confounding, sampling, measurement) and point to the global Threats section.
\myparagraph{Threats to Validity for RQ3}
Lorem ipsum dolor sit amet consectetur adipiscing elit.
Quisque faucibus ex sapien vitae pellentesque sem placerat. 
In id cursus mi pretium tellus duis convallis. 

% --- SUBSECTION: EMPIRICAL FINDINGS FOR RQ3 ---
\subsection{RQ3 Empirical Findings}
\label{subsec:rq3_findings}

% Paragraph 1 – Overview and preprocessing:
% Describe the subset of data used for RQ3, any preprocessing, exclusions, and
% basic descriptive stats (e.g., counts, ranges).
Lorem ipsum dolor sit amet consectetur adipiscing elit.
Quisque faucibus ex sapien vitae pellentesque sem placerat. 
In id cursus mi pretium tellus duis convallis. 
Tempus leo eu aenean sed diam urna tempor.

% Paragraphs 2–3 – Main findings with figures/tables:
% Present the core results in logical order. Every figure/table referenced must
% be explicitly cited in the text ("As shown in \Cref{fig:...}"). Stay strictly
% descriptive; no interpretation beyond factual patterns.
Lorem ipsum dolor sit amet consectetur adipiscing elit. 
Quisque faucibus ex sapien vitae pellentesque sem placerat.
In id cursus mi pretium tellus duis convallis. 
Tempus leo eu aenean sed diam urna tempor.
Pulvinar vivamus fringilla lacus nec metus bibendum egestas. 
Iaculis massa nisl malesuada lacinia integer nunc posuere.
Ut hendrerit semper vel class aptent taciti sociosqu. 

% Paragraph 4 – Trends, patterns, and null results:
% Summarize observed trends, highlight interesting edge cases, and explicitly
% mention negative or null results. Stay descriptive; save interpretation for the
% Discussion.
Lorem ipsum dolor sit amet consectetur adipiscing elit.
Quisque faucibus ex sapien vitae pellentesque sem placerat. 
In id cursus mi pretium tellus duis convallis. 
Tempus leo eu aenean sed diam urna tempor. 

% Formal 'Answer to RQ3' box:
% 2–3 sentences of neutral synthesis directly answering the RQ text, no speculation.
% =========================================================================
% EMPIRICAL FINDING DISPLAY BOX
% DEPENDENCY WARNING: This structural container relies directly on the 
% \newfinding command (defined in typeset/misc.tex). 
% Calling \newfinding inside this block:
%   1. Increments the central sequential 'finding' counter register.
%   2. Globally registers a target link for downstream \label and \ref calls.
%   3. Injects the bold string "Finding X." directly into the box canvas stream.
% =========================================================================

\begin{tcolorbox}[
    colback=yellow!30,    % Warm yellow alert canvas interior shading tint
    colframe=black,       % High-contrast solid black border outline
    sharp corners,        % Disables rounded bounding radii to mimic a standard \fcolorbox
    boxrule=0.5pt,        % Ultra-clean frame line width boundary scale
    % Comprehensive layout padding overrides to prevent text edge crowding:
    left=2mm, 
    right=2mm, 
    top=2mm, 
    bottom=5pt, 
    width=0.98\linewidth  % Consistently centers and bounds the box safely inside column margins
]
    % Note: Typically, you will invoke only ONE unique \newfinding call per box container.
    % Multiple consecutive calls here are used to demonstrate sequence step increments.
    \newfinding In id cursus mi pretium tellus duis convallis.
    
    \newfinding Pulvinar vivamus fringilla lacus nec metus bibendum.
    
    \newfinding Iaculis massa nisl malesuada lacinia integer.
    \label{finding: template-finding}
\end{tcolorbox}

% Pointer sentence to Discussion:
% Single sentence forwarding the reader to the Discussion for qualitative
% interpretation and implications.
Lorem ipsum dolor sit amet consectetur adipiscing elit.
Quisque faucibus ex sapien vitae pellentesque sem placerat. 
In id cursus mi pretium tellus duis convallis. 
Tempus leo eu aenean sed diam urna tempor. 

% =========================================================================
% SECTION: DISCUSSION
% Narrative Flow: Narrow (Interpretations) -> Broad (Implications) 
%                -> Defensive (Limitations) -> Forward-Looking (Future Work)
% =========================================================================
\section{Discussion}
\label{sec:discussion}

% --- PARAGRAPH 1: DISCUSSION OPENING ---
% Target: 7–9 sentences. 
% Content: Restate the high‑level takeaway of the study, tying back to the Introduction’s big gap and summarizing the overall pattern across RQ1–RQ3. This paragraph sets the tone for the rest of the Discussion.
Lorem ipsum dolor sit amet consectetur adipiscing elit. 
Quisque faucibus ex sapien vitae pellentesque sem placerat. 
In id cursus mi pretium tellus duis convallis. 
Tempus leo eu aenean sed diam urna tempor. 
Pulvinar vivamus fringilla lacus nec metus bibendum egestas. 
Iaculis massa nisl malesuada lacinia integer nunc posuere. 
Ut hendrerit semper vel class aptent taciti sociosqu. 

% --- SUBSECTION: INTERPRETATION OF RQS ---
% Per‑RQ mini‑discussions: each \myparagraph is 8–12 sentences total, structured as interpretation → mechanism → relation to prior work → scoped implication/limit.
\subsection{Interpretation of Research Questions}
\label{subsec:interpretations}

\myparagraph{RQ$_1$ -- Baseline Prevalence Metrics}
% Interpretation: 2–3 sentences summarizing how the findings answer RQ1.
Lorem ipsum dolor sit amet consectetur adipiscing elit. Quisque faucibus ex sapien vitae pellentesque sem placerat. In id cursus mi pretium tellus duis convallis.
% Mechanism/Why: 2–3 sentences explaining why these patterns likely occur, grounded in theory or domain knowledge.
Lorem ipsum dolor sit amet consectetur adipiscing elit. Quisque faucibus ex sapien vitae pellentesque sem placerat. In id cursus mi pretium tellus duis convallis.
% Relation to prior work: 2–3 sentences comparing RQ1 findings to specific prior studies; note agreements and contradictions.
Lorem ipsum dolor sit amet consectetur adipiscing elit. Quisque faucibus ex sapien vitae pellentesque sem placerat. In id cursus mi pretium tellus duis convallis.
% Scoped actionable/limit: 2–3 sentences stating implications or recommendations for practice, plus one key limitation that tempers those recommendations.
Lorem ipsum dolor sit amet consectetur adipiscing elit. Quisque faucibus ex sapien vitae pellentesque sem placerat. In id cursus mi pretium tellus duis convallis. Tempus leo eu aenean sed diam urna tempor.

\myparagraph{RQ$_2$ -- The Mechanics of Fragile Reuse}
% Interpretation: 2–3 sentences summarizing how the findings answer RQ2.
Lorem ipsum dolor sit amet consectetur adipiscing elit.
% Mechanism/Why: 2–3 sentences explaining why these patterns likely occur, grounded in theory or domain knowledge.
Lorem ipsum dolor sit amet consectetur adipiscing elit. Quisque faucibus ex sapien vitae pellentesque sem placerat. In id cursus mi pretium tellus duis convallis.
% Relation to prior work: 2–3 sentences comparing RQ2 findings to specific prior studies.
Lorem ipsum dolor sit amet consectetur adipiscing elit. Quisque faucibus ex sapien vitae pellentesque sem placerat. In id cursus mi pretium tellus duis convallis.
% Scoped actionable/limit: 2–3 sentences stating implications or recommendations for practice.
Lorem ipsum dolor sit amet consectetur adipiscing elit. Quisque faucibus ex sapien vitae pellentesque sem placerat. In id cursus mi pretium tellus duis convallis. Tempus leo eu aenean sed diam urna tempor.

\myparagraph{RQ$_3$ -- Quantification of Pipeline Half-Life}
% Interpretation: 2–3 sentences summarizing how the findings answer RQ3.
Lorem ipsum dolor sit amet consectetur adipiscing elit.
% Mechanism/Why: 2–3 sentences explaining why these patterns likely occur, grounded in theory or domain knowledge.
Lorem ipsum dolor sit amet consectetur adipiscing elit. Quisque faucibus ex sapien vitae pellentesque sem placerat. In id cursus mi pretium tellus duis convallis.
% Relation to prior work: 2–3 sentences comparing RQ3 findings to specific prior studies.
Lorem ipsum dolor sit amet consectetur adipiscing elit. Quisque faucibus ex sapien vitae pellentesque sem placerat. In id cursus mi pretium tellus duis convallis.
% Scoped actionable/limit: 2–3 sentences stating implications or recommendations for practice.
Lorem ipsum dolor sit amet consectetur adipiscing elit. Quisque faucibus ex sapien vitae pellentesque sem placerat. In id cursus mi pretium tellus duis convallis. Tempus leo eu aenean sed diam urna tempor.

% --- SUBSECTION: WIDER IMPLICATIONS ---
\subsection{Wider Implications for Science and Engineering}
\label{subsec:implications}

% Paragraph 1 – Theoretical implications: 
% Target: 5–7 sentences. 
% Content: Synthesize across RQs to state how the study refines or challenges existing theories/models. Focus on general principles.
Lorem ipsum dolor sit amet consectetur adipiscing elit. 
Quisque faucibus ex sapien vitae pellentesque sem placerat. 
In id cursus mi pretium tellus duis convallis. 
Tempus leo eu aenean sed diam urna tempor. 
Pulvinar vivamus fringilla lacus nec metus bibendum egestas. 
Iaculis massa nisl malesuada lacinia integer nunc posuere. 
Ut hendrerit semper vel class aptent taciti sociosqu.

% Paragraph 2 – Practical/industrial implications: 
% Target: 5–7 sentences. 
% Content: Translate findings into concrete recommendations for teams, tools, and processes, mentioning feasibility and constraints.
Lorem ipsum dolor sit amet consectetur adipiscing elit. 
Quisque faucibus ex sapien vitae pellentesque sem placerat. 
In id cursus mi pretium tellus duis convallis. 
Tempus leo eu aenean sed diam urna tempor. 
Pulvinar vivamus fringilla lacus nec metus bibendum egestas. 
Iaculis massa nisl malesuada lacinia integer nunc posuere. 
Ut hendrerit semper vel class aptent taciti sociosqu.

% Optional 3rd Paragraph – Policy / Educational Frameworks:
% NOTE: Optional, beyond the two standard implication types (Theoretical and
%   Practical/Industrial) defined by the style guide. Delete if not needed.
% Target: 5–7 sentences on implications for standards, curricula, or broader scientific practice.
Lorem ipsum dolor sit amet consectetur adipiscing elit. 
Quisque faucibus ex sapien vitae pellentesque sem placerat. 
In id cursus mi pretium tellus duis convallis. 
Tempus leo eu aenean sed diam urna tempor. 
Pulvinar vivamus fringilla lacus nec metus bibendum egestas. 

% --- SUBSECTION: LIMITATIONS AND BIASES ---
\subsection{Limitations and Alternative Explanations}
\label{subsec:limitations}

% Paragraph 1 – Main limitations: 
% Target: 5–7 sentences. 
% Content: Summarize key limitations across design, data, and measures; explain how each might bias results or limit generality.
Lorem ipsum dolor sit amet consectetur adipiscing elit. 
Quisque faucibus ex sapien vitae pellentesque sem placerat. 
In id cursus mi pretium tellus duis convallis. 
Tempus leo eu aenean sed diam urna tempor. 
Pulvinar vivamus fringilla lacus nec metus bibendum egestas. 

% Paragraph 2 – Alternative explanations: 
% Target: 5–7 sentences. 
% Content: Present plausible alternative interpretations of major findings and briefly note what further data/methods would be needed to test them.
Lorem ipsum dolor sit amet consectetur adipiscing elit. 
Quisque faucibus ex sapien vitae pellentesque sem placerat. 
In id cursus mi pretium tellus duis convallis. 
Tempus leo eu aenean sed diam urna tempor. 
Pulvinar vivamus fringilla lacus nec metus bibendum egestas. 

% --- SUBSECTION: FUTURE DIRECTIONS ---
\subsection{Future Research Trajectories}
\label{subsec:future_directions}

% Paragraph 1 – Next logical steps: 
% Target: 4–6 sentences. 
% Content: Propose immediate follow‑up studies (replications, scaling, boundary conditions), each explicitly linked to a specific result or limitation.
Lorem ipsum dolor sit amet consectetur adipiscing elit. 
Quisque faucibus ex sapien vitae pellentesque sem placerat. 
In id cursus mi pretium tellus duis convallis. 
Tempus leo eu aenean sed diam urna tempor. 

% Paragraph 2 – Methodological refinements: 
% Target: 3–5 sentences. 
% Content: Suggest improvements to design, metrics, or data collection that would strengthen future evidence.
Lorem ipsum dolor sit amet consectetur adipiscing elit. 
Quisque faucibus ex sapien vitae pellentesque sem placerat. 
In id cursus mi pretium tellus duis convallis. 

% Paragraph 3 – New questions: 
% Target: 4–6 sentences. 
% Content: Outline new research questions opened by unexpected findings or implications; explain why each is important for empirical SE.
Lorem ipsum dolor sit amet consectetur adipiscing elit. 
Quisque faucibus ex sapien vitae pellentesque sem placerat. 
In id cursus mi pretium tellus duis convallis. 
Tempus leo eu aenean sed diam urna tempor. 

% Final closing block – Take-Home Message:
% Target: 4–6 sentences. 
% Content: Provide a synthesized, confident statement of what the study ultimately shows and why it matters. No new ideas or data; this should echo the core contribution and high‑level impact.
\myparagraph{Take-Home Message}
Lorem ipsum dolor sit amet consectetur adipiscing elit. 
Quisque faucibus ex sapien vitae pellentesque sem placerat. 
In id cursus mi pretium tellus duis convallis. 
Tempus leo eu aenean sed diam urna tempor. 
Pulvinar vivamus fringilla lacus nec metus bibendum egestas. 
Iaculis massa nisl malesuada lacinia integer nunc posuere. 

% =========================================================================
% SECTION: THREATS TO VALIDITY
% Function: Systematically classify study vulnerabilities into three subsections
%   (Internal, External, Construct validity). Be candid; each block names a threat,
%   its potential effect on results, and the mitigation taken.
% =========================================================================
\section{Threats to Validity}
\label{sec:threats}

% Section opening:
% Target: 3–4 sentences. 
% Content: Explain the purpose of this section (systematically analyzing internal, external, and construct validity threats) and note that detailed RQ‑specific threats were mentioned in each RQ methodology subsection.
Lorem ipsum dolor sit amet consectetur adipiscing elit. 
Quisque faucibus ex sapien vitae pellentesque sem placerat. 
In id cursus mi pretium tellus duis convallis. 
Our analysis follows the standard empirical software engineering classification distinguishing between internal, external, and construct validity.

% --- SUBSECTION: INTERNAL VALIDITY ---
\subsection{Internal Validity}
\label{subsec:internal_validity}

% First block (Confounding Variables):
% Target: 3–5 sentences.
% Content: Explain potential confounding variables and how they may distort causal
%   interpretations; name the mitigation strategies used. Use causal signposts
%   ("As a result", "Consequently").
Lorem ipsum dolor sit amet consectetur adipiscing elit. 
Quisque faucibus ex sapien vitae pellentesque sem placerat. 
In id cursus mi pretium tellus duis convallis. 
Tempus leo eu aenean sed diam urna tempor. 
Pulvinar vivamus fringilla lacus nec metus bibendum egestas. 

% Second block (History, Maturation, and Instrumentation):
% Target: 3–4 sentences.
% Content: Discuss history, maturation, and instrumentation effects and their
%   potential influence on outcomes; note how design choices attempt to control them.
Lorem ipsum dolor sit amet consectetur adipiscing elit. 
Quisque faucibus ex sapien vitae pellentesque sem placerat. 
In id cursus mi pretium tellus duis convallis. 
Tempus leo eu aenean sed diam urna tempor. 

% Third block (Attrition and Missing Data):
% Target: 2–3 sentences.
% Content: Describe attrition/missing data patterns and how they were handled
%   (e.g., exclusion, imputation, sensitivity analysis).
Lorem ipsum dolor sit amet consectetur adipiscing elit. 
Quisque faucibus ex sapien vitae pellentesque sem placerat. 
In id cursus mi pretium tellus duis convallis. 

% --- SUBSECTION: EXTERNAL VALIDITY ---
\subsection{External Validity}
\label{subsec:external_validity}

% First block (Context & Generalizability):
% Target: 3–4 sentences.
% Content: Clarify the study setting and domains, and explicitly state the boundary
%   conditions — where the results do and do not generalize.
Lorem ipsum dolor sit amet consectetur adipiscing elit. 
Quisque faucibus ex sapien vitae pellentesque sem placerat. 
In id cursus mi pretium tellus duis convallis. 
Tempus leo eu aenean sed diam urna tempor. 

% Second block (Sampling Bias):
% Target: 2–3 sentences.
% Content: Discuss how selection of projects/participants may skew results and how
%   that limits external validity.
Lorem ipsum dolor sit amet consectetur adipiscing elit. 
Quisque faucibus ex sapien vitae pellentesque sem placerat. 
In id cursus mi pretium tellus duis convallis. 

% --- SUBSECTION: CONSTRUCT VALIDITY ---
\subsection{Construct Validity}
\label{subsec:construct_validity}

% First block (Indirect Metrics and Operationalization):
% Target: 4–6 sentences.
% Content: Justify that the chosen metrics capture the intended constructs,
%   referencing prior work and triangulation where possible.
Lorem ipsum dolor sit amet consectetur adipiscing elit. 
Quisque faucibus ex sapien vitae pellentesque sem placerat. 
In id cursus mi pretium tellus duis convallis. 
Tempus leo eu aenean sed diam urna tempor. 
Pulvinar vivamus fringilla lacus nec metus bibendum egestas. 
Iaculis massa nisl malesuada lacinia integer nunc posuere. 
Ut hendrerit semper vel class aptent taciti sociosqu. 
Ad litora torquent per conubia nostra inceptos himenaeos.

% Second block (Measurement Error and Technical Bias):
% Target: 3–5 sentences.
% Content: Address logging inaccuracies, self-report biases, and tool limitations;
%   explain how they might affect interpretations.
Lorem ipsum dolor sit amet consectetur adipiscing elit. 
Quisque faucibus ex sapien vitae pellentesque sem placerat. 
In id cursus mi pretium tellus duis convallis. 
Tempus leo eu aenean sed diam urna tempor. 
Pulvinar vivamus fringilla lacus nec metus bibendum egestas. 
Iaculis massa nisl malesuada lacinia integer nunc posuere. 

% =========================================================================
% SECTION: CONCLUSION
% =========================================================================
\section{Conclusion}
\label{sec:conclusion}

% Paragraph 1 – Aims and findings synthesis: 
% Target: 2–4 sentences. 
% Content: Restate the overall problem and aims in fresh language, then synthesize the main findings across RQs (no detailed numbers).
Lorem ipsum dolor sit amet consectetur adipiscing elit. 
Quisque faucibus ex sapien vitae pellentesque sem placerat. 
In id cursus mi pretium tellus duis convallis. 
Tempus leo eu aenean sed diam urna tempor. 
Pulvinar vivamus fringilla lacus nec metus bibendum egestas. 
Iaculis massa nisl malesuada lacinia integer nunc posuere. 
Ut hendrerit semper vel class aptent taciti sociosqu.

% Paragraph 2 – Contributions and practical implications: 
% Target: 2–4 sentences. 
% Content: Explicitly name the key contributions and what they enable or change for researchers and practitioners.
Lorem ipsum dolor sit amet consectetur adipiscing elit. 
Quisque faucibus ex sapien vitae pellentesque sem placerat. 
In id cursus mi pretium tellus duis convallis. 
Tempus leo eu aenean sed diam urna tempor. 
Pulvinar vivamus fringilla lacus nec metus bibendum egestas. 
Iaculis massa nisl malesuada lacinia integer nunc posuere. 
Ut hendrerit semper vel class aptent taciti sociosqu.

% Paragraph 3 – Limitations, future work, final takeaway: 
% Target: 3–4 sentences. 
% Style Rule: The final sentence must be a confident, synthesized statement—not an apology.
% Content: Echo the most important limitation/future direction at a high level and end with a strong single sentence that captures the study’s core message.
Lorem ipsum dolor sit amet consectetur adipiscing elit. 
Quisque faucibus ex sapien vitae pellentesque sem placerat. 
In id cursus mi pretium tellus duis convallis. 
Tempus leo eu aenean sed diam urna tempor. 
Pulvinar vivamus fringilla lacus nec metus bibendum egestas. 
Iaculis massa nisl malesuada lacinia integer nunc posuere. 
Ut hendrerit semper vel class aptent taciti sociosqu. 

% =========================================================================
% ACKNOWLEDGMENTS
% =========================================================================
\section*{Acknowledgments}
% Target: 1–3 sentences. 
% Content: Thank funding sources, collaborators, and infrastructure. No technical content.
Lorem ipsum dolor sit amet consectetur adipiscing elit. 
Quisque faucibus ex sapien vitae pellentesque sem placerat. 
In id cursus mi pretium tellus duis convallis. 

\bibliographystyle{ACM-Reference-Format}
\bibliography{citations}    % For all other submissions
%% =========================================================================
% DOCUMENT CLASS & SETUP
% =========================================================================

% \documentclass[sigconf]{acmart}       % For final camera-ready peer-reviewed submissions
\documentclass[sigconf,nonacm]{acmart} % For arXiv/preprint: removes copyright blocks & ACM footers

% =========================================================================
% GLOBAL PARAGRAPH AND PROSE RULES
% =========================================================================
% 1. Target Length: Aim for 3–7 sentences per paragraph. Avoid single-sentence 
%    paragraphs and dense walls of text.
% 2. Functional Focus: Each paragraph must serve a singular rhetorical function 
%    (e.g., context, problem, method, finding, interpretation, implication, 
%    limitation, or future work).
% 3. Explicit Signposting: Maintain visible logical flow across paragraph 
%    boundaries using structural transitions (e.g., "However", "As a result", 
%    "Consequently", "In contrast").
% 4. Micro-Sections: Use \myparagraph{...} blocks to signify named mini-paragraphs 
%    with specialized micro-rhetorical goals (e.g., "Contributions", "Significance").

% =========================================================================
% CONDITIONAL DEBUGGING FLAGS
% =========================================================================

\newif\ifDEBUG  % Declares a boolean switch called "DEBUG" for conditional compilation
\DEBUGtrue      % Set DEBUG to true (activates draft/todo notes)
\DEBUGfalse     % Set DEBUG to false (silences draft notes for clean compilation)
% NET EFFECT: \DEBUGfalse runs last, so debug is OFF by default. While OFF,
% \TODO{}, \CITE{}, and \NMS{} (co-author notes) render as nothing.
% To see draft markers, comment out the \DEBUGfalse line above.

% =========================================================================
% MODULAR PREAMBLE IMPORTS
% =========================================================================

% =========================================================================
% misc/data.tex
% Global Data Macro Definitions
% =========================================================================
        % Imports global stats, variables, and data macros
% =========================================================================
% core typography & layout engine
% =========================================================================

% \usepackage{amssymb,amsmath,amsfonts} % Optional math packages (currently commented)
\usepackage[T1]{fontenc}               % Uses modern 8-bit font encoding for crisp glyph rendering
\usepackage{balance}                   % Equalizes the length of columns on the final page of a 2-column layout

% =========================================================================
% code listing & custom graphical callout containers
% =========================================================================

\usepackage[most]{tcolorbox}           % Advanced custom framed boxes with advanced skins/listings hooks
\usepackage{listings}                  % Native source code block layout engine
\usepackage{xcolor}                    % Comprehensive driver-independent color naming & mixing support

% =========================================================================
% standard acm-approved helper packages
% =========================================================================

% --- structural & page geometry control
\usepackage{afterpage}                 % Postpones execution of commands until the next page break
\usepackage{float}                     % Enhances figure place control; introduces the rigid [H] specifier
\usepackage{lscape}                    % Places specific content pages in landscape orientation
\usepackage{rotating}                  % Rotates text, tables, and figures along arbitrary angle limits

% --- mathematics & computational logic
\usepackage{algorithmic}               % Structural pseudocode layout environment builder

% --- advanced tabular structures
\usepackage{array}                     % Extended configuration controls for table column definitions
\usepackage{booktabs}                  % Renders publication-grade professional horizontal table rules
\usepackage{makecell}                  % Multiline table cell formatting with manual line breaks
\usepackage{multirow}                  % Enables cells spanning multiple horizontal rows
\usepackage{tabularx}                  % Auto-calculates column widths to fit specified horizontal constraints
\usepackage{longtable}                 % Handles multi-page tables natively within the standard stream

% --- captions & bibliographic citations
\usepackage[skip=4pt]{caption}         % Uniform control over caption spacing across figures and tables
\usepackage[compress]{cite}            % Automatically sorts and collapses numeric sequence reference tags
\usepackage{natbib}[numbers]           % Explicit configuration for structured numeric reference citation types

% --- document hyperlinking & reference automation
\usepackage{hyperref}                  % Injects interactive hyperlinks and global document indexes
\usepackage{cleveref}                  % Smart inline references (adds contextual prefixes like "Section" or "Fig.")

% --- list layout control
\usepackage{enumitem}                  % Fine-grained layout control over itemize, enumerate, and description lists

% --- graphic assets, icons, & utility packages
% \usepackage{filecontents}            % Pass-through payload file management (currently commented)
\usepackage{fontawesome5}              % Vector icon engine for system design and repository markers
\usepackage{graphicx}                  % Core graphics scaling, cropping, and rendering extension
\usepackage{lipsum}                    % Generates dummy text placeholders for layout prototyping
\usepackage{siunitx}                   % Uniform typesetting formatting engine for technical units and data measurements
\usepackage{soul}                      % Provides tracking, kerning, and strike-through/text-highlighting features
\usepackage{textcomp}                  % Provides specialized miscellaneous symbols and text glyphs
\usepackage[normalem]{ulem}            % Smooth underlining control; [normalem] prevents breaking \emph highlights
\usepackage{xspace}                    % Intelligent trailing space macro insertion engine

% =========================================================================
% page design & heading spacing tweaks
% =========================================================================

% \usepackage[expectsections, unnumberedsections, fencedCode]{markdown} % Direct Markdown parsing engine (commented)
% \usepackage[vskip=1em,font=itshape,leftmargin=2em,rightmargin=2em]{quoting} % Advanced blockquote layout engine (commented)

% Compacted section structural layout overrides (currently commented)
% \titlespacing*{\section}{0pt}{6pt}{4pt}
% \titlespacing*{\subsection}{0pt}{5pt}{3pt}
% \titlespacing*{\subsubsection}{0pt}{4pt}{2pt}

% =========================================================================
% CUSTOM UTILITY MACROS
% =========================================================================

% --- standard abbreviated latin syntax styling
\newcommand{\ie}{\textit{i.e.},\xspace}  % Uniformly formats "id est" with correct italicization and spacing
\newcommand{\eg}{\textit{e.g.},\xspace}  % Uniformly formats "exempli gratia" with correct italicization and spacing

% --- layout constructs
\newcommand{\myparagraph}[1]{\vspace{0.20cm}\noindent\textbf{#1} \noindent{}} % Spaced bold headers for inline subsections

% --- generative framework / prompts callout containers
% =========================================================================
% LLM PROMPT MACRO TEMPLATE
% Purpose: Formats Large Language Model system instructions, context windows,
%          or user prompts within a distinct, clean visual container.
%
% Usage: \myprompt{Argument #1: Box Heading}{Argument #2: Prompt Text Payload}
% Example: \myprompt{System Prompt}{Identify software reuse patterns...}
% =========================================================================

\newcommand{\myprompt}[2]{%
  \begin{tcolorbox}[
    enhanced,              % Activates advanced skin engine hooks (requires tcolorbox library 'skins')
    colback=gray!10,       % Light gray canvas fill background (10% opacity black over white canvas)
    colframe=black!20,     % Muted gray border line frame color (20% opacity black)
    boxrule=0.4pt,         % Ultra-thin, crisp frame line boundary thickness
    arc=1mm,               % Sets a subtle, modern 1mm corner radius curve
    % Tight geometric padding rules to prevent large text blocks from forcing column overflows:
    left=2mm,              % Internal text padding buffer from the left border
    right=2mm,             % Internal text padding buffer from the right border
    top=1mm,               % Internal text padding buffer beneath the title separator rule
    bottom=1mm,            % Internal text padding buffer above the lower frame floor
    fonttitle=\bfseries,   % Sets the box title header typography font to boldface
    coltitle=black,        % Keeps the title string text plain black for contrast over gray background
    title={#1},            % Evaluates Positional Argument #1 as the structural header text
    % Elastic vertical paragraph spacing rules that stretch or shrink to prevent layout underfull breaks:
    before skip=3pt plus 2pt minus 1pt,
    after skip=3pt plus 2pt minus 1pt
  ]
  \small #2                % Reduces text size inside the box to \small for better scannability of dense text payloads
  \end{tcolorbox}%         % The trailing '%' character strips accidental trailing whitespaces from compilation output
}

% --- structured empirical finding blocks
\newcounter{finding}                   % Declares custom structural metric serialization register
\newcommand{\newfinding}{%             % Instantiates inline empirical milestone callouts
    \refstepcounter{finding}           % Increments register metrics sequentially and binds anchor links for \label
    \textbf{Finding \thefinding.}   % Prints stylized numeric marker index into current structural block stream
}

% =========================================================================
% REVIEWER INTERACTION & COLLABORATIVE DEBUG SWITCHES
% =========================================================================

\ifDEBUG
    % --- active draft editing marker functions
    \newcommand{\TODO}[1]{\hl{TODO: #1}} % Visual yellow structural highlight marker
    \newcommand{\CITE}[1]{\hl{CITE: #1}} % Visual yellow missing source tracking marker
    
    % --- colored co-author editorial discussion nodes
    \newcommand{\NMS}[1]{\textcolor{orange}{[NMS: #1]}} % Nicholas M. Synovic
\else
    % --- production compilation overrides (silences markers entirely)
    \newcommand{\TODO}[1]{}
    \newcommand{\CITE}[1]{}
    
    \newcommand{\NMS}[1]{}
\fi % Imports formatting overrides, extra packages, and layout tweaks

% =========================================================================
% CUSTOM LOGO MACROS
% =========================================================================

% Delays definition until document start to prevent package conflicts
\AtBeginDocument{%
    % Safely defines \BibTeX logo if not already defined by a loaded package
    \providecommand\BibTeX{{%
        \normalfont B\kern-0.5em{%   % Backspaces slightly after 'B'
            \scshape                 % Switches to small caps for 'ib'
            i\kern-0.25em b}\kern-0.8em\TeX}}} 
% Micro-spaces 'ib' and pulls 'TeX' close

% =========================================================================
% RIGHTS & COPYRIGHT MANAGEMENT (Suppressed while 'nonacm' is active)
% =========================================================================

\setcopyright{acmlicensed}          % Specifies the ACM licensing agreement type
\copyrightyear{2026}                % Sets the explicit copyright year
\acmYear{2026}                      % Sets the publishing year index
\acmDOI{XXXXXXX.XXXXXXX}            % Placeholder for the Digital Object Identifier string
\acmISBN{978-1-4503-XXXX-X/18/06}   % Placeholder for the standard publication ISBN

% Defines conference acronym, full metadata description, dates, and geographic location
\acmConference[Conference acronym 'XX]{Make sure to enter the correct
  conference title from your rights confirmation email}{June 03--05,
  2018}{Woodstock, NY}

% =========================================================================
% DOCUMENT CONTENT START
% =========================================================================

\begin{document}

\title{\textbf{INSERT TITLE}} % Formats the core paper title in bold face

% -------------------------------------------------------------------------
% AUTHOR METADATA
% -------------------------------------------------------------------------

\author{Nicholas M. Synovic}
\email{nsynovic@luc.edu}
\orcid{0000-0003-0413-4594} % Unique digital author identifier
\affiliation{
  \institution{Loyola University Chicago}
  \streetaddress{1032 W Sheridan Rd}
  \city{Chicago}
  \state{Illinois}
  \country{USA}
  \postcode{60660}
}

% Replaces the full author list in page running headers with a concise version to save space
\renewcommand{\shortauthors}{Synovic}

% -------------------------------------------------------------------------
% ABSTRACT SECTION
% -------------------------------------------------------------------------

% ABSTRACT: one self-contained, three-part narrative block (6–8 sentences total).
% HARD RULES: no citations and no raw statistical tables/numbers in the abstract.
\begin{abstract}
% Part 1 — Context & Problem (Sentences 1–3):
% Establish the broad SE domain and why it matters, then pivot to the core
% friction point or consequence. Function: context → problem.
Lorem ipsum dolor sit amet consectetur adipiscing elit. 
Quisque faucibus ex sapien vitae pellentesque sem placerat. 
In id cursus mi pretium tellus duis convallis. 

% Part 2 — Approach & Scope (Sentences 4–5):
% State "In this paper, we…" and summarize the design/method at one-sentence
% granularity, including evaluation scope (data size, system types, participants).
Tempus leo eu aenean sed diam urna tempor. 
Pulvinar vivamus fringilla lacus nec metus bibendum egestas. 

% Part 3 — Impact & Insights (Sentences 6–8):
% Deliver the headline qualitative/quantitative outcome, the supporting insight,
% and one concrete implication for practice or future research.
% Keep it high-level: no detailed numbers or statistics.
Iaculis massa nisl malesuada lacinia integer nunc posuere. 
Ut hendrerit semper vel class aptent taciti sociosqu. 
Ad litora torquent per conubia nostra inceptos himenaeos.
\end{abstract}

% =========================================================================
% ACM COMPUTING CLASSIFICATION SYSTEM (https://dl.acm.org/ccs)
% =========================================================================
\begin{CCSXML}
\end{CCSXML}

% The \ccsdesc macros render XML classifications into the document footer
\ccsdesc[500]{Software and its engineering}

% =========================================================================
% KEYWORDS INDEXING
% =========================================================================

\keywords{Replace me}

\maketitle

% =========================================================================
% SECTION: INTRODUCTION
% Function: Sell the paper. Six-paragraph arc; each block has ONE rhetorical job.
% Arc (one block each): broad context/hook → narrowed gap → goal & core insight
%   → contributions (bulleted) → significance/scope → roadmap.
% Rules: 3–7 sentences per paragraph; explicit transitions between blocks
%   ("However", "As a result"); no unquantified hype ("novel", "significant").
% =========================================================================
% RULE: every figure must be cited in-text via \Cref{fig:...} (e.g. \Cref{fig:template_figure}).
% =========================================================================
% TEMPLATE FIGURE PLACEHOLDER
% Uses a standard boxed paragraph container to stand in for missing visual assets.
% =========================================================================

\begin{figure}[h] % [h] requests placement "here" inline with the surrounding text stream
    \centering    % Horizontally centers the figure artifact within the column/page margins
    
    % \fbox draws a visible frame border around its inner content block
    \fbox{
        % \parbox structure creates an explicit, fixed-size paragraph boundary container
        % Arguments format: [inner-alignment][height][content-alignment]{width}
        \parbox[c][2in][c]{3in}{
            \centering Missing Figure % Centered placeholder text string inside the bounding box
        }
    }
    
    % Core descriptive label rendered immediately beneath the figure box frame
    \caption{Template figure placeholder}
    
    % Unique cross-referencing anchor tag; reference this asset elsewhere via \Cref{fig:template_figure}
    \label{fig:template_figure}
\end{figure}
\section{Introduction}
\label{sec:introduction}

% --- PARAGRAPH 1: CONTEXT AND HOOK ---
% Function: Intro paragraph 1.
% Target: 3–6 sentences.
% Content: Establish the broad software engineering domain, why it matters (reliability, cost, UX, scientific impact), and a hook that shows the domain is non‑trivial. No detailed technicalities here.
Lorem ipsum dolor sit amet consectetur adipiscing elit.
Quisque faucibus ex sapien vitae pellentesque sem placerat. 
In id cursus mi pretium tellus duis convallis.
Tempus leo eu aenean sed diam urna tempor. 
Pulvinar vivamus fringilla lacus nec metus bibendum egestas.
Iaculis massa nisl malesuada lacinia integer nunc posuere. 

% --- PARAGRAPH 2: NARROWED PROBLEM AND GAP ---
% Function: Intro paragraph 2.
% Target: 4–7 sentences.
% Content: Summarize current practice/prior work, then pivot with a clear ‘However/Yet/Despite this’ sentence that states the precise gap or friction point. End by emphasizing why this gap is consequential.
Lorem ipsum dolor sit amet consectetur adipiscing elit. 
Quisque faucibus ex sapien vitae pellentesque sem placerat.
In id cursus mi pretium tellus duis convallis. 
Tempus leo eu aenean sed diam urna tempor.
Pulvinar vivamus fringilla lacus nec metus bibendum egestas. 
Iaculis massa nisl malesuada lacinia integer nunc posuere.
Ut hendrerit semper vel class aptent taciti sociosqu. 

% --- PARAGRAPH 3: GOAL AND APPROACH (THE CORE INSIGHT) ---
% Function: Intro paragraph 3.
% Target: 3–6 sentences.
% Content: Present the overall research goal and high‑level approach (system, study, framework). Include plain‑language intuition about why your approach should address the stated gap.
Lorem ipsum dolor sit amet consectetur adipiscing elit. 
Quisque faucibus ex sapien vitae pellentesque sem placerat.
In id cursus mi pretium tellus duis convallis. 
Tempus leo eu aenean sed diam urna tempor.
Pulvinar vivamus fringilla lacus nec metus bibendum egestas. 
Iaculis massa nisl malesuada lacinia integer nunc posuere.

% --- PARAGRAPH 4: CONTRIBUTIONS ---
% Structure: Use 3–4 bullets, each describing one type of contribution (artifact, empirical study, insight, open science). 
% Formatting Rule: Keep each item one sentence plus a short clarifying clause, avoiding vague phrases (‘novel’, ‘significant’) without specifics.
\myparagraph{Contributions}
This paper makes the following core contributions:
\begin{itemize}[leftmargin=16pt, rightmargin=5pt, topsep=0pt]
    % Contribution 1: The Artifact (Tool / Framework / Dataset)
    \item Lorem ipsum dolor sit amet consectetur adipiscing elit.
    % Contribution 2: The Empirical Scale & Matrix Analysis
    \item Lorem ipsum dolor sit amet consectetur adipiscing elit.
    % Contribution 3: The Findings, Open Science & Replication Package
    \item Lorem ipsum dolor sit amet consectetur adipiscing elit.
\end{itemize}

% --- PARAGRAPH 5: SIGNIFICANCE & SCOPE ---
% Function: Intro paragraph 5 — "headline results + scope" preview (required).
% Target: 3–5 sentences.
% Content: Provide 1–2 headline quantitative or qualitative outcomes, plus a brief statement of scope/assumptions (where it works, where it might not). Do not repeat details from the Results; this is a teaser, not the findings.
\myparagraph{Significance}
Lorem ipsum dolor sit amet consectetur adipiscing elit. 
Quisque faucibus ex sapien vitae pellentesque sem placerat.
In id cursus mi pretium tellus duis convallis. 
Tempus leo eu aenean sed diam urna tempor.
Pulvinar vivamus fringilla lacus nec metus bibendum egestas. 

% --- PARAGRAPH 6: PAPER ROADMAP ---
% Function: Manuscript outline (Intro paragraph 6).
% Target: 2–4 sentences.
% Content: Walk the reader through the paper in document order so the roadmap
%   doubles as a structural map. The actual order is:
%   Background (sec:background) → Research Questions (sec:research_questions)
%   → RQ1 (sec:rq1_evaluation) → RQ2 (sec:rq2_evaluation)
%   → RQ3 (sec:rq3_evaluation) → Discussion (sec:discussion)
%   → Threats (sec:threats) → Conclusion (sec:conclusion).
%   NOTE: there is no standalone "Methodology" section — methodology is embedded
%   inside each RQ section. Use lowercase section names + \Cref macros; never
%   hardcode numbers.
\myparagraph{Paper Outline}
This paper is organized as follows: 
\Cref{sec:background} provides essential background and reviews related work, and 
\Cref{sec:research_questions} formalizes the research questions.
\Cref{sec:rq1_evaluation}, \Cref{sec:rq2_evaluation}, and \Cref{sec:rq3_evaluation} each present the methodology and empirical findings for one research question.
The paper concludes in \Cref{sec:discussion} with a discussion of implications for the software engineering community and future work, followed by an examination of potential threats to validity in \Cref{sec:threats} and closing remarks in \Cref{sec:conclusion}.

% =========================================================================
% SECTION: BACKGROUND AND RELATED WORK
% Narrative Arc: "What world are we in?" -> "What fails here?" -> "How do we fix it?"
% =========================================================================
\section{Background and Related Work}
\label{sec:background}

% --- MAJOR SIGNPOSTING & RE-CONTEXTUALIZATION ---
% Entry paragraph (3–5 sentences): re‑state the problem space with more technical depth than the Introduction and briefly signpost the three subsections (domain overview, existing approaches, conceptual foundation).
Lorem ipsum dolor sit amet consectetur adipiscing elit. 
Quisque faucibus ex sapien vitae pellentesque sem placerat.
In id cursus mi pretium tellus duis convallis. 
Tempus leo eu aenean sed diam urna tempor.
Pulvinar vivamus fringilla lacus nec metus bibendum egestas. 

% =========================================================================
% SUBSECTION 1: THE DOMAIN OVERVIEW (Domain overview)
% =========================================================================
\subsection{Subsection 1}
\label{subsec:subsection-1}

% Paragraph 1 – Domain definitions and workflows: 
% Technical definitions and standard workflows. Clarify entities (projects, pipelines, registries), typical processes, and key terminology.
Lorem ipsum dolor sit amet consectetur adipiscing elit. 
Quisque faucibus ex sapien vitae pellentesque sem placerat.
In id cursus mi pretium tellus duis convallis. 
Tempus leo eu aenean sed diam urna tempor.
Pulvinar vivamus fringilla lacus nec metus bibendum egestas. 
Iaculis massa nisl malesuada lacinia integer nunc posuere.

% Paragraph 2 – Cause‑and‑effect linkages: 
% Explain how practices and constraints in this domain lead to specific challenges (e.g., fragility, reproducibility issues). Use ‘As a result/Consequently’ to make causal links explicit.
Lorem ipsum dolor sit amet consectetur adipiscing elit.
Quisque faucibus ex sapien vitae pellentesque sem placerat. 
In id cursus mi pretium tellus duis convallis.
Tempus leo eu aenean sed diam urna tempor. 
Pulvinar vivamus fringilla lacus nec metus bibendum egestas.
Iaculis massa nisl malesuada lacinia integer nunc posuere. 

% Paragraph 3 – Transition to existing approaches: 
% Close the subsection by hinting at what practitioners or researchers currently do to cope with these challenges, setting up Subsection 2.
Lorem ipsum dolor sit amet consectetur adipiscing elit. 
Quisque faucibus ex sapien vitae pellentesque sem placerat.
In id cursus mi pretium tellus duis convallis. 
Tempus leo eu aenean sed diam urna tempor.
Pulvinar vivamus fringilla lacus nec metus bibendum egestas. 
Iaculis massa nisl malesuada lacinia integer nunc posuere.

% =========================================================================
% SUBSECTION 2: EXISTING APPROACHES & PRIOR WORK
% =========================================================================
\subsection{Subsection 2}
\label{subsec:subsection-2}

% Paragraph 1 – Subsection entry and literature grouping: 
% Target: 3–5 sentences. 
% Content: Explicitly tie back to Subsection 1 and describe high‑level categories of existing approaches (e.g., tooling, guidelines, empirical studies). Group literature by theme, not chronologically.
Lorem ipsum dolor sit amet consectetur adipiscing elit. 
Quisque faucibus ex sapien vitae pellentesque sem placerat.
In id cursus mi pretium tellus duis convallis. 
Tempus leo eu aenean sed diam urna tempor.
Pulvinar vivamus fringilla lacus nec metus bibendum egestas. 
Iaculis massa nisl malesuada lacinia integer nunc posuere.

% Paragraph 2 – Contrastive analysis and limitations: 
% Content: Describe what these approaches achieve and where they fall short. Use contrast markers (‘However’, ‘In contrast’) and connect limitations to your gap.
Lorem ipsum dolor sit amet consectetur adipiscing elit.
Quisque faucibus ex sapien vitae pellentesque sem placerat. 
In id cursus mi pretium tellus duis convallis.
Tempus leo eu aenean sed diam urna tempor. 
Pulvinar vivamus fringilla lacus nec metus bibendum egestas.
Iaculis massa nisl malesuada lacinia integer nunc posuere. 

% Paragraph 3 – Exit transition to conceptual lens: 
% Content: Summarize the overall shortcomings of existing work and point forward to the conceptual foundation you will use to reason about them (Subsection 3).
Lorem ipsum dolor sit amet consectetur adipiscing elit. 
Quisque faucibus ex sapien vitae pellentesque sem placerat.
In id cursus mi pretium tellus duis convallis. 
Tempus leo eu aenean sed diam urna tempor.
Pulvinar vivamus fringilla lacus nec metus bibendum egestas. 
Iaculis massa nisl malesuada lacinia integer nunc posuere.

% =========================================================================
% SUBSECTION 3: CONCEPTUAL FOUNDATION
% =========================================================================
\subsection{Subsection 3}
\label{subsec:subsection-3}

% Paragraph 1 – Subsection entry and theoretical framing: 
% Content: Introduce the conceptual/theoretical lens (e.g., socio‑technical, diffusion, reliability) that underpins your study. Motivate why this lens is appropriate given the domain and prior work.
Lorem ipsum dolor sit amet consectetur adipiscing elit. 
Quisque faucibus ex sapien vitae pellentesque sem placerat.
In id cursus mi pretium tellus duis convallis. 
Tempus leo eu aenean sed diam urna tempor.
Pulvinar vivamus fringilla lacus nec metus bibendum egestas. 
Iaculis massa nisl malesuada lacinia integer nunc posuere.

% Paragraph 2 – Technical hook into methodology: 
% Content: Tie the chosen lens directly to your upcoming methodology, hinting at key constructs, relationships, or hypotheses that will be operationalized in the Methods and RQ sections.
Lorem ipsum dolor sit amet consectetur adipiscing elit. 
Quisque faucibus ex sapien vitae pellentesque sem placerat.
In id cursus mi pretium tellus duis convallis. 
Tempus leo eu aenean sed diam urna tempor.
Pulvinar vivamus fringilla lacus nec metus bibendum egestas. 
Iaculis massa nisl malesuada lacinia integer nunc posuere.

% =========================================================================
% SECTION: RESEARCH QUESTIONS (The Structural Bridge)
% =========================================================================
\section{Research Questions}
\label{sec:research_questions}

% --- INTRODUCTORY SIGNPOSTING & GAP ANCHORING ---
% Target: 1–3 sentences. 
% Content: Remind the reader of the key gap from Background and state that you now formalize the evaluation goals as research questions.
Lorem ipsum dolor sit amet consectetur adipiscing elit. 
Quisque faucibus ex sapien vitae pellentesque sem placerat.
In id cursus mi pretium tellus duis convallis. 

% --- THE FORMAL RQ LIST ---
% Function: The formal, enumerated research questions (the structural bridge).
% Content: List 3–4 concise RQs, each isolating ONE dimension (e.g., prevalence,
%   patterns, impact). Keep each RQ to a single sentence prefaced by a bolded
%   dimensional anchor (e.g., [Prevalence]).
% Format: per the style guide, the colon sits OUTSIDE the bold: \textbf{RQ\arabic*}:
\noindent We address the following research questions:
\begin{enumerate}[label=\textbf{RQ\arabic*}:, leftmargin=*]
    \item \textbf{[Dimension 1, e.g., Prevalence].} Lorem ipsum dolor sit amet consectetur adipiscing elit?
    \item \textbf{[Dimension 2, e.g., Patterns].} Lorem ipsum dolor sit amet consectetur adipiscing elit?
    \item \textbf{[Dimension 3, e.g., Impact].} Lorem ipsum dolor sit amet consectetur adipiscing elit?
\end{enumerate}

% --- MAPPING PARAGRAPH ---
% Target: 2–3 sentences. 
% Content: Map each RQ to its evaluation section (RQ1, RQ2, RQ3) and briefly mention which data and methods answer which question. This proves the downstream structure aligns with stated goals.
Lorem ipsum dolor sit amet consectetur adipiscing elit. 
Quisque faucibus ex sapien vitae pellentesque sem placerat.
In id cursus mi pretium tellus duis convallis. 

% =========================================================================
% SECTION: EVALUATION OF RESEARCH QUESTION 1 (RQ1)
% RQ1 block arc: overview & motivation → RQ1‑specific methodology → RQ1 findings → formal answer box → pointer to Discussion.
% =========================================================================
\section{Evaluation of RQ1: [Dimension Title, e.g., Prevalence]}
\label{sec:rq1_evaluation}

% --- RQ1 OVERVIEW PARAGRAPH ---
% Target: 6–8 sentences. 
% Content: Restate RQ1 verbatim, explain its role relative to other RQs (baseline, patterns, impact), and describe why this dimension matters. End with a sentence that introduces the upcoming methodology subsection.
Lorem ipsum dolor sit amet consectetur adipiscing elit. 
Quisque faucibus ex sapien vitae pellentesque sem placerat. 
In id cursus mi pretium tellus duis convallis.
Tempus leo eu aenean sed diam urna tempor. 
Pulvinar vivamus fringilla lacus nec metus bibendum egestas.
Iaculis massa nisl malesuada lacinia integer nunc posuere. 
Ut hendrerit semper vel class aptent taciti sociosqu. 
Ad litora torquent per conubia nostra inceptos himenaeos.

% --- SUBSECTION: METHODOLOGY FOR RQ1 ---
\subsection{RQ1 Methodology and Operationalization}
\label{subsec:rq1_methodology}
% RULE: cite every figure in-text via \Cref{fig:...}; never rely on placement alone.
% =========================================================================
% TEMPLATE FIGURE PLACEHOLDER
% Uses a standard boxed paragraph container to stand in for missing visual assets.
% =========================================================================

\begin{figure}[h] % [h] requests placement "here" inline with the surrounding text stream
    \centering    % Horizontally centers the figure artifact within the column/page margins
    
    % \fbox draws a visible frame border around its inner content block
    \fbox{
        % \parbox structure creates an explicit, fixed-size paragraph boundary container
        % Arguments format: [inner-alignment][height][content-alignment]{width}
        \parbox[c][2in][c]{3in}{
            \centering Missing Figure % Centered placeholder text string inside the bounding box
        }
    }
    
    % Core descriptive label rendered immediately beneath the figure box frame
    \caption{Template figure placeholder}
    
    % Unique cross-referencing anchor tag; reference this asset elsewhere via \Cref{fig:template_figure}
    \label{fig:template_figure}
\end{figure}

% Paragraph 1 – Design & rationale: 
% Explain the study design (experiment, case study, mining, survey) and why it is appropriate for RQ1.
Lorem ipsum dolor sit amet consectetur adipiscing elit. 
Quisque faucibus ex sapien vitae pellentesque sem placerat. 
In id cursus mi pretium tellus duis convallis. 
Tempus leo eu aenean sed diam urna tempor.

% Paragraph 2 – Data sources & sampling: 
% Describe datasets, projects, participants, and inclusion/exclusion rules. Provide enough detail for replication.
Lorem ipsum dolor sit amet consectetur adipiscing elit. 
Quisque faucibus ex sapien vitae pellentesque sem placerat. 
In id cursus mi pretium tellus duis convallis. 
Tempus leo eu aenean sed diam urna tempor.

% Paragraph 3 – Procedure: 
% Outline the step‑by‑step execution (pipelines run, tasks assigned, scripts used), emphasizing chronology and reproducibility.
Lorem ipsum dolor sit amet consectetur adipiscing elit. 
Quisque faucibus ex sapien vitae pellentesque sem placerat. 
In id cursus mi pretium tellus duis convallis. 
Tempus leo eu aenean sed diam urna tempor.

% Paragraph 4 – Measures & analysis plan: 
% Enumerate metrics and constructs, describe how they are calculated, and specify statistical/analytical methods used to answer RQ1.
Lorem ipsum dolor sit amet consectetur adipiscing elit. 
Quisque faucibus ex sapien vitae pellentesque sem placerat. 
In id cursus mi pretium tellus duis convallis. 
Tempus leo eu aenean sed diam urna tempor.

% Threats to Validity for RQ1:
% Short paragraph. Note key RQ1‑specific threats (e.g., confounding, sampling, measurement) and point to the global Threats section for deeper analysis.
\myparagraph{Threats to Validity for RQ1}
Lorem ipsum dolor sit amet consectetur adipiscing elit. 
Quisque faucibus ex sapien vitae pellentesque sem placerat. 
In id cursus mi pretium tellus duis convallis. 

% --- SUBSECTION: EMPIRICAL FINDINGS FOR RQ1 ---
\subsection{RQ1 Empirical Findings}
\label{subsec:rq1_findings}

% Paragraph 1 – Overview and preprocessing: 
% Describe the subset of data used for RQ1, any preprocessing, exclusions, and basic descriptive stats (e.g., counts, ranges).
Lorem ipsum dolor sit amet consectetur adipiscing elit. 
Quisque faucibus ex sapien vitae pellentesque sem placerat. 
In id cursus mi pretium tellus duis convallis. 
Tempus leo eu aenean sed diam urna tempor.

% Paragraphs 2–3 – Main findings with figures/tables: 
% Present the core results in logical order (e.g., primary metrics, distributions, comparative stats). Every figure/table referenced must be explicitly cited in the text (‘As shown in Figure X…’). No interpretation beyond factual patterns.
Lorem ipsum dolor sit amet consectetur adipiscing elit. 
Quisque faucibus ex sapien vitae pellentesque sem placerat.
In id cursus mi pretium tellus duis convallis. 
Tempus leo eu aenean sed diam urna tempor.
Pulvinar vivamus fringilla lacus nec metus bibendum egestas. 
Iaculis massa nisl malesuada lacinia integer nunc posuere.
Ut hendrerit semper vel class aptent taciti sociosqu. 

% Paragraph 4 – Trends, patterns, and null results: 
% Summarize observed trends, highlight interesting edge cases, and explicitly mention negative or null results. Stay descriptive; save interpretation for the Discussion.
Lorem ipsum dolor sit amet consectetur adipiscing elit. 
Quisque faucibus ex sapien vitae pellentesque sem placerat. 
In id cursus mi pretium tellus duis convallis. 
Tempus leo eu aenean sed diam urna tempor.

% Formal ‘Answer to RQ1’ box: 
% 2–3 sentences of neutral synthesis directly answering the RQ text, with no speculation.
% =========================================================================
% EMPIRICAL FINDING DISPLAY BOX
% DEPENDENCY WARNING: This structural container relies directly on the 
% \newfinding command (defined in typeset/misc.tex). 
% Calling \newfinding inside this block:
%   1. Increments the central sequential 'finding' counter register.
%   2. Globally registers a target link for downstream \label and \ref calls.
%   3. Injects the bold string "Finding X." directly into the box canvas stream.
% =========================================================================

\begin{tcolorbox}[
    colback=yellow!30,    % Warm yellow alert canvas interior shading tint
    colframe=black,       % High-contrast solid black border outline
    sharp corners,        % Disables rounded bounding radii to mimic a standard \fcolorbox
    boxrule=0.5pt,        % Ultra-clean frame line width boundary scale
    % Comprehensive layout padding overrides to prevent text edge crowding:
    left=2mm, 
    right=2mm, 
    top=2mm, 
    bottom=5pt, 
    width=0.98\linewidth  % Consistently centers and bounds the box safely inside column margins
]
    % Note: Typically, you will invoke only ONE unique \newfinding call per box container.
    % Multiple consecutive calls here are used to demonstrate sequence step increments.
    \newfinding In id cursus mi pretium tellus duis convallis.
    
    \newfinding Pulvinar vivamus fringilla lacus nec metus bibendum.
    
    \newfinding Iaculis massa nisl malesuada lacinia integer.
    \label{finding: template-finding}
\end{tcolorbox}

% Pointer sentence to Discussion: 
% Single sentence that forwards the reader to the Discussion for qualitative interpretation and implications.
Lorem ipsum dolor sit amet consectetur adipiscing elit.

% =========================================================================
% SECTION: EVALUATION OF RESEARCH QUESTION 2 (RQ2)
% RQ2 block arc: overview & motivation → RQ2‑specific methodology → RQ2 findings → formal answer box → pointer to Discussion.
% =========================================================================
\section{Evaluation of RQ2: [Dimension Title, e.g., Patterns]}
\label{sec:rq2_evaluation}

% --- RQ2 OVERVIEW PARAGRAPH ---
% Target: 6–8 sentences.
% Content: Restate RQ2 verbatim, explain its role relative to the other RQs (it
%   builds on RQ1's baseline by characterizing patterns/mechanics), and explain
%   why this dimension matters. End with a sentence introducing the methodology.
Lorem ipsum dolor sit amet consectetur adipiscing elit.
Quisque faucibus ex sapien vitae pellentesque sem placerat. 
In id cursus mi pretium tellus duis convallis.
Tempus leo eu aenean sed diam urna tempor. 
Pulvinar vivamus fringilla lacus nec metus bibendum egestas.
Iaculis massa nisl malesuada lacinia integer nunc posuere. 
Ut hendrerit semper vel class aptent taciti sociosqu. 
Ad litora torquent per conubia nostra inceptos himenaeos.

% --- SUBSECTION: METHODOLOGY FOR RQ2 ---
\subsection{RQ2 Methodology and Categorization Pipeline}
\label{subsec:rq2_methodology}
% RULE: cite every figure in-text via \Cref{fig:...}; never rely on placement alone.
% =========================================================================
% TEMPLATE FIGURE PLACEHOLDER
% Uses a standard boxed paragraph container to stand in for missing visual assets.
% =========================================================================

\begin{figure}[h] % [h] requests placement "here" inline with the surrounding text stream
    \centering    % Horizontally centers the figure artifact within the column/page margins
    
    % \fbox draws a visible frame border around its inner content block
    \fbox{
        % \parbox structure creates an explicit, fixed-size paragraph boundary container
        % Arguments format: [inner-alignment][height][content-alignment]{width}
        \parbox[c][2in][c]{3in}{
            \centering Missing Figure % Centered placeholder text string inside the bounding box
        }
    }
    
    % Core descriptive label rendered immediately beneath the figure box frame
    \caption{Template figure placeholder}
    
    % Unique cross-referencing anchor tag; reference this asset elsewhere via \Cref{fig:template_figure}
    \label{fig:template_figure}
\end{figure}

% Paragraph 1 – Design & rationale:
% Explain the study design (experiment, case study, mining, survey) and why it
% is appropriate for RQ2.
Lorem ipsum dolor sit amet consectetur adipiscing elit.
Quisque faucibus ex sapien vitae pellentesque sem placerat. 
In id cursus mi pretium tellus duis convallis. 
Tempus leo eu aenean sed diam urna tempor.

% Paragraph 2 – Data sources & sampling:
% Describe datasets, projects, participants, and inclusion/exclusion rules.
% Provide enough detail for replication.
Lorem ipsum dolor sit amet consectetur adipiscing elit.
Quisque faucibus ex sapien vitae pellentesque sem placerat. 
In id cursus mi pretium tellus duis convallis. 
Tempus leo eu aenean sed diam urna tempor.

% Paragraph 3 – Procedure:
% Outline the step-by-step execution (pipelines run, tasks assigned, scripts
% used), emphasizing chronology and reproducibility.
Lorem ipsum dolor sit amet consectetur adipiscing elit.
Quisque faucibus ex sapien vitae pellentesque sem placerat. 
In id cursus mi pretium tellus duis convallis. 
Tempus leo eu aenean sed diam urna tempor.

% Paragraph 4 – Measures & analysis plan:
% Enumerate metrics and constructs, describe how they are calculated, and specify
% the statistical/analytical methods used to answer RQ2.
Lorem ipsum dolor sit amet consectetur adipiscing elit.
Quisque faucibus ex sapien vitae pellentesque sem placerat. 
In id cursus mi pretium tellus duis convallis. 
Tempus leo eu aenean sed diam urna tempor.

% Threats to Validity for RQ2:
% Short paragraph (2–3 sentences). Note key RQ2-specific threats (e.g.,
% confounding, sampling, measurement) and point to the global Threats section.
\myparagraph{Threats to Validity for RQ2}
Lorem ipsum dolor sit amet consectetur adipiscing elit.
Quisque faucibus ex sapien vitae pellentesque sem placerat. 
In id cursus mi pretium tellus duis convallis. 

% --- SUBSECTION: EMPIRICAL FINDINGS FOR RQ2 ---
\subsection{RQ2 Empirical Findings}
\label{subsec:rq2_findings}

% Paragraph 1 – Overview and preprocessing:
% Describe the subset of data used for RQ2, any preprocessing, exclusions, and
% basic descriptive stats (e.g., counts, ranges).
Lorem ipsum dolor sit amet consectetur adipiscing elit.
Quisque faucibus ex sapien vitae pellentesque sem placerat. 
In id cursus mi pretium tellus duis convallis. 
Tempus leo eu aenean sed diam urna tempor.

% Paragraphs 2–3 – Main findings with figures/tables:
% Present the core results in logical order. Every figure/table referenced must
% be explicitly cited in the text ("As shown in \Cref{fig:...}"). Stay strictly
% descriptive; no interpretation beyond factual patterns.
Lorem ipsum dolor sit amet consectetur adipiscing elit. 
Quisque faucibus ex sapien vitae pellentesque sem placerat.
In id cursus mi pretium tellus duis convallis. 
Tempus leo eu aenean sed diam urna tempor.
Pulvinar vivamus fringilla lacus nec metus bibendum egestas. 
Iaculis massa nisl malesuada lacinia integer nunc posuere.
Ut hendrerit semper vel class aptent taciti sociosqu. 

% Paragraph 4 – Trends, patterns, and null results:
% Summarize observed trends, highlight interesting edge cases, and explicitly
% mention negative or null results. Stay descriptive; save interpretation for the
% Discussion.
Lorem ipsum dolor sit amet consectetur adipiscing elit.
Quisque faucibus ex sapien vitae pellentesque sem placerat. 
In id cursus mi pretium tellus duis convallis. 
Tempus leo eu aenean sed diam urna tempor. 

% Formal 'Answer to RQ2' box:
% 2–3 sentences of neutral synthesis directly answering the RQ text, no speculation.
% =========================================================================
% EMPIRICAL FINDING DISPLAY BOX
% DEPENDENCY WARNING: This structural container relies directly on the 
% \newfinding command (defined in typeset/misc.tex). 
% Calling \newfinding inside this block:
%   1. Increments the central sequential 'finding' counter register.
%   2. Globally registers a target link for downstream \label and \ref calls.
%   3. Injects the bold string "Finding X." directly into the box canvas stream.
% =========================================================================

\begin{tcolorbox}[
    colback=yellow!30,    % Warm yellow alert canvas interior shading tint
    colframe=black,       % High-contrast solid black border outline
    sharp corners,        % Disables rounded bounding radii to mimic a standard \fcolorbox
    boxrule=0.5pt,        % Ultra-clean frame line width boundary scale
    % Comprehensive layout padding overrides to prevent text edge crowding:
    left=2mm, 
    right=2mm, 
    top=2mm, 
    bottom=5pt, 
    width=0.98\linewidth  % Consistently centers and bounds the box safely inside column margins
]
    % Note: Typically, you will invoke only ONE unique \newfinding call per box container.
    % Multiple consecutive calls here are used to demonstrate sequence step increments.
    \newfinding In id cursus mi pretium tellus duis convallis.
    
    \newfinding Pulvinar vivamus fringilla lacus nec metus bibendum.
    
    \newfinding Iaculis massa nisl malesuada lacinia integer.
    \label{finding: template-finding}
\end{tcolorbox}

% Pointer sentence to Discussion:
% Single sentence forwarding the reader to the Discussion for qualitative
% interpretation and implications.
Lorem ipsum dolor sit amet consectetur adipiscing elit.

% =========================================================================
% SECTION: EVALUATION OF RESEARCH QUESTION 3 (RQ3)
% RQ3 block arc: overview & motivation → RQ3‑specific methodology → RQ3 findings → formal answer box → pointer to Discussion.
% =========================================================================
\section{Evaluation of RQ3: [Dimension Title, e.g., Impact]}
\label{sec:rq3_evaluation}

% --- RQ3 OVERVIEW PARAGRAPH ---
% Target: 6–8 sentences.
% Content: Restate RQ3 verbatim, explain its role relative to the other RQs (it
%   typically caps the arc by assessing impact/consequence), and explain why this
%   dimension matters. End with a sentence introducing the methodology subsection.
Lorem ipsum dolor sit amet consectetur adipiscing elit.
Quisque faucibus ex sapien vitae pellentesque sem placerat. 
In id cursus mi pretium tellus duis convallis.
Tempus leo eu aenean sed diam urna tempor. 
Pulvinar vivamus fringilla lacus nec metus bibendum egestas.
Iaculis massa nisl malesuada lacinia integer nunc posuere. 
Ut hendrerit semper vel class aptent taciti sociosqu. 
Ad litora torquent per conubia nostra inceptos himenaeos.

% --- SUBSECTION: METHODOLOGY FOR RQ3 ---
\subsection{RQ3 Methodology and Analytical Matrix}
\label{subsec:rq3_methodology}
% RULE: cite every figure in-text via \Cref{fig:...}; never rely on placement alone.
% =========================================================================
% TEMPLATE FIGURE PLACEHOLDER
% Uses a standard boxed paragraph container to stand in for missing visual assets.
% =========================================================================

\begin{figure}[h] % [h] requests placement "here" inline with the surrounding text stream
    \centering    % Horizontally centers the figure artifact within the column/page margins
    
    % \fbox draws a visible frame border around its inner content block
    \fbox{
        % \parbox structure creates an explicit, fixed-size paragraph boundary container
        % Arguments format: [inner-alignment][height][content-alignment]{width}
        \parbox[c][2in][c]{3in}{
            \centering Missing Figure % Centered placeholder text string inside the bounding box
        }
    }
    
    % Core descriptive label rendered immediately beneath the figure box frame
    \caption{Template figure placeholder}
    
    % Unique cross-referencing anchor tag; reference this asset elsewhere via \Cref{fig:template_figure}
    \label{fig:template_figure}
\end{figure}

% Paragraph 1 – Design & rationale:
% Explain the study design (experiment, case study, mining, survey) and why it
% is appropriate for RQ3.
Lorem ipsum dolor sit amet consectetur adipiscing elit.
Quisque faucibus ex sapien vitae pellentesque sem placerat. 
In id cursus mi pretium tellus duis convallis. 
Tempus leo eu aenean sed diam urna tempor.

% Paragraph 2 – Data sources & sampling:
% Describe datasets, projects, participants, and inclusion/exclusion rules.
% Provide enough detail for replication.
Lorem ipsum dolor sit amet consectetur adipiscing elit.
Quisque faucibus ex sapien vitae pellentesque sem placerat. 
In id cursus mi pretium tellus duis convallis. 
Tempus leo eu aenean sed diam urna tempor. 

% Paragraph 3 – Procedure:
% Outline the step-by-step execution (pipelines run, tasks assigned, scripts
% used), emphasizing chronology and reproducibility.
Lorem ipsum dolor sit amet consectetur adipiscing elit.
Quisque faucibus ex sapien vitae pellentesque sem placerat. 
In id cursus mi pretium tellus duis convallis. 
Tempus leo eu aenean sed diam urna tempor.

% Paragraph 4 – Measures & analysis plan:
% Enumerate metrics and constructs, describe how they are calculated, and specify
% the statistical/analytical methods used to answer RQ3.
Lorem ipsum dolor sit amet consectetur adipiscing elit.
Quisque faucibus ex sapien vitae pellentesque sem placerat. 
In id cursus mi pretium tellus duis convallis. 
Tempus leo eu aenean sed diam urna tempor.

% Threats to Validity for RQ3:
% Short paragraph (2–3 sentences). Note key RQ3-specific threats (e.g.,
% confounding, sampling, measurement) and point to the global Threats section.
\myparagraph{Threats to Validity for RQ3}
Lorem ipsum dolor sit amet consectetur adipiscing elit.
Quisque faucibus ex sapien vitae pellentesque sem placerat. 
In id cursus mi pretium tellus duis convallis. 

% --- SUBSECTION: EMPIRICAL FINDINGS FOR RQ3 ---
\subsection{RQ3 Empirical Findings}
\label{subsec:rq3_findings}

% Paragraph 1 – Overview and preprocessing:
% Describe the subset of data used for RQ3, any preprocessing, exclusions, and
% basic descriptive stats (e.g., counts, ranges).
Lorem ipsum dolor sit amet consectetur adipiscing elit.
Quisque faucibus ex sapien vitae pellentesque sem placerat. 
In id cursus mi pretium tellus duis convallis. 
Tempus leo eu aenean sed diam urna tempor.

% Paragraphs 2–3 – Main findings with figures/tables:
% Present the core results in logical order. Every figure/table referenced must
% be explicitly cited in the text ("As shown in \Cref{fig:...}"). Stay strictly
% descriptive; no interpretation beyond factual patterns.
Lorem ipsum dolor sit amet consectetur adipiscing elit. 
Quisque faucibus ex sapien vitae pellentesque sem placerat.
In id cursus mi pretium tellus duis convallis. 
Tempus leo eu aenean sed diam urna tempor.
Pulvinar vivamus fringilla lacus nec metus bibendum egestas. 
Iaculis massa nisl malesuada lacinia integer nunc posuere.
Ut hendrerit semper vel class aptent taciti sociosqu. 

% Paragraph 4 – Trends, patterns, and null results:
% Summarize observed trends, highlight interesting edge cases, and explicitly
% mention negative or null results. Stay descriptive; save interpretation for the
% Discussion.
Lorem ipsum dolor sit amet consectetur adipiscing elit.
Quisque faucibus ex sapien vitae pellentesque sem placerat. 
In id cursus mi pretium tellus duis convallis. 
Tempus leo eu aenean sed diam urna tempor. 

% Formal 'Answer to RQ3' box:
% 2–3 sentences of neutral synthesis directly answering the RQ text, no speculation.
% =========================================================================
% EMPIRICAL FINDING DISPLAY BOX
% DEPENDENCY WARNING: This structural container relies directly on the 
% \newfinding command (defined in typeset/misc.tex). 
% Calling \newfinding inside this block:
%   1. Increments the central sequential 'finding' counter register.
%   2. Globally registers a target link for downstream \label and \ref calls.
%   3. Injects the bold string "Finding X." directly into the box canvas stream.
% =========================================================================

\begin{tcolorbox}[
    colback=yellow!30,    % Warm yellow alert canvas interior shading tint
    colframe=black,       % High-contrast solid black border outline
    sharp corners,        % Disables rounded bounding radii to mimic a standard \fcolorbox
    boxrule=0.5pt,        % Ultra-clean frame line width boundary scale
    % Comprehensive layout padding overrides to prevent text edge crowding:
    left=2mm, 
    right=2mm, 
    top=2mm, 
    bottom=5pt, 
    width=0.98\linewidth  % Consistently centers and bounds the box safely inside column margins
]
    % Note: Typically, you will invoke only ONE unique \newfinding call per box container.
    % Multiple consecutive calls here are used to demonstrate sequence step increments.
    \newfinding In id cursus mi pretium tellus duis convallis.
    
    \newfinding Pulvinar vivamus fringilla lacus nec metus bibendum.
    
    \newfinding Iaculis massa nisl malesuada lacinia integer.
    \label{finding: template-finding}
\end{tcolorbox}

% Pointer sentence to Discussion:
% Single sentence forwarding the reader to the Discussion for qualitative
% interpretation and implications.
Lorem ipsum dolor sit amet consectetur adipiscing elit.
Quisque faucibus ex sapien vitae pellentesque sem placerat. 
In id cursus mi pretium tellus duis convallis. 
Tempus leo eu aenean sed diam urna tempor. 

% =========================================================================
% SECTION: DISCUSSION
% Narrative Flow: Narrow (Interpretations) -> Broad (Implications) 
%                -> Defensive (Limitations) -> Forward-Looking (Future Work)
% =========================================================================
\section{Discussion}
\label{sec:discussion}

% --- PARAGRAPH 1: DISCUSSION OPENING ---
% Target: 7–9 sentences. 
% Content: Restate the high‑level takeaway of the study, tying back to the Introduction’s big gap and summarizing the overall pattern across RQ1–RQ3. This paragraph sets the tone for the rest of the Discussion.
Lorem ipsum dolor sit amet consectetur adipiscing elit. 
Quisque faucibus ex sapien vitae pellentesque sem placerat. 
In id cursus mi pretium tellus duis convallis. 
Tempus leo eu aenean sed diam urna tempor. 
Pulvinar vivamus fringilla lacus nec metus bibendum egestas. 
Iaculis massa nisl malesuada lacinia integer nunc posuere. 
Ut hendrerit semper vel class aptent taciti sociosqu. 

% --- SUBSECTION: INTERPRETATION OF RQS ---
% Per‑RQ mini‑discussions: each \myparagraph is 8–12 sentences total, structured as interpretation → mechanism → relation to prior work → scoped implication/limit.
\subsection{Interpretation of Research Questions}
\label{subsec:interpretations}

\myparagraph{RQ$_1$ -- Baseline Prevalence Metrics}
% Interpretation: 2–3 sentences summarizing how the findings answer RQ1.
Lorem ipsum dolor sit amet consectetur adipiscing elit. Quisque faucibus ex sapien vitae pellentesque sem placerat. In id cursus mi pretium tellus duis convallis.
% Mechanism/Why: 2–3 sentences explaining why these patterns likely occur, grounded in theory or domain knowledge.
Lorem ipsum dolor sit amet consectetur adipiscing elit. Quisque faucibus ex sapien vitae pellentesque sem placerat. In id cursus mi pretium tellus duis convallis.
% Relation to prior work: 2–3 sentences comparing RQ1 findings to specific prior studies; note agreements and contradictions.
Lorem ipsum dolor sit amet consectetur adipiscing elit. Quisque faucibus ex sapien vitae pellentesque sem placerat. In id cursus mi pretium tellus duis convallis.
% Scoped actionable/limit: 2–3 sentences stating implications or recommendations for practice, plus one key limitation that tempers those recommendations.
Lorem ipsum dolor sit amet consectetur adipiscing elit. Quisque faucibus ex sapien vitae pellentesque sem placerat. In id cursus mi pretium tellus duis convallis. Tempus leo eu aenean sed diam urna tempor.

\myparagraph{RQ$_2$ -- The Mechanics of Fragile Reuse}
% Interpretation: 2–3 sentences summarizing how the findings answer RQ2.
Lorem ipsum dolor sit amet consectetur adipiscing elit.
% Mechanism/Why: 2–3 sentences explaining why these patterns likely occur, grounded in theory or domain knowledge.
Lorem ipsum dolor sit amet consectetur adipiscing elit. Quisque faucibus ex sapien vitae pellentesque sem placerat. In id cursus mi pretium tellus duis convallis.
% Relation to prior work: 2–3 sentences comparing RQ2 findings to specific prior studies.
Lorem ipsum dolor sit amet consectetur adipiscing elit. Quisque faucibus ex sapien vitae pellentesque sem placerat. In id cursus mi pretium tellus duis convallis.
% Scoped actionable/limit: 2–3 sentences stating implications or recommendations for practice.
Lorem ipsum dolor sit amet consectetur adipiscing elit. Quisque faucibus ex sapien vitae pellentesque sem placerat. In id cursus mi pretium tellus duis convallis. Tempus leo eu aenean sed diam urna tempor.

\myparagraph{RQ$_3$ -- Quantification of Pipeline Half-Life}
% Interpretation: 2–3 sentences summarizing how the findings answer RQ3.
Lorem ipsum dolor sit amet consectetur adipiscing elit.
% Mechanism/Why: 2–3 sentences explaining why these patterns likely occur, grounded in theory or domain knowledge.
Lorem ipsum dolor sit amet consectetur adipiscing elit. Quisque faucibus ex sapien vitae pellentesque sem placerat. In id cursus mi pretium tellus duis convallis.
% Relation to prior work: 2–3 sentences comparing RQ3 findings to specific prior studies.
Lorem ipsum dolor sit amet consectetur adipiscing elit. Quisque faucibus ex sapien vitae pellentesque sem placerat. In id cursus mi pretium tellus duis convallis.
% Scoped actionable/limit: 2–3 sentences stating implications or recommendations for practice.
Lorem ipsum dolor sit amet consectetur adipiscing elit. Quisque faucibus ex sapien vitae pellentesque sem placerat. In id cursus mi pretium tellus duis convallis. Tempus leo eu aenean sed diam urna tempor.

% --- SUBSECTION: WIDER IMPLICATIONS ---
\subsection{Wider Implications for Science and Engineering}
\label{subsec:implications}

% Paragraph 1 – Theoretical implications: 
% Target: 5–7 sentences. 
% Content: Synthesize across RQs to state how the study refines or challenges existing theories/models. Focus on general principles.
Lorem ipsum dolor sit amet consectetur adipiscing elit. 
Quisque faucibus ex sapien vitae pellentesque sem placerat. 
In id cursus mi pretium tellus duis convallis. 
Tempus leo eu aenean sed diam urna tempor. 
Pulvinar vivamus fringilla lacus nec metus bibendum egestas. 
Iaculis massa nisl malesuada lacinia integer nunc posuere. 
Ut hendrerit semper vel class aptent taciti sociosqu.

% Paragraph 2 – Practical/industrial implications: 
% Target: 5–7 sentences. 
% Content: Translate findings into concrete recommendations for teams, tools, and processes, mentioning feasibility and constraints.
Lorem ipsum dolor sit amet consectetur adipiscing elit. 
Quisque faucibus ex sapien vitae pellentesque sem placerat. 
In id cursus mi pretium tellus duis convallis. 
Tempus leo eu aenean sed diam urna tempor. 
Pulvinar vivamus fringilla lacus nec metus bibendum egestas. 
Iaculis massa nisl malesuada lacinia integer nunc posuere. 
Ut hendrerit semper vel class aptent taciti sociosqu.

% Optional 3rd Paragraph – Policy / Educational Frameworks:
% NOTE: Optional, beyond the two standard implication types (Theoretical and
%   Practical/Industrial) defined by the style guide. Delete if not needed.
% Target: 5–7 sentences on implications for standards, curricula, or broader scientific practice.
Lorem ipsum dolor sit amet consectetur adipiscing elit. 
Quisque faucibus ex sapien vitae pellentesque sem placerat. 
In id cursus mi pretium tellus duis convallis. 
Tempus leo eu aenean sed diam urna tempor. 
Pulvinar vivamus fringilla lacus nec metus bibendum egestas. 

% --- SUBSECTION: LIMITATIONS AND BIASES ---
\subsection{Limitations and Alternative Explanations}
\label{subsec:limitations}

% Paragraph 1 – Main limitations: 
% Target: 5–7 sentences. 
% Content: Summarize key limitations across design, data, and measures; explain how each might bias results or limit generality.
Lorem ipsum dolor sit amet consectetur adipiscing elit. 
Quisque faucibus ex sapien vitae pellentesque sem placerat. 
In id cursus mi pretium tellus duis convallis. 
Tempus leo eu aenean sed diam urna tempor. 
Pulvinar vivamus fringilla lacus nec metus bibendum egestas. 

% Paragraph 2 – Alternative explanations: 
% Target: 5–7 sentences. 
% Content: Present plausible alternative interpretations of major findings and briefly note what further data/methods would be needed to test them.
Lorem ipsum dolor sit amet consectetur adipiscing elit. 
Quisque faucibus ex sapien vitae pellentesque sem placerat. 
In id cursus mi pretium tellus duis convallis. 
Tempus leo eu aenean sed diam urna tempor. 
Pulvinar vivamus fringilla lacus nec metus bibendum egestas. 

% --- SUBSECTION: FUTURE DIRECTIONS ---
\subsection{Future Research Trajectories}
\label{subsec:future_directions}

% Paragraph 1 – Next logical steps: 
% Target: 4–6 sentences. 
% Content: Propose immediate follow‑up studies (replications, scaling, boundary conditions), each explicitly linked to a specific result or limitation.
Lorem ipsum dolor sit amet consectetur adipiscing elit. 
Quisque faucibus ex sapien vitae pellentesque sem placerat. 
In id cursus mi pretium tellus duis convallis. 
Tempus leo eu aenean sed diam urna tempor. 

% Paragraph 2 – Methodological refinements: 
% Target: 3–5 sentences. 
% Content: Suggest improvements to design, metrics, or data collection that would strengthen future evidence.
Lorem ipsum dolor sit amet consectetur adipiscing elit. 
Quisque faucibus ex sapien vitae pellentesque sem placerat. 
In id cursus mi pretium tellus duis convallis. 

% Paragraph 3 – New questions: 
% Target: 4–6 sentences. 
% Content: Outline new research questions opened by unexpected findings or implications; explain why each is important for empirical SE.
Lorem ipsum dolor sit amet consectetur adipiscing elit. 
Quisque faucibus ex sapien vitae pellentesque sem placerat. 
In id cursus mi pretium tellus duis convallis. 
Tempus leo eu aenean sed diam urna tempor. 

% Final closing block – Take-Home Message:
% Target: 4–6 sentences. 
% Content: Provide a synthesized, confident statement of what the study ultimately shows and why it matters. No new ideas or data; this should echo the core contribution and high‑level impact.
\myparagraph{Take-Home Message}
Lorem ipsum dolor sit amet consectetur adipiscing elit. 
Quisque faucibus ex sapien vitae pellentesque sem placerat. 
In id cursus mi pretium tellus duis convallis. 
Tempus leo eu aenean sed diam urna tempor. 
Pulvinar vivamus fringilla lacus nec metus bibendum egestas. 
Iaculis massa nisl malesuada lacinia integer nunc posuere. 

% =========================================================================
% SECTION: THREATS TO VALIDITY
% Function: Systematically classify study vulnerabilities into three subsections
%   (Internal, External, Construct validity). Be candid; each block names a threat,
%   its potential effect on results, and the mitigation taken.
% =========================================================================
\section{Threats to Validity}
\label{sec:threats}

% Section opening:
% Target: 3–4 sentences. 
% Content: Explain the purpose of this section (systematically analyzing internal, external, and construct validity threats) and note that detailed RQ‑specific threats were mentioned in each RQ methodology subsection.
Lorem ipsum dolor sit amet consectetur adipiscing elit. 
Quisque faucibus ex sapien vitae pellentesque sem placerat. 
In id cursus mi pretium tellus duis convallis. 
Our analysis follows the standard empirical software engineering classification distinguishing between internal, external, and construct validity.

% --- SUBSECTION: INTERNAL VALIDITY ---
\subsection{Internal Validity}
\label{subsec:internal_validity}

% First block (Confounding Variables):
% Target: 3–5 sentences.
% Content: Explain potential confounding variables and how they may distort causal
%   interpretations; name the mitigation strategies used. Use causal signposts
%   ("As a result", "Consequently").
Lorem ipsum dolor sit amet consectetur adipiscing elit. 
Quisque faucibus ex sapien vitae pellentesque sem placerat. 
In id cursus mi pretium tellus duis convallis. 
Tempus leo eu aenean sed diam urna tempor. 
Pulvinar vivamus fringilla lacus nec metus bibendum egestas. 

% Second block (History, Maturation, and Instrumentation):
% Target: 3–4 sentences.
% Content: Discuss history, maturation, and instrumentation effects and their
%   potential influence on outcomes; note how design choices attempt to control them.
Lorem ipsum dolor sit amet consectetur adipiscing elit. 
Quisque faucibus ex sapien vitae pellentesque sem placerat. 
In id cursus mi pretium tellus duis convallis. 
Tempus leo eu aenean sed diam urna tempor. 

% Third block (Attrition and Missing Data):
% Target: 2–3 sentences.
% Content: Describe attrition/missing data patterns and how they were handled
%   (e.g., exclusion, imputation, sensitivity analysis).
Lorem ipsum dolor sit amet consectetur adipiscing elit. 
Quisque faucibus ex sapien vitae pellentesque sem placerat. 
In id cursus mi pretium tellus duis convallis. 

% --- SUBSECTION: EXTERNAL VALIDITY ---
\subsection{External Validity}
\label{subsec:external_validity}

% First block (Context & Generalizability):
% Target: 3–4 sentences.
% Content: Clarify the study setting and domains, and explicitly state the boundary
%   conditions — where the results do and do not generalize.
Lorem ipsum dolor sit amet consectetur adipiscing elit. 
Quisque faucibus ex sapien vitae pellentesque sem placerat. 
In id cursus mi pretium tellus duis convallis. 
Tempus leo eu aenean sed diam urna tempor. 

% Second block (Sampling Bias):
% Target: 2–3 sentences.
% Content: Discuss how selection of projects/participants may skew results and how
%   that limits external validity.
Lorem ipsum dolor sit amet consectetur adipiscing elit. 
Quisque faucibus ex sapien vitae pellentesque sem placerat. 
In id cursus mi pretium tellus duis convallis. 

% --- SUBSECTION: CONSTRUCT VALIDITY ---
\subsection{Construct Validity}
\label{subsec:construct_validity}

% First block (Indirect Metrics and Operationalization):
% Target: 4–6 sentences.
% Content: Justify that the chosen metrics capture the intended constructs,
%   referencing prior work and triangulation where possible.
Lorem ipsum dolor sit amet consectetur adipiscing elit. 
Quisque faucibus ex sapien vitae pellentesque sem placerat. 
In id cursus mi pretium tellus duis convallis. 
Tempus leo eu aenean sed diam urna tempor. 
Pulvinar vivamus fringilla lacus nec metus bibendum egestas. 
Iaculis massa nisl malesuada lacinia integer nunc posuere. 
Ut hendrerit semper vel class aptent taciti sociosqu. 
Ad litora torquent per conubia nostra inceptos himenaeos.

% Second block (Measurement Error and Technical Bias):
% Target: 3–5 sentences.
% Content: Address logging inaccuracies, self-report biases, and tool limitations;
%   explain how they might affect interpretations.
Lorem ipsum dolor sit amet consectetur adipiscing elit. 
Quisque faucibus ex sapien vitae pellentesque sem placerat. 
In id cursus mi pretium tellus duis convallis. 
Tempus leo eu aenean sed diam urna tempor. 
Pulvinar vivamus fringilla lacus nec metus bibendum egestas. 
Iaculis massa nisl malesuada lacinia integer nunc posuere. 

% =========================================================================
% SECTION: CONCLUSION
% =========================================================================
\section{Conclusion}
\label{sec:conclusion}

% Paragraph 1 – Aims and findings synthesis: 
% Target: 2–4 sentences. 
% Content: Restate the overall problem and aims in fresh language, then synthesize the main findings across RQs (no detailed numbers).
Lorem ipsum dolor sit amet consectetur adipiscing elit. 
Quisque faucibus ex sapien vitae pellentesque sem placerat. 
In id cursus mi pretium tellus duis convallis. 
Tempus leo eu aenean sed diam urna tempor. 
Pulvinar vivamus fringilla lacus nec metus bibendum egestas. 
Iaculis massa nisl malesuada lacinia integer nunc posuere. 
Ut hendrerit semper vel class aptent taciti sociosqu.

% Paragraph 2 – Contributions and practical implications: 
% Target: 2–4 sentences. 
% Content: Explicitly name the key contributions and what they enable or change for researchers and practitioners.
Lorem ipsum dolor sit amet consectetur adipiscing elit. 
Quisque faucibus ex sapien vitae pellentesque sem placerat. 
In id cursus mi pretium tellus duis convallis. 
Tempus leo eu aenean sed diam urna tempor. 
Pulvinar vivamus fringilla lacus nec metus bibendum egestas. 
Iaculis massa nisl malesuada lacinia integer nunc posuere. 
Ut hendrerit semper vel class aptent taciti sociosqu.

% Paragraph 3 – Limitations, future work, final takeaway: 
% Target: 3–4 sentences. 
% Style Rule: The final sentence must be a confident, synthesized statement—not an apology.
% Content: Echo the most important limitation/future direction at a high level and end with a strong single sentence that captures the study’s core message.
Lorem ipsum dolor sit amet consectetur adipiscing elit. 
Quisque faucibus ex sapien vitae pellentesque sem placerat. 
In id cursus mi pretium tellus duis convallis. 
Tempus leo eu aenean sed diam urna tempor. 
Pulvinar vivamus fringilla lacus nec metus bibendum egestas. 
Iaculis massa nisl malesuada lacinia integer nunc posuere. 
Ut hendrerit semper vel class aptent taciti sociosqu. 

% =========================================================================
% ACKNOWLEDGMENTS
% =========================================================================
\section*{Acknowledgments}
% Target: 1–3 sentences. 
% Content: Thank funding sources, collaborators, and infrastructure. No technical content.
Lorem ipsum dolor sit amet consectetur adipiscing elit. 
Quisque faucibus ex sapien vitae pellentesque sem placerat. 
In id cursus mi pretium tellus duis convallis. 

\bibliographystyle{ACM-Reference-Format}
\bibliography{citations}    % For all other submissions
%% =========================================================================
% DOCUMENT CLASS & SETUP
% =========================================================================

% \documentclass[sigconf]{acmart}       % For final camera-ready peer-reviewed submissions
\documentclass[sigconf,nonacm]{acmart} % For arXiv/preprint: removes copyright blocks & ACM footers

% =========================================================================
% GLOBAL PARAGRAPH AND PROSE RULES
% =========================================================================
% 1. Target Length: Aim for 3–7 sentences per paragraph. Avoid single-sentence 
%    paragraphs and dense walls of text.
% 2. Functional Focus: Each paragraph must serve a singular rhetorical function 
%    (e.g., context, problem, method, finding, interpretation, implication, 
%    limitation, or future work).
% 3. Explicit Signposting: Maintain visible logical flow across paragraph 
%    boundaries using structural transitions (e.g., "However", "As a result", 
%    "Consequently", "In contrast").
% 4. Micro-Sections: Use \myparagraph{...} blocks to signify named mini-paragraphs 
%    with specialized micro-rhetorical goals (e.g., "Contributions", "Significance").

% =========================================================================
% CONDITIONAL DEBUGGING FLAGS
% =========================================================================

\newif\ifDEBUG  % Declares a boolean switch called "DEBUG" for conditional compilation
\DEBUGtrue      % Set DEBUG to true (activates draft/todo notes)
\DEBUGfalse     % Set DEBUG to false (silences draft notes for clean compilation)
% NET EFFECT: \DEBUGfalse runs last, so debug is OFF by default. While OFF,
% \TODO{}, \CITE{}, and \NMS{} (co-author notes) render as nothing.
% To see draft markers, comment out the \DEBUGfalse line above.

% =========================================================================
% MODULAR PREAMBLE IMPORTS
% =========================================================================

% =========================================================================
% misc/data.tex
% Global Data Macro Definitions
% =========================================================================
        % Imports global stats, variables, and data macros
% =========================================================================
% core typography & layout engine
% =========================================================================

% \usepackage{amssymb,amsmath,amsfonts} % Optional math packages (currently commented)
\usepackage[T1]{fontenc}               % Uses modern 8-bit font encoding for crisp glyph rendering
\usepackage{balance}                   % Equalizes the length of columns on the final page of a 2-column layout

% =========================================================================
% code listing & custom graphical callout containers
% =========================================================================

\usepackage[most]{tcolorbox}           % Advanced custom framed boxes with advanced skins/listings hooks
\usepackage{listings}                  % Native source code block layout engine
\usepackage{xcolor}                    % Comprehensive driver-independent color naming & mixing support

% =========================================================================
% standard acm-approved helper packages
% =========================================================================

% --- structural & page geometry control
\usepackage{afterpage}                 % Postpones execution of commands until the next page break
\usepackage{float}                     % Enhances figure place control; introduces the rigid [H] specifier
\usepackage{lscape}                    % Places specific content pages in landscape orientation
\usepackage{rotating}                  % Rotates text, tables, and figures along arbitrary angle limits

% --- mathematics & computational logic
\usepackage{algorithmic}               % Structural pseudocode layout environment builder

% --- advanced tabular structures
\usepackage{array}                     % Extended configuration controls for table column definitions
\usepackage{booktabs}                  % Renders publication-grade professional horizontal table rules
\usepackage{makecell}                  % Multiline table cell formatting with manual line breaks
\usepackage{multirow}                  % Enables cells spanning multiple horizontal rows
\usepackage{tabularx}                  % Auto-calculates column widths to fit specified horizontal constraints
\usepackage{longtable}                 % Handles multi-page tables natively within the standard stream

% --- captions & bibliographic citations
\usepackage[skip=4pt]{caption}         % Uniform control over caption spacing across figures and tables
\usepackage[compress]{cite}            % Automatically sorts and collapses numeric sequence reference tags
\usepackage{natbib}[numbers]           % Explicit configuration for structured numeric reference citation types

% --- document hyperlinking & reference automation
\usepackage{hyperref}                  % Injects interactive hyperlinks and global document indexes
\usepackage{cleveref}                  % Smart inline references (adds contextual prefixes like "Section" or "Fig.")

% --- list layout control
\usepackage{enumitem}                  % Fine-grained layout control over itemize, enumerate, and description lists

% --- graphic assets, icons, & utility packages
% \usepackage{filecontents}            % Pass-through payload file management (currently commented)
\usepackage{fontawesome5}              % Vector icon engine for system design and repository markers
\usepackage{graphicx}                  % Core graphics scaling, cropping, and rendering extension
\usepackage{lipsum}                    % Generates dummy text placeholders for layout prototyping
\usepackage{siunitx}                   % Uniform typesetting formatting engine for technical units and data measurements
\usepackage{soul}                      % Provides tracking, kerning, and strike-through/text-highlighting features
\usepackage{textcomp}                  % Provides specialized miscellaneous symbols and text glyphs
\usepackage[normalem]{ulem}            % Smooth underlining control; [normalem] prevents breaking \emph highlights
\usepackage{xspace}                    % Intelligent trailing space macro insertion engine

% =========================================================================
% page design & heading spacing tweaks
% =========================================================================

% \usepackage[expectsections, unnumberedsections, fencedCode]{markdown} % Direct Markdown parsing engine (commented)
% \usepackage[vskip=1em,font=itshape,leftmargin=2em,rightmargin=2em]{quoting} % Advanced blockquote layout engine (commented)

% Compacted section structural layout overrides (currently commented)
% \titlespacing*{\section}{0pt}{6pt}{4pt}
% \titlespacing*{\subsection}{0pt}{5pt}{3pt}
% \titlespacing*{\subsubsection}{0pt}{4pt}{2pt}

% =========================================================================
% CUSTOM UTILITY MACROS
% =========================================================================

% --- standard abbreviated latin syntax styling
\newcommand{\ie}{\textit{i.e.},\xspace}  % Uniformly formats "id est" with correct italicization and spacing
\newcommand{\eg}{\textit{e.g.},\xspace}  % Uniformly formats "exempli gratia" with correct italicization and spacing

% --- layout constructs
\newcommand{\myparagraph}[1]{\vspace{0.20cm}\noindent\textbf{#1} \noindent{}} % Spaced bold headers for inline subsections

% --- generative framework / prompts callout containers
% =========================================================================
% LLM PROMPT MACRO TEMPLATE
% Purpose: Formats Large Language Model system instructions, context windows,
%          or user prompts within a distinct, clean visual container.
%
% Usage: \myprompt{Argument #1: Box Heading}{Argument #2: Prompt Text Payload}
% Example: \myprompt{System Prompt}{Identify software reuse patterns...}
% =========================================================================

\newcommand{\myprompt}[2]{%
  \begin{tcolorbox}[
    enhanced,              % Activates advanced skin engine hooks (requires tcolorbox library 'skins')
    colback=gray!10,       % Light gray canvas fill background (10% opacity black over white canvas)
    colframe=black!20,     % Muted gray border line frame color (20% opacity black)
    boxrule=0.4pt,         % Ultra-thin, crisp frame line boundary thickness
    arc=1mm,               % Sets a subtle, modern 1mm corner radius curve
    % Tight geometric padding rules to prevent large text blocks from forcing column overflows:
    left=2mm,              % Internal text padding buffer from the left border
    right=2mm,             % Internal text padding buffer from the right border
    top=1mm,               % Internal text padding buffer beneath the title separator rule
    bottom=1mm,            % Internal text padding buffer above the lower frame floor
    fonttitle=\bfseries,   % Sets the box title header typography font to boldface
    coltitle=black,        % Keeps the title string text plain black for contrast over gray background
    title={#1},            % Evaluates Positional Argument #1 as the structural header text
    % Elastic vertical paragraph spacing rules that stretch or shrink to prevent layout underfull breaks:
    before skip=3pt plus 2pt minus 1pt,
    after skip=3pt plus 2pt minus 1pt
  ]
  \small #2                % Reduces text size inside the box to \small for better scannability of dense text payloads
  \end{tcolorbox}%         % The trailing '%' character strips accidental trailing whitespaces from compilation output
}

% --- structured empirical finding blocks
\newcounter{finding}                   % Declares custom structural metric serialization register
\newcommand{\newfinding}{%             % Instantiates inline empirical milestone callouts
    \refstepcounter{finding}           % Increments register metrics sequentially and binds anchor links for \label
    \textbf{Finding \thefinding.}   % Prints stylized numeric marker index into current structural block stream
}

% =========================================================================
% REVIEWER INTERACTION & COLLABORATIVE DEBUG SWITCHES
% =========================================================================

\ifDEBUG
    % --- active draft editing marker functions
    \newcommand{\TODO}[1]{\hl{TODO: #1}} % Visual yellow structural highlight marker
    \newcommand{\CITE}[1]{\hl{CITE: #1}} % Visual yellow missing source tracking marker
    
    % --- colored co-author editorial discussion nodes
    \newcommand{\NMS}[1]{\textcolor{orange}{[NMS: #1]}} % Nicholas M. Synovic
\else
    % --- production compilation overrides (silences markers entirely)
    \newcommand{\TODO}[1]{}
    \newcommand{\CITE}[1]{}
    
    \newcommand{\NMS}[1]{}
\fi % Imports formatting overrides, extra packages, and layout tweaks

% =========================================================================
% CUSTOM LOGO MACROS
% =========================================================================

% Delays definition until document start to prevent package conflicts
\AtBeginDocument{%
    % Safely defines \BibTeX logo if not already defined by a loaded package
    \providecommand\BibTeX{{%
        \normalfont B\kern-0.5em{%   % Backspaces slightly after 'B'
            \scshape                 % Switches to small caps for 'ib'
            i\kern-0.25em b}\kern-0.8em\TeX}}} 
% Micro-spaces 'ib' and pulls 'TeX' close

% =========================================================================
% RIGHTS & COPYRIGHT MANAGEMENT (Suppressed while 'nonacm' is active)
% =========================================================================

\setcopyright{acmlicensed}          % Specifies the ACM licensing agreement type
\copyrightyear{2026}                % Sets the explicit copyright year
\acmYear{2026}                      % Sets the publishing year index
\acmDOI{XXXXXXX.XXXXXXX}            % Placeholder for the Digital Object Identifier string
\acmISBN{978-1-4503-XXXX-X/18/06}   % Placeholder for the standard publication ISBN

% Defines conference acronym, full metadata description, dates, and geographic location
\acmConference[Conference acronym 'XX]{Make sure to enter the correct
  conference title from your rights confirmation email}{June 03--05,
  2018}{Woodstock, NY}

% =========================================================================
% DOCUMENT CONTENT START
% =========================================================================

\begin{document}

\title{\textbf{INSERT TITLE}} % Formats the core paper title in bold face

% -------------------------------------------------------------------------
% AUTHOR METADATA
% -------------------------------------------------------------------------

\author{Nicholas M. Synovic}
\email{nsynovic@luc.edu}
\orcid{0000-0003-0413-4594} % Unique digital author identifier
\affiliation{
  \institution{Loyola University Chicago}
  \streetaddress{1032 W Sheridan Rd}
  \city{Chicago}
  \state{Illinois}
  \country{USA}
  \postcode{60660}
}

% Replaces the full author list in page running headers with a concise version to save space
\renewcommand{\shortauthors}{Synovic}

% -------------------------------------------------------------------------
% ABSTRACT SECTION
% -------------------------------------------------------------------------

% ABSTRACT: one self-contained, three-part narrative block (6–8 sentences total).
% HARD RULES: no citations and no raw statistical tables/numbers in the abstract.
\begin{abstract}
% Part 1 — Context & Problem (Sentences 1–3):
% Establish the broad SE domain and why it matters, then pivot to the core
% friction point or consequence. Function: context → problem.
Lorem ipsum dolor sit amet consectetur adipiscing elit. 
Quisque faucibus ex sapien vitae pellentesque sem placerat. 
In id cursus mi pretium tellus duis convallis. 

% Part 2 — Approach & Scope (Sentences 4–5):
% State "In this paper, we…" and summarize the design/method at one-sentence
% granularity, including evaluation scope (data size, system types, participants).
Tempus leo eu aenean sed diam urna tempor. 
Pulvinar vivamus fringilla lacus nec metus bibendum egestas. 

% Part 3 — Impact & Insights (Sentences 6–8):
% Deliver the headline qualitative/quantitative outcome, the supporting insight,
% and one concrete implication for practice or future research.
% Keep it high-level: no detailed numbers or statistics.
Iaculis massa nisl malesuada lacinia integer nunc posuere. 
Ut hendrerit semper vel class aptent taciti sociosqu. 
Ad litora torquent per conubia nostra inceptos himenaeos.
\end{abstract}

% =========================================================================
% ACM COMPUTING CLASSIFICATION SYSTEM (https://dl.acm.org/ccs)
% =========================================================================
\begin{CCSXML}
\end{CCSXML}

% The \ccsdesc macros render XML classifications into the document footer
\ccsdesc[500]{Software and its engineering}

% =========================================================================
% KEYWORDS INDEXING
% =========================================================================

\keywords{Replace me}

\maketitle

% =========================================================================
% SECTION: INTRODUCTION
% Function: Sell the paper. Six-paragraph arc; each block has ONE rhetorical job.
% Arc (one block each): broad context/hook → narrowed gap → goal & core insight
%   → contributions (bulleted) → significance/scope → roadmap.
% Rules: 3–7 sentences per paragraph; explicit transitions between blocks
%   ("However", "As a result"); no unquantified hype ("novel", "significant").
% =========================================================================
% RULE: every figure must be cited in-text via \Cref{fig:...} (e.g. \Cref{fig:template_figure}).
% =========================================================================
% TEMPLATE FIGURE PLACEHOLDER
% Uses a standard boxed paragraph container to stand in for missing visual assets.
% =========================================================================

\begin{figure}[h] % [h] requests placement "here" inline with the surrounding text stream
    \centering    % Horizontally centers the figure artifact within the column/page margins
    
    % \fbox draws a visible frame border around its inner content block
    \fbox{
        % \parbox structure creates an explicit, fixed-size paragraph boundary container
        % Arguments format: [inner-alignment][height][content-alignment]{width}
        \parbox[c][2in][c]{3in}{
            \centering Missing Figure % Centered placeholder text string inside the bounding box
        }
    }
    
    % Core descriptive label rendered immediately beneath the figure box frame
    \caption{Template figure placeholder}
    
    % Unique cross-referencing anchor tag; reference this asset elsewhere via \Cref{fig:template_figure}
    \label{fig:template_figure}
\end{figure}
\section{Introduction}
\label{sec:introduction}

% --- PARAGRAPH 1: CONTEXT AND HOOK ---
% Function: Intro paragraph 1.
% Target: 3–6 sentences.
% Content: Establish the broad software engineering domain, why it matters (reliability, cost, UX, scientific impact), and a hook that shows the domain is non‑trivial. No detailed technicalities here.
Lorem ipsum dolor sit amet consectetur adipiscing elit.
Quisque faucibus ex sapien vitae pellentesque sem placerat. 
In id cursus mi pretium tellus duis convallis.
Tempus leo eu aenean sed diam urna tempor. 
Pulvinar vivamus fringilla lacus nec metus bibendum egestas.
Iaculis massa nisl malesuada lacinia integer nunc posuere. 

% --- PARAGRAPH 2: NARROWED PROBLEM AND GAP ---
% Function: Intro paragraph 2.
% Target: 4–7 sentences.
% Content: Summarize current practice/prior work, then pivot with a clear ‘However/Yet/Despite this’ sentence that states the precise gap or friction point. End by emphasizing why this gap is consequential.
Lorem ipsum dolor sit amet consectetur adipiscing elit. 
Quisque faucibus ex sapien vitae pellentesque sem placerat.
In id cursus mi pretium tellus duis convallis. 
Tempus leo eu aenean sed diam urna tempor.
Pulvinar vivamus fringilla lacus nec metus bibendum egestas. 
Iaculis massa nisl malesuada lacinia integer nunc posuere.
Ut hendrerit semper vel class aptent taciti sociosqu. 

% --- PARAGRAPH 3: GOAL AND APPROACH (THE CORE INSIGHT) ---
% Function: Intro paragraph 3.
% Target: 3–6 sentences.
% Content: Present the overall research goal and high‑level approach (system, study, framework). Include plain‑language intuition about why your approach should address the stated gap.
Lorem ipsum dolor sit amet consectetur adipiscing elit. 
Quisque faucibus ex sapien vitae pellentesque sem placerat.
In id cursus mi pretium tellus duis convallis. 
Tempus leo eu aenean sed diam urna tempor.
Pulvinar vivamus fringilla lacus nec metus bibendum egestas. 
Iaculis massa nisl malesuada lacinia integer nunc posuere.

% --- PARAGRAPH 4: CONTRIBUTIONS ---
% Structure: Use 3–4 bullets, each describing one type of contribution (artifact, empirical study, insight, open science). 
% Formatting Rule: Keep each item one sentence plus a short clarifying clause, avoiding vague phrases (‘novel’, ‘significant’) without specifics.
\myparagraph{Contributions}
This paper makes the following core contributions:
\begin{itemize}[leftmargin=16pt, rightmargin=5pt, topsep=0pt]
    % Contribution 1: The Artifact (Tool / Framework / Dataset)
    \item Lorem ipsum dolor sit amet consectetur adipiscing elit.
    % Contribution 2: The Empirical Scale & Matrix Analysis
    \item Lorem ipsum dolor sit amet consectetur adipiscing elit.
    % Contribution 3: The Findings, Open Science & Replication Package
    \item Lorem ipsum dolor sit amet consectetur adipiscing elit.
\end{itemize}

% --- PARAGRAPH 5: SIGNIFICANCE & SCOPE ---
% Function: Intro paragraph 5 — "headline results + scope" preview (required).
% Target: 3–5 sentences.
% Content: Provide 1–2 headline quantitative or qualitative outcomes, plus a brief statement of scope/assumptions (where it works, where it might not). Do not repeat details from the Results; this is a teaser, not the findings.
\myparagraph{Significance}
Lorem ipsum dolor sit amet consectetur adipiscing elit. 
Quisque faucibus ex sapien vitae pellentesque sem placerat.
In id cursus mi pretium tellus duis convallis. 
Tempus leo eu aenean sed diam urna tempor.
Pulvinar vivamus fringilla lacus nec metus bibendum egestas. 

% --- PARAGRAPH 6: PAPER ROADMAP ---
% Function: Manuscript outline (Intro paragraph 6).
% Target: 2–4 sentences.
% Content: Walk the reader through the paper in document order so the roadmap
%   doubles as a structural map. The actual order is:
%   Background (sec:background) → Research Questions (sec:research_questions)
%   → RQ1 (sec:rq1_evaluation) → RQ2 (sec:rq2_evaluation)
%   → RQ3 (sec:rq3_evaluation) → Discussion (sec:discussion)
%   → Threats (sec:threats) → Conclusion (sec:conclusion).
%   NOTE: there is no standalone "Methodology" section — methodology is embedded
%   inside each RQ section. Use lowercase section names + \Cref macros; never
%   hardcode numbers.
\myparagraph{Paper Outline}
This paper is organized as follows: 
\Cref{sec:background} provides essential background and reviews related work, and 
\Cref{sec:research_questions} formalizes the research questions.
\Cref{sec:rq1_evaluation}, \Cref{sec:rq2_evaluation}, and \Cref{sec:rq3_evaluation} each present the methodology and empirical findings for one research question.
The paper concludes in \Cref{sec:discussion} with a discussion of implications for the software engineering community and future work, followed by an examination of potential threats to validity in \Cref{sec:threats} and closing remarks in \Cref{sec:conclusion}.

% =========================================================================
% SECTION: BACKGROUND AND RELATED WORK
% Narrative Arc: "What world are we in?" -> "What fails here?" -> "How do we fix it?"
% =========================================================================
\section{Background and Related Work}
\label{sec:background}

% --- MAJOR SIGNPOSTING & RE-CONTEXTUALIZATION ---
% Entry paragraph (3–5 sentences): re‑state the problem space with more technical depth than the Introduction and briefly signpost the three subsections (domain overview, existing approaches, conceptual foundation).
Lorem ipsum dolor sit amet consectetur adipiscing elit. 
Quisque faucibus ex sapien vitae pellentesque sem placerat.
In id cursus mi pretium tellus duis convallis. 
Tempus leo eu aenean sed diam urna tempor.
Pulvinar vivamus fringilla lacus nec metus bibendum egestas. 

% =========================================================================
% SUBSECTION 1: THE DOMAIN OVERVIEW (Domain overview)
% =========================================================================
\subsection{Subsection 1}
\label{subsec:subsection-1}

% Paragraph 1 – Domain definitions and workflows: 
% Technical definitions and standard workflows. Clarify entities (projects, pipelines, registries), typical processes, and key terminology.
Lorem ipsum dolor sit amet consectetur adipiscing elit. 
Quisque faucibus ex sapien vitae pellentesque sem placerat.
In id cursus mi pretium tellus duis convallis. 
Tempus leo eu aenean sed diam urna tempor.
Pulvinar vivamus fringilla lacus nec metus bibendum egestas. 
Iaculis massa nisl malesuada lacinia integer nunc posuere.

% Paragraph 2 – Cause‑and‑effect linkages: 
% Explain how practices and constraints in this domain lead to specific challenges (e.g., fragility, reproducibility issues). Use ‘As a result/Consequently’ to make causal links explicit.
Lorem ipsum dolor sit amet consectetur adipiscing elit.
Quisque faucibus ex sapien vitae pellentesque sem placerat. 
In id cursus mi pretium tellus duis convallis.
Tempus leo eu aenean sed diam urna tempor. 
Pulvinar vivamus fringilla lacus nec metus bibendum egestas.
Iaculis massa nisl malesuada lacinia integer nunc posuere. 

% Paragraph 3 – Transition to existing approaches: 
% Close the subsection by hinting at what practitioners or researchers currently do to cope with these challenges, setting up Subsection 2.
Lorem ipsum dolor sit amet consectetur adipiscing elit. 
Quisque faucibus ex sapien vitae pellentesque sem placerat.
In id cursus mi pretium tellus duis convallis. 
Tempus leo eu aenean sed diam urna tempor.
Pulvinar vivamus fringilla lacus nec metus bibendum egestas. 
Iaculis massa nisl malesuada lacinia integer nunc posuere.

% =========================================================================
% SUBSECTION 2: EXISTING APPROACHES & PRIOR WORK
% =========================================================================
\subsection{Subsection 2}
\label{subsec:subsection-2}

% Paragraph 1 – Subsection entry and literature grouping: 
% Target: 3–5 sentences. 
% Content: Explicitly tie back to Subsection 1 and describe high‑level categories of existing approaches (e.g., tooling, guidelines, empirical studies). Group literature by theme, not chronologically.
Lorem ipsum dolor sit amet consectetur adipiscing elit. 
Quisque faucibus ex sapien vitae pellentesque sem placerat.
In id cursus mi pretium tellus duis convallis. 
Tempus leo eu aenean sed diam urna tempor.
Pulvinar vivamus fringilla lacus nec metus bibendum egestas. 
Iaculis massa nisl malesuada lacinia integer nunc posuere.

% Paragraph 2 – Contrastive analysis and limitations: 
% Content: Describe what these approaches achieve and where they fall short. Use contrast markers (‘However’, ‘In contrast’) and connect limitations to your gap.
Lorem ipsum dolor sit amet consectetur adipiscing elit.
Quisque faucibus ex sapien vitae pellentesque sem placerat. 
In id cursus mi pretium tellus duis convallis.
Tempus leo eu aenean sed diam urna tempor. 
Pulvinar vivamus fringilla lacus nec metus bibendum egestas.
Iaculis massa nisl malesuada lacinia integer nunc posuere. 

% Paragraph 3 – Exit transition to conceptual lens: 
% Content: Summarize the overall shortcomings of existing work and point forward to the conceptual foundation you will use to reason about them (Subsection 3).
Lorem ipsum dolor sit amet consectetur adipiscing elit. 
Quisque faucibus ex sapien vitae pellentesque sem placerat.
In id cursus mi pretium tellus duis convallis. 
Tempus leo eu aenean sed diam urna tempor.
Pulvinar vivamus fringilla lacus nec metus bibendum egestas. 
Iaculis massa nisl malesuada lacinia integer nunc posuere.

% =========================================================================
% SUBSECTION 3: CONCEPTUAL FOUNDATION
% =========================================================================
\subsection{Subsection 3}
\label{subsec:subsection-3}

% Paragraph 1 – Subsection entry and theoretical framing: 
% Content: Introduce the conceptual/theoretical lens (e.g., socio‑technical, diffusion, reliability) that underpins your study. Motivate why this lens is appropriate given the domain and prior work.
Lorem ipsum dolor sit amet consectetur adipiscing elit. 
Quisque faucibus ex sapien vitae pellentesque sem placerat.
In id cursus mi pretium tellus duis convallis. 
Tempus leo eu aenean sed diam urna tempor.
Pulvinar vivamus fringilla lacus nec metus bibendum egestas. 
Iaculis massa nisl malesuada lacinia integer nunc posuere.

% Paragraph 2 – Technical hook into methodology: 
% Content: Tie the chosen lens directly to your upcoming methodology, hinting at key constructs, relationships, or hypotheses that will be operationalized in the Methods and RQ sections.
Lorem ipsum dolor sit amet consectetur adipiscing elit. 
Quisque faucibus ex sapien vitae pellentesque sem placerat.
In id cursus mi pretium tellus duis convallis. 
Tempus leo eu aenean sed diam urna tempor.
Pulvinar vivamus fringilla lacus nec metus bibendum egestas. 
Iaculis massa nisl malesuada lacinia integer nunc posuere.

% =========================================================================
% SECTION: RESEARCH QUESTIONS (The Structural Bridge)
% =========================================================================
\section{Research Questions}
\label{sec:research_questions}

% --- INTRODUCTORY SIGNPOSTING & GAP ANCHORING ---
% Target: 1–3 sentences. 
% Content: Remind the reader of the key gap from Background and state that you now formalize the evaluation goals as research questions.
Lorem ipsum dolor sit amet consectetur adipiscing elit. 
Quisque faucibus ex sapien vitae pellentesque sem placerat.
In id cursus mi pretium tellus duis convallis. 

% --- THE FORMAL RQ LIST ---
% Function: The formal, enumerated research questions (the structural bridge).
% Content: List 3–4 concise RQs, each isolating ONE dimension (e.g., prevalence,
%   patterns, impact). Keep each RQ to a single sentence prefaced by a bolded
%   dimensional anchor (e.g., [Prevalence]).
% Format: per the style guide, the colon sits OUTSIDE the bold: \textbf{RQ\arabic*}:
\noindent We address the following research questions:
\begin{enumerate}[label=\textbf{RQ\arabic*}:, leftmargin=*]
    \item \textbf{[Dimension 1, e.g., Prevalence].} Lorem ipsum dolor sit amet consectetur adipiscing elit?
    \item \textbf{[Dimension 2, e.g., Patterns].} Lorem ipsum dolor sit amet consectetur adipiscing elit?
    \item \textbf{[Dimension 3, e.g., Impact].} Lorem ipsum dolor sit amet consectetur adipiscing elit?
\end{enumerate}

% --- MAPPING PARAGRAPH ---
% Target: 2–3 sentences. 
% Content: Map each RQ to its evaluation section (RQ1, RQ2, RQ3) and briefly mention which data and methods answer which question. This proves the downstream structure aligns with stated goals.
Lorem ipsum dolor sit amet consectetur adipiscing elit. 
Quisque faucibus ex sapien vitae pellentesque sem placerat.
In id cursus mi pretium tellus duis convallis. 

% =========================================================================
% SECTION: EVALUATION OF RESEARCH QUESTION 1 (RQ1)
% RQ1 block arc: overview & motivation → RQ1‑specific methodology → RQ1 findings → formal answer box → pointer to Discussion.
% =========================================================================
\section{Evaluation of RQ1: [Dimension Title, e.g., Prevalence]}
\label{sec:rq1_evaluation}

% --- RQ1 OVERVIEW PARAGRAPH ---
% Target: 6–8 sentences. 
% Content: Restate RQ1 verbatim, explain its role relative to other RQs (baseline, patterns, impact), and describe why this dimension matters. End with a sentence that introduces the upcoming methodology subsection.
Lorem ipsum dolor sit amet consectetur adipiscing elit. 
Quisque faucibus ex sapien vitae pellentesque sem placerat. 
In id cursus mi pretium tellus duis convallis.
Tempus leo eu aenean sed diam urna tempor. 
Pulvinar vivamus fringilla lacus nec metus bibendum egestas.
Iaculis massa nisl malesuada lacinia integer nunc posuere. 
Ut hendrerit semper vel class aptent taciti sociosqu. 
Ad litora torquent per conubia nostra inceptos himenaeos.

% --- SUBSECTION: METHODOLOGY FOR RQ1 ---
\subsection{RQ1 Methodology and Operationalization}
\label{subsec:rq1_methodology}
% RULE: cite every figure in-text via \Cref{fig:...}; never rely on placement alone.
% =========================================================================
% TEMPLATE FIGURE PLACEHOLDER
% Uses a standard boxed paragraph container to stand in for missing visual assets.
% =========================================================================

\begin{figure}[h] % [h] requests placement "here" inline with the surrounding text stream
    \centering    % Horizontally centers the figure artifact within the column/page margins
    
    % \fbox draws a visible frame border around its inner content block
    \fbox{
        % \parbox structure creates an explicit, fixed-size paragraph boundary container
        % Arguments format: [inner-alignment][height][content-alignment]{width}
        \parbox[c][2in][c]{3in}{
            \centering Missing Figure % Centered placeholder text string inside the bounding box
        }
    }
    
    % Core descriptive label rendered immediately beneath the figure box frame
    \caption{Template figure placeholder}
    
    % Unique cross-referencing anchor tag; reference this asset elsewhere via \Cref{fig:template_figure}
    \label{fig:template_figure}
\end{figure}

% Paragraph 1 – Design & rationale: 
% Explain the study design (experiment, case study, mining, survey) and why it is appropriate for RQ1.
Lorem ipsum dolor sit amet consectetur adipiscing elit. 
Quisque faucibus ex sapien vitae pellentesque sem placerat. 
In id cursus mi pretium tellus duis convallis. 
Tempus leo eu aenean sed diam urna tempor.

% Paragraph 2 – Data sources & sampling: 
% Describe datasets, projects, participants, and inclusion/exclusion rules. Provide enough detail for replication.
Lorem ipsum dolor sit amet consectetur adipiscing elit. 
Quisque faucibus ex sapien vitae pellentesque sem placerat. 
In id cursus mi pretium tellus duis convallis. 
Tempus leo eu aenean sed diam urna tempor.

% Paragraph 3 – Procedure: 
% Outline the step‑by‑step execution (pipelines run, tasks assigned, scripts used), emphasizing chronology and reproducibility.
Lorem ipsum dolor sit amet consectetur adipiscing elit. 
Quisque faucibus ex sapien vitae pellentesque sem placerat. 
In id cursus mi pretium tellus duis convallis. 
Tempus leo eu aenean sed diam urna tempor.

% Paragraph 4 – Measures & analysis plan: 
% Enumerate metrics and constructs, describe how they are calculated, and specify statistical/analytical methods used to answer RQ1.
Lorem ipsum dolor sit amet consectetur adipiscing elit. 
Quisque faucibus ex sapien vitae pellentesque sem placerat. 
In id cursus mi pretium tellus duis convallis. 
Tempus leo eu aenean sed diam urna tempor.

% Threats to Validity for RQ1:
% Short paragraph. Note key RQ1‑specific threats (e.g., confounding, sampling, measurement) and point to the global Threats section for deeper analysis.
\myparagraph{Threats to Validity for RQ1}
Lorem ipsum dolor sit amet consectetur adipiscing elit. 
Quisque faucibus ex sapien vitae pellentesque sem placerat. 
In id cursus mi pretium tellus duis convallis. 

% --- SUBSECTION: EMPIRICAL FINDINGS FOR RQ1 ---
\subsection{RQ1 Empirical Findings}
\label{subsec:rq1_findings}

% Paragraph 1 – Overview and preprocessing: 
% Describe the subset of data used for RQ1, any preprocessing, exclusions, and basic descriptive stats (e.g., counts, ranges).
Lorem ipsum dolor sit amet consectetur adipiscing elit. 
Quisque faucibus ex sapien vitae pellentesque sem placerat. 
In id cursus mi pretium tellus duis convallis. 
Tempus leo eu aenean sed diam urna tempor.

% Paragraphs 2–3 – Main findings with figures/tables: 
% Present the core results in logical order (e.g., primary metrics, distributions, comparative stats). Every figure/table referenced must be explicitly cited in the text (‘As shown in Figure X…’). No interpretation beyond factual patterns.
Lorem ipsum dolor sit amet consectetur adipiscing elit. 
Quisque faucibus ex sapien vitae pellentesque sem placerat.
In id cursus mi pretium tellus duis convallis. 
Tempus leo eu aenean sed diam urna tempor.
Pulvinar vivamus fringilla lacus nec metus bibendum egestas. 
Iaculis massa nisl malesuada lacinia integer nunc posuere.
Ut hendrerit semper vel class aptent taciti sociosqu. 

% Paragraph 4 – Trends, patterns, and null results: 
% Summarize observed trends, highlight interesting edge cases, and explicitly mention negative or null results. Stay descriptive; save interpretation for the Discussion.
Lorem ipsum dolor sit amet consectetur adipiscing elit. 
Quisque faucibus ex sapien vitae pellentesque sem placerat. 
In id cursus mi pretium tellus duis convallis. 
Tempus leo eu aenean sed diam urna tempor.

% Formal ‘Answer to RQ1’ box: 
% 2–3 sentences of neutral synthesis directly answering the RQ text, with no speculation.
% =========================================================================
% EMPIRICAL FINDING DISPLAY BOX
% DEPENDENCY WARNING: This structural container relies directly on the 
% \newfinding command (defined in typeset/misc.tex). 
% Calling \newfinding inside this block:
%   1. Increments the central sequential 'finding' counter register.
%   2. Globally registers a target link for downstream \label and \ref calls.
%   3. Injects the bold string "Finding X." directly into the box canvas stream.
% =========================================================================

\begin{tcolorbox}[
    colback=yellow!30,    % Warm yellow alert canvas interior shading tint
    colframe=black,       % High-contrast solid black border outline
    sharp corners,        % Disables rounded bounding radii to mimic a standard \fcolorbox
    boxrule=0.5pt,        % Ultra-clean frame line width boundary scale
    % Comprehensive layout padding overrides to prevent text edge crowding:
    left=2mm, 
    right=2mm, 
    top=2mm, 
    bottom=5pt, 
    width=0.98\linewidth  % Consistently centers and bounds the box safely inside column margins
]
    % Note: Typically, you will invoke only ONE unique \newfinding call per box container.
    % Multiple consecutive calls here are used to demonstrate sequence step increments.
    \newfinding In id cursus mi pretium tellus duis convallis.
    
    \newfinding Pulvinar vivamus fringilla lacus nec metus bibendum.
    
    \newfinding Iaculis massa nisl malesuada lacinia integer.
    \label{finding: template-finding}
\end{tcolorbox}

% Pointer sentence to Discussion: 
% Single sentence that forwards the reader to the Discussion for qualitative interpretation and implications.
Lorem ipsum dolor sit amet consectetur adipiscing elit.

% =========================================================================
% SECTION: EVALUATION OF RESEARCH QUESTION 2 (RQ2)
% RQ2 block arc: overview & motivation → RQ2‑specific methodology → RQ2 findings → formal answer box → pointer to Discussion.
% =========================================================================
\section{Evaluation of RQ2: [Dimension Title, e.g., Patterns]}
\label{sec:rq2_evaluation}

% --- RQ2 OVERVIEW PARAGRAPH ---
% Target: 6–8 sentences.
% Content: Restate RQ2 verbatim, explain its role relative to the other RQs (it
%   builds on RQ1's baseline by characterizing patterns/mechanics), and explain
%   why this dimension matters. End with a sentence introducing the methodology.
Lorem ipsum dolor sit amet consectetur adipiscing elit.
Quisque faucibus ex sapien vitae pellentesque sem placerat. 
In id cursus mi pretium tellus duis convallis.
Tempus leo eu aenean sed diam urna tempor. 
Pulvinar vivamus fringilla lacus nec metus bibendum egestas.
Iaculis massa nisl malesuada lacinia integer nunc posuere. 
Ut hendrerit semper vel class aptent taciti sociosqu. 
Ad litora torquent per conubia nostra inceptos himenaeos.

% --- SUBSECTION: METHODOLOGY FOR RQ2 ---
\subsection{RQ2 Methodology and Categorization Pipeline}
\label{subsec:rq2_methodology}
% RULE: cite every figure in-text via \Cref{fig:...}; never rely on placement alone.
% =========================================================================
% TEMPLATE FIGURE PLACEHOLDER
% Uses a standard boxed paragraph container to stand in for missing visual assets.
% =========================================================================

\begin{figure}[h] % [h] requests placement "here" inline with the surrounding text stream
    \centering    % Horizontally centers the figure artifact within the column/page margins
    
    % \fbox draws a visible frame border around its inner content block
    \fbox{
        % \parbox structure creates an explicit, fixed-size paragraph boundary container
        % Arguments format: [inner-alignment][height][content-alignment]{width}
        \parbox[c][2in][c]{3in}{
            \centering Missing Figure % Centered placeholder text string inside the bounding box
        }
    }
    
    % Core descriptive label rendered immediately beneath the figure box frame
    \caption{Template figure placeholder}
    
    % Unique cross-referencing anchor tag; reference this asset elsewhere via \Cref{fig:template_figure}
    \label{fig:template_figure}
\end{figure}

% Paragraph 1 – Design & rationale:
% Explain the study design (experiment, case study, mining, survey) and why it
% is appropriate for RQ2.
Lorem ipsum dolor sit amet consectetur adipiscing elit.
Quisque faucibus ex sapien vitae pellentesque sem placerat. 
In id cursus mi pretium tellus duis convallis. 
Tempus leo eu aenean sed diam urna tempor.

% Paragraph 2 – Data sources & sampling:
% Describe datasets, projects, participants, and inclusion/exclusion rules.
% Provide enough detail for replication.
Lorem ipsum dolor sit amet consectetur adipiscing elit.
Quisque faucibus ex sapien vitae pellentesque sem placerat. 
In id cursus mi pretium tellus duis convallis. 
Tempus leo eu aenean sed diam urna tempor.

% Paragraph 3 – Procedure:
% Outline the step-by-step execution (pipelines run, tasks assigned, scripts
% used), emphasizing chronology and reproducibility.
Lorem ipsum dolor sit amet consectetur adipiscing elit.
Quisque faucibus ex sapien vitae pellentesque sem placerat. 
In id cursus mi pretium tellus duis convallis. 
Tempus leo eu aenean sed diam urna tempor.

% Paragraph 4 – Measures & analysis plan:
% Enumerate metrics and constructs, describe how they are calculated, and specify
% the statistical/analytical methods used to answer RQ2.
Lorem ipsum dolor sit amet consectetur adipiscing elit.
Quisque faucibus ex sapien vitae pellentesque sem placerat. 
In id cursus mi pretium tellus duis convallis. 
Tempus leo eu aenean sed diam urna tempor.

% Threats to Validity for RQ2:
% Short paragraph (2–3 sentences). Note key RQ2-specific threats (e.g.,
% confounding, sampling, measurement) and point to the global Threats section.
\myparagraph{Threats to Validity for RQ2}
Lorem ipsum dolor sit amet consectetur adipiscing elit.
Quisque faucibus ex sapien vitae pellentesque sem placerat. 
In id cursus mi pretium tellus duis convallis. 

% --- SUBSECTION: EMPIRICAL FINDINGS FOR RQ2 ---
\subsection{RQ2 Empirical Findings}
\label{subsec:rq2_findings}

% Paragraph 1 – Overview and preprocessing:
% Describe the subset of data used for RQ2, any preprocessing, exclusions, and
% basic descriptive stats (e.g., counts, ranges).
Lorem ipsum dolor sit amet consectetur adipiscing elit.
Quisque faucibus ex sapien vitae pellentesque sem placerat. 
In id cursus mi pretium tellus duis convallis. 
Tempus leo eu aenean sed diam urna tempor.

% Paragraphs 2–3 – Main findings with figures/tables:
% Present the core results in logical order. Every figure/table referenced must
% be explicitly cited in the text ("As shown in \Cref{fig:...}"). Stay strictly
% descriptive; no interpretation beyond factual patterns.
Lorem ipsum dolor sit amet consectetur adipiscing elit. 
Quisque faucibus ex sapien vitae pellentesque sem placerat.
In id cursus mi pretium tellus duis convallis. 
Tempus leo eu aenean sed diam urna tempor.
Pulvinar vivamus fringilla lacus nec metus bibendum egestas. 
Iaculis massa nisl malesuada lacinia integer nunc posuere.
Ut hendrerit semper vel class aptent taciti sociosqu. 

% Paragraph 4 – Trends, patterns, and null results:
% Summarize observed trends, highlight interesting edge cases, and explicitly
% mention negative or null results. Stay descriptive; save interpretation for the
% Discussion.
Lorem ipsum dolor sit amet consectetur adipiscing elit.
Quisque faucibus ex sapien vitae pellentesque sem placerat. 
In id cursus mi pretium tellus duis convallis. 
Tempus leo eu aenean sed diam urna tempor. 

% Formal 'Answer to RQ2' box:
% 2–3 sentences of neutral synthesis directly answering the RQ text, no speculation.
% =========================================================================
% EMPIRICAL FINDING DISPLAY BOX
% DEPENDENCY WARNING: This structural container relies directly on the 
% \newfinding command (defined in typeset/misc.tex). 
% Calling \newfinding inside this block:
%   1. Increments the central sequential 'finding' counter register.
%   2. Globally registers a target link for downstream \label and \ref calls.
%   3. Injects the bold string "Finding X." directly into the box canvas stream.
% =========================================================================

\begin{tcolorbox}[
    colback=yellow!30,    % Warm yellow alert canvas interior shading tint
    colframe=black,       % High-contrast solid black border outline
    sharp corners,        % Disables rounded bounding radii to mimic a standard \fcolorbox
    boxrule=0.5pt,        % Ultra-clean frame line width boundary scale
    % Comprehensive layout padding overrides to prevent text edge crowding:
    left=2mm, 
    right=2mm, 
    top=2mm, 
    bottom=5pt, 
    width=0.98\linewidth  % Consistently centers and bounds the box safely inside column margins
]
    % Note: Typically, you will invoke only ONE unique \newfinding call per box container.
    % Multiple consecutive calls here are used to demonstrate sequence step increments.
    \newfinding In id cursus mi pretium tellus duis convallis.
    
    \newfinding Pulvinar vivamus fringilla lacus nec metus bibendum.
    
    \newfinding Iaculis massa nisl malesuada lacinia integer.
    \label{finding: template-finding}
\end{tcolorbox}

% Pointer sentence to Discussion:
% Single sentence forwarding the reader to the Discussion for qualitative
% interpretation and implications.
Lorem ipsum dolor sit amet consectetur adipiscing elit.

% =========================================================================
% SECTION: EVALUATION OF RESEARCH QUESTION 3 (RQ3)
% RQ3 block arc: overview & motivation → RQ3‑specific methodology → RQ3 findings → formal answer box → pointer to Discussion.
% =========================================================================
\section{Evaluation of RQ3: [Dimension Title, e.g., Impact]}
\label{sec:rq3_evaluation}

% --- RQ3 OVERVIEW PARAGRAPH ---
% Target: 6–8 sentences.
% Content: Restate RQ3 verbatim, explain its role relative to the other RQs (it
%   typically caps the arc by assessing impact/consequence), and explain why this
%   dimension matters. End with a sentence introducing the methodology subsection.
Lorem ipsum dolor sit amet consectetur adipiscing elit.
Quisque faucibus ex sapien vitae pellentesque sem placerat. 
In id cursus mi pretium tellus duis convallis.
Tempus leo eu aenean sed diam urna tempor. 
Pulvinar vivamus fringilla lacus nec metus bibendum egestas.
Iaculis massa nisl malesuada lacinia integer nunc posuere. 
Ut hendrerit semper vel class aptent taciti sociosqu. 
Ad litora torquent per conubia nostra inceptos himenaeos.

% --- SUBSECTION: METHODOLOGY FOR RQ3 ---
\subsection{RQ3 Methodology and Analytical Matrix}
\label{subsec:rq3_methodology}
% RULE: cite every figure in-text via \Cref{fig:...}; never rely on placement alone.
% =========================================================================
% TEMPLATE FIGURE PLACEHOLDER
% Uses a standard boxed paragraph container to stand in for missing visual assets.
% =========================================================================

\begin{figure}[h] % [h] requests placement "here" inline with the surrounding text stream
    \centering    % Horizontally centers the figure artifact within the column/page margins
    
    % \fbox draws a visible frame border around its inner content block
    \fbox{
        % \parbox structure creates an explicit, fixed-size paragraph boundary container
        % Arguments format: [inner-alignment][height][content-alignment]{width}
        \parbox[c][2in][c]{3in}{
            \centering Missing Figure % Centered placeholder text string inside the bounding box
        }
    }
    
    % Core descriptive label rendered immediately beneath the figure box frame
    \caption{Template figure placeholder}
    
    % Unique cross-referencing anchor tag; reference this asset elsewhere via \Cref{fig:template_figure}
    \label{fig:template_figure}
\end{figure}

% Paragraph 1 – Design & rationale:
% Explain the study design (experiment, case study, mining, survey) and why it
% is appropriate for RQ3.
Lorem ipsum dolor sit amet consectetur adipiscing elit.
Quisque faucibus ex sapien vitae pellentesque sem placerat. 
In id cursus mi pretium tellus duis convallis. 
Tempus leo eu aenean sed diam urna tempor.

% Paragraph 2 – Data sources & sampling:
% Describe datasets, projects, participants, and inclusion/exclusion rules.
% Provide enough detail for replication.
Lorem ipsum dolor sit amet consectetur adipiscing elit.
Quisque faucibus ex sapien vitae pellentesque sem placerat. 
In id cursus mi pretium tellus duis convallis. 
Tempus leo eu aenean sed diam urna tempor. 

% Paragraph 3 – Procedure:
% Outline the step-by-step execution (pipelines run, tasks assigned, scripts
% used), emphasizing chronology and reproducibility.
Lorem ipsum dolor sit amet consectetur adipiscing elit.
Quisque faucibus ex sapien vitae pellentesque sem placerat. 
In id cursus mi pretium tellus duis convallis. 
Tempus leo eu aenean sed diam urna tempor.

% Paragraph 4 – Measures & analysis plan:
% Enumerate metrics and constructs, describe how they are calculated, and specify
% the statistical/analytical methods used to answer RQ3.
Lorem ipsum dolor sit amet consectetur adipiscing elit.
Quisque faucibus ex sapien vitae pellentesque sem placerat. 
In id cursus mi pretium tellus duis convallis. 
Tempus leo eu aenean sed diam urna tempor.

% Threats to Validity for RQ3:
% Short paragraph (2–3 sentences). Note key RQ3-specific threats (e.g.,
% confounding, sampling, measurement) and point to the global Threats section.
\myparagraph{Threats to Validity for RQ3}
Lorem ipsum dolor sit amet consectetur adipiscing elit.
Quisque faucibus ex sapien vitae pellentesque sem placerat. 
In id cursus mi pretium tellus duis convallis. 

% --- SUBSECTION: EMPIRICAL FINDINGS FOR RQ3 ---
\subsection{RQ3 Empirical Findings}
\label{subsec:rq3_findings}

% Paragraph 1 – Overview and preprocessing:
% Describe the subset of data used for RQ3, any preprocessing, exclusions, and
% basic descriptive stats (e.g., counts, ranges).
Lorem ipsum dolor sit amet consectetur adipiscing elit.
Quisque faucibus ex sapien vitae pellentesque sem placerat. 
In id cursus mi pretium tellus duis convallis. 
Tempus leo eu aenean sed diam urna tempor.

% Paragraphs 2–3 – Main findings with figures/tables:
% Present the core results in logical order. Every figure/table referenced must
% be explicitly cited in the text ("As shown in \Cref{fig:...}"). Stay strictly
% descriptive; no interpretation beyond factual patterns.
Lorem ipsum dolor sit amet consectetur adipiscing elit. 
Quisque faucibus ex sapien vitae pellentesque sem placerat.
In id cursus mi pretium tellus duis convallis. 
Tempus leo eu aenean sed diam urna tempor.
Pulvinar vivamus fringilla lacus nec metus bibendum egestas. 
Iaculis massa nisl malesuada lacinia integer nunc posuere.
Ut hendrerit semper vel class aptent taciti sociosqu. 

% Paragraph 4 – Trends, patterns, and null results:
% Summarize observed trends, highlight interesting edge cases, and explicitly
% mention negative or null results. Stay descriptive; save interpretation for the
% Discussion.
Lorem ipsum dolor sit amet consectetur adipiscing elit.
Quisque faucibus ex sapien vitae pellentesque sem placerat. 
In id cursus mi pretium tellus duis convallis. 
Tempus leo eu aenean sed diam urna tempor. 

% Formal 'Answer to RQ3' box:
% 2–3 sentences of neutral synthesis directly answering the RQ text, no speculation.
% =========================================================================
% EMPIRICAL FINDING DISPLAY BOX
% DEPENDENCY WARNING: This structural container relies directly on the 
% \newfinding command (defined in typeset/misc.tex). 
% Calling \newfinding inside this block:
%   1. Increments the central sequential 'finding' counter register.
%   2. Globally registers a target link for downstream \label and \ref calls.
%   3. Injects the bold string "Finding X." directly into the box canvas stream.
% =========================================================================

\begin{tcolorbox}[
    colback=yellow!30,    % Warm yellow alert canvas interior shading tint
    colframe=black,       % High-contrast solid black border outline
    sharp corners,        % Disables rounded bounding radii to mimic a standard \fcolorbox
    boxrule=0.5pt,        % Ultra-clean frame line width boundary scale
    % Comprehensive layout padding overrides to prevent text edge crowding:
    left=2mm, 
    right=2mm, 
    top=2mm, 
    bottom=5pt, 
    width=0.98\linewidth  % Consistently centers and bounds the box safely inside column margins
]
    % Note: Typically, you will invoke only ONE unique \newfinding call per box container.
    % Multiple consecutive calls here are used to demonstrate sequence step increments.
    \newfinding In id cursus mi pretium tellus duis convallis.
    
    \newfinding Pulvinar vivamus fringilla lacus nec metus bibendum.
    
    \newfinding Iaculis massa nisl malesuada lacinia integer.
    \label{finding: template-finding}
\end{tcolorbox}

% Pointer sentence to Discussion:
% Single sentence forwarding the reader to the Discussion for qualitative
% interpretation and implications.
Lorem ipsum dolor sit amet consectetur adipiscing elit.
Quisque faucibus ex sapien vitae pellentesque sem placerat. 
In id cursus mi pretium tellus duis convallis. 
Tempus leo eu aenean sed diam urna tempor. 

% =========================================================================
% SECTION: DISCUSSION
% Narrative Flow: Narrow (Interpretations) -> Broad (Implications) 
%                -> Defensive (Limitations) -> Forward-Looking (Future Work)
% =========================================================================
\section{Discussion}
\label{sec:discussion}

% --- PARAGRAPH 1: DISCUSSION OPENING ---
% Target: 7–9 sentences. 
% Content: Restate the high‑level takeaway of the study, tying back to the Introduction’s big gap and summarizing the overall pattern across RQ1–RQ3. This paragraph sets the tone for the rest of the Discussion.
Lorem ipsum dolor sit amet consectetur adipiscing elit. 
Quisque faucibus ex sapien vitae pellentesque sem placerat. 
In id cursus mi pretium tellus duis convallis. 
Tempus leo eu aenean sed diam urna tempor. 
Pulvinar vivamus fringilla lacus nec metus bibendum egestas. 
Iaculis massa nisl malesuada lacinia integer nunc posuere. 
Ut hendrerit semper vel class aptent taciti sociosqu. 

% --- SUBSECTION: INTERPRETATION OF RQS ---
% Per‑RQ mini‑discussions: each \myparagraph is 8–12 sentences total, structured as interpretation → mechanism → relation to prior work → scoped implication/limit.
\subsection{Interpretation of Research Questions}
\label{subsec:interpretations}

\myparagraph{RQ$_1$ -- Baseline Prevalence Metrics}
% Interpretation: 2–3 sentences summarizing how the findings answer RQ1.
Lorem ipsum dolor sit amet consectetur adipiscing elit. Quisque faucibus ex sapien vitae pellentesque sem placerat. In id cursus mi pretium tellus duis convallis.
% Mechanism/Why: 2–3 sentences explaining why these patterns likely occur, grounded in theory or domain knowledge.
Lorem ipsum dolor sit amet consectetur adipiscing elit. Quisque faucibus ex sapien vitae pellentesque sem placerat. In id cursus mi pretium tellus duis convallis.
% Relation to prior work: 2–3 sentences comparing RQ1 findings to specific prior studies; note agreements and contradictions.
Lorem ipsum dolor sit amet consectetur adipiscing elit. Quisque faucibus ex sapien vitae pellentesque sem placerat. In id cursus mi pretium tellus duis convallis.
% Scoped actionable/limit: 2–3 sentences stating implications or recommendations for practice, plus one key limitation that tempers those recommendations.
Lorem ipsum dolor sit amet consectetur adipiscing elit. Quisque faucibus ex sapien vitae pellentesque sem placerat. In id cursus mi pretium tellus duis convallis. Tempus leo eu aenean sed diam urna tempor.

\myparagraph{RQ$_2$ -- The Mechanics of Fragile Reuse}
% Interpretation: 2–3 sentences summarizing how the findings answer RQ2.
Lorem ipsum dolor sit amet consectetur adipiscing elit.
% Mechanism/Why: 2–3 sentences explaining why these patterns likely occur, grounded in theory or domain knowledge.
Lorem ipsum dolor sit amet consectetur adipiscing elit. Quisque faucibus ex sapien vitae pellentesque sem placerat. In id cursus mi pretium tellus duis convallis.
% Relation to prior work: 2–3 sentences comparing RQ2 findings to specific prior studies.
Lorem ipsum dolor sit amet consectetur adipiscing elit. Quisque faucibus ex sapien vitae pellentesque sem placerat. In id cursus mi pretium tellus duis convallis.
% Scoped actionable/limit: 2–3 sentences stating implications or recommendations for practice.
Lorem ipsum dolor sit amet consectetur adipiscing elit. Quisque faucibus ex sapien vitae pellentesque sem placerat. In id cursus mi pretium tellus duis convallis. Tempus leo eu aenean sed diam urna tempor.

\myparagraph{RQ$_3$ -- Quantification of Pipeline Half-Life}
% Interpretation: 2–3 sentences summarizing how the findings answer RQ3.
Lorem ipsum dolor sit amet consectetur adipiscing elit.
% Mechanism/Why: 2–3 sentences explaining why these patterns likely occur, grounded in theory or domain knowledge.
Lorem ipsum dolor sit amet consectetur adipiscing elit. Quisque faucibus ex sapien vitae pellentesque sem placerat. In id cursus mi pretium tellus duis convallis.
% Relation to prior work: 2–3 sentences comparing RQ3 findings to specific prior studies.
Lorem ipsum dolor sit amet consectetur adipiscing elit. Quisque faucibus ex sapien vitae pellentesque sem placerat. In id cursus mi pretium tellus duis convallis.
% Scoped actionable/limit: 2–3 sentences stating implications or recommendations for practice.
Lorem ipsum dolor sit amet consectetur adipiscing elit. Quisque faucibus ex sapien vitae pellentesque sem placerat. In id cursus mi pretium tellus duis convallis. Tempus leo eu aenean sed diam urna tempor.

% --- SUBSECTION: WIDER IMPLICATIONS ---
\subsection{Wider Implications for Science and Engineering}
\label{subsec:implications}

% Paragraph 1 – Theoretical implications: 
% Target: 5–7 sentences. 
% Content: Synthesize across RQs to state how the study refines or challenges existing theories/models. Focus on general principles.
Lorem ipsum dolor sit amet consectetur adipiscing elit. 
Quisque faucibus ex sapien vitae pellentesque sem placerat. 
In id cursus mi pretium tellus duis convallis. 
Tempus leo eu aenean sed diam urna tempor. 
Pulvinar vivamus fringilla lacus nec metus bibendum egestas. 
Iaculis massa nisl malesuada lacinia integer nunc posuere. 
Ut hendrerit semper vel class aptent taciti sociosqu.

% Paragraph 2 – Practical/industrial implications: 
% Target: 5–7 sentences. 
% Content: Translate findings into concrete recommendations for teams, tools, and processes, mentioning feasibility and constraints.
Lorem ipsum dolor sit amet consectetur adipiscing elit. 
Quisque faucibus ex sapien vitae pellentesque sem placerat. 
In id cursus mi pretium tellus duis convallis. 
Tempus leo eu aenean sed diam urna tempor. 
Pulvinar vivamus fringilla lacus nec metus bibendum egestas. 
Iaculis massa nisl malesuada lacinia integer nunc posuere. 
Ut hendrerit semper vel class aptent taciti sociosqu.

% Optional 3rd Paragraph – Policy / Educational Frameworks:
% NOTE: Optional, beyond the two standard implication types (Theoretical and
%   Practical/Industrial) defined by the style guide. Delete if not needed.
% Target: 5–7 sentences on implications for standards, curricula, or broader scientific practice.
Lorem ipsum dolor sit amet consectetur adipiscing elit. 
Quisque faucibus ex sapien vitae pellentesque sem placerat. 
In id cursus mi pretium tellus duis convallis. 
Tempus leo eu aenean sed diam urna tempor. 
Pulvinar vivamus fringilla lacus nec metus bibendum egestas. 

% --- SUBSECTION: LIMITATIONS AND BIASES ---
\subsection{Limitations and Alternative Explanations}
\label{subsec:limitations}

% Paragraph 1 – Main limitations: 
% Target: 5–7 sentences. 
% Content: Summarize key limitations across design, data, and measures; explain how each might bias results or limit generality.
Lorem ipsum dolor sit amet consectetur adipiscing elit. 
Quisque faucibus ex sapien vitae pellentesque sem placerat. 
In id cursus mi pretium tellus duis convallis. 
Tempus leo eu aenean sed diam urna tempor. 
Pulvinar vivamus fringilla lacus nec metus bibendum egestas. 

% Paragraph 2 – Alternative explanations: 
% Target: 5–7 sentences. 
% Content: Present plausible alternative interpretations of major findings and briefly note what further data/methods would be needed to test them.
Lorem ipsum dolor sit amet consectetur adipiscing elit. 
Quisque faucibus ex sapien vitae pellentesque sem placerat. 
In id cursus mi pretium tellus duis convallis. 
Tempus leo eu aenean sed diam urna tempor. 
Pulvinar vivamus fringilla lacus nec metus bibendum egestas. 

% --- SUBSECTION: FUTURE DIRECTIONS ---
\subsection{Future Research Trajectories}
\label{subsec:future_directions}

% Paragraph 1 – Next logical steps: 
% Target: 4–6 sentences. 
% Content: Propose immediate follow‑up studies (replications, scaling, boundary conditions), each explicitly linked to a specific result or limitation.
Lorem ipsum dolor sit amet consectetur adipiscing elit. 
Quisque faucibus ex sapien vitae pellentesque sem placerat. 
In id cursus mi pretium tellus duis convallis. 
Tempus leo eu aenean sed diam urna tempor. 

% Paragraph 2 – Methodological refinements: 
% Target: 3–5 sentences. 
% Content: Suggest improvements to design, metrics, or data collection that would strengthen future evidence.
Lorem ipsum dolor sit amet consectetur adipiscing elit. 
Quisque faucibus ex sapien vitae pellentesque sem placerat. 
In id cursus mi pretium tellus duis convallis. 

% Paragraph 3 – New questions: 
% Target: 4–6 sentences. 
% Content: Outline new research questions opened by unexpected findings or implications; explain why each is important for empirical SE.
Lorem ipsum dolor sit amet consectetur adipiscing elit. 
Quisque faucibus ex sapien vitae pellentesque sem placerat. 
In id cursus mi pretium tellus duis convallis. 
Tempus leo eu aenean sed diam urna tempor. 

% Final closing block – Take-Home Message:
% Target: 4–6 sentences. 
% Content: Provide a synthesized, confident statement of what the study ultimately shows and why it matters. No new ideas or data; this should echo the core contribution and high‑level impact.
\myparagraph{Take-Home Message}
Lorem ipsum dolor sit amet consectetur adipiscing elit. 
Quisque faucibus ex sapien vitae pellentesque sem placerat. 
In id cursus mi pretium tellus duis convallis. 
Tempus leo eu aenean sed diam urna tempor. 
Pulvinar vivamus fringilla lacus nec metus bibendum egestas. 
Iaculis massa nisl malesuada lacinia integer nunc posuere. 

% =========================================================================
% SECTION: THREATS TO VALIDITY
% Function: Systematically classify study vulnerabilities into three subsections
%   (Internal, External, Construct validity). Be candid; each block names a threat,
%   its potential effect on results, and the mitigation taken.
% =========================================================================
\section{Threats to Validity}
\label{sec:threats}

% Section opening:
% Target: 3–4 sentences. 
% Content: Explain the purpose of this section (systematically analyzing internal, external, and construct validity threats) and note that detailed RQ‑specific threats were mentioned in each RQ methodology subsection.
Lorem ipsum dolor sit amet consectetur adipiscing elit. 
Quisque faucibus ex sapien vitae pellentesque sem placerat. 
In id cursus mi pretium tellus duis convallis. 
Our analysis follows the standard empirical software engineering classification distinguishing between internal, external, and construct validity.

% --- SUBSECTION: INTERNAL VALIDITY ---
\subsection{Internal Validity}
\label{subsec:internal_validity}

% First block (Confounding Variables):
% Target: 3–5 sentences.
% Content: Explain potential confounding variables and how they may distort causal
%   interpretations; name the mitigation strategies used. Use causal signposts
%   ("As a result", "Consequently").
Lorem ipsum dolor sit amet consectetur adipiscing elit. 
Quisque faucibus ex sapien vitae pellentesque sem placerat. 
In id cursus mi pretium tellus duis convallis. 
Tempus leo eu aenean sed diam urna tempor. 
Pulvinar vivamus fringilla lacus nec metus bibendum egestas. 

% Second block (History, Maturation, and Instrumentation):
% Target: 3–4 sentences.
% Content: Discuss history, maturation, and instrumentation effects and their
%   potential influence on outcomes; note how design choices attempt to control them.
Lorem ipsum dolor sit amet consectetur adipiscing elit. 
Quisque faucibus ex sapien vitae pellentesque sem placerat. 
In id cursus mi pretium tellus duis convallis. 
Tempus leo eu aenean sed diam urna tempor. 

% Third block (Attrition and Missing Data):
% Target: 2–3 sentences.
% Content: Describe attrition/missing data patterns and how they were handled
%   (e.g., exclusion, imputation, sensitivity analysis).
Lorem ipsum dolor sit amet consectetur adipiscing elit. 
Quisque faucibus ex sapien vitae pellentesque sem placerat. 
In id cursus mi pretium tellus duis convallis. 

% --- SUBSECTION: EXTERNAL VALIDITY ---
\subsection{External Validity}
\label{subsec:external_validity}

% First block (Context & Generalizability):
% Target: 3–4 sentences.
% Content: Clarify the study setting and domains, and explicitly state the boundary
%   conditions — where the results do and do not generalize.
Lorem ipsum dolor sit amet consectetur adipiscing elit. 
Quisque faucibus ex sapien vitae pellentesque sem placerat. 
In id cursus mi pretium tellus duis convallis. 
Tempus leo eu aenean sed diam urna tempor. 

% Second block (Sampling Bias):
% Target: 2–3 sentences.
% Content: Discuss how selection of projects/participants may skew results and how
%   that limits external validity.
Lorem ipsum dolor sit amet consectetur adipiscing elit. 
Quisque faucibus ex sapien vitae pellentesque sem placerat. 
In id cursus mi pretium tellus duis convallis. 

% --- SUBSECTION: CONSTRUCT VALIDITY ---
\subsection{Construct Validity}
\label{subsec:construct_validity}

% First block (Indirect Metrics and Operationalization):
% Target: 4–6 sentences.
% Content: Justify that the chosen metrics capture the intended constructs,
%   referencing prior work and triangulation where possible.
Lorem ipsum dolor sit amet consectetur adipiscing elit. 
Quisque faucibus ex sapien vitae pellentesque sem placerat. 
In id cursus mi pretium tellus duis convallis. 
Tempus leo eu aenean sed diam urna tempor. 
Pulvinar vivamus fringilla lacus nec metus bibendum egestas. 
Iaculis massa nisl malesuada lacinia integer nunc posuere. 
Ut hendrerit semper vel class aptent taciti sociosqu. 
Ad litora torquent per conubia nostra inceptos himenaeos.

% Second block (Measurement Error and Technical Bias):
% Target: 3–5 sentences.
% Content: Address logging inaccuracies, self-report biases, and tool limitations;
%   explain how they might affect interpretations.
Lorem ipsum dolor sit amet consectetur adipiscing elit. 
Quisque faucibus ex sapien vitae pellentesque sem placerat. 
In id cursus mi pretium tellus duis convallis. 
Tempus leo eu aenean sed diam urna tempor. 
Pulvinar vivamus fringilla lacus nec metus bibendum egestas. 
Iaculis massa nisl malesuada lacinia integer nunc posuere. 

% =========================================================================
% SECTION: CONCLUSION
% =========================================================================
\section{Conclusion}
\label{sec:conclusion}

% Paragraph 1 – Aims and findings synthesis: 
% Target: 2–4 sentences. 
% Content: Restate the overall problem and aims in fresh language, then synthesize the main findings across RQs (no detailed numbers).
Lorem ipsum dolor sit amet consectetur adipiscing elit. 
Quisque faucibus ex sapien vitae pellentesque sem placerat. 
In id cursus mi pretium tellus duis convallis. 
Tempus leo eu aenean sed diam urna tempor. 
Pulvinar vivamus fringilla lacus nec metus bibendum egestas. 
Iaculis massa nisl malesuada lacinia integer nunc posuere. 
Ut hendrerit semper vel class aptent taciti sociosqu.

% Paragraph 2 – Contributions and practical implications: 
% Target: 2–4 sentences. 
% Content: Explicitly name the key contributions and what they enable or change for researchers and practitioners.
Lorem ipsum dolor sit amet consectetur adipiscing elit. 
Quisque faucibus ex sapien vitae pellentesque sem placerat. 
In id cursus mi pretium tellus duis convallis. 
Tempus leo eu aenean sed diam urna tempor. 
Pulvinar vivamus fringilla lacus nec metus bibendum egestas. 
Iaculis massa nisl malesuada lacinia integer nunc posuere. 
Ut hendrerit semper vel class aptent taciti sociosqu.

% Paragraph 3 – Limitations, future work, final takeaway: 
% Target: 3–4 sentences. 
% Style Rule: The final sentence must be a confident, synthesized statement—not an apology.
% Content: Echo the most important limitation/future direction at a high level and end with a strong single sentence that captures the study’s core message.
Lorem ipsum dolor sit amet consectetur adipiscing elit. 
Quisque faucibus ex sapien vitae pellentesque sem placerat. 
In id cursus mi pretium tellus duis convallis. 
Tempus leo eu aenean sed diam urna tempor. 
Pulvinar vivamus fringilla lacus nec metus bibendum egestas. 
Iaculis massa nisl malesuada lacinia integer nunc posuere. 
Ut hendrerit semper vel class aptent taciti sociosqu. 

% =========================================================================
% ACKNOWLEDGMENTS
% =========================================================================
\section*{Acknowledgments}
% Target: 1–3 sentences. 
% Content: Thank funding sources, collaborators, and infrastructure. No technical content.
Lorem ipsum dolor sit amet consectetur adipiscing elit. 
Quisque faucibus ex sapien vitae pellentesque sem placerat. 
In id cursus mi pretium tellus duis convallis. 

\bibliographystyle{ACM-Reference-Format}
\bibliography{citations}    % For all other submissions
%% =========================================================================
% DOCUMENT CLASS & SETUP
% =========================================================================

% \documentclass[sigconf]{acmart}       % For final camera-ready peer-reviewed submissions
\documentclass[sigconf,nonacm]{acmart} % For arXiv/preprint: removes copyright blocks & ACM footers

% =========================================================================
% GLOBAL PARAGRAPH AND PROSE RULES
% =========================================================================
% 1. Target Length: Aim for 3–7 sentences per paragraph. Avoid single-sentence 
%    paragraphs and dense walls of text.
% 2. Functional Focus: Each paragraph must serve a singular rhetorical function 
%    (e.g., context, problem, method, finding, interpretation, implication, 
%    limitation, or future work).
% 3. Explicit Signposting: Maintain visible logical flow across paragraph 
%    boundaries using structural transitions (e.g., "However", "As a result", 
%    "Consequently", "In contrast").
% 4. Micro-Sections: Use \myparagraph{...} blocks to signify named mini-paragraphs 
%    with specialized micro-rhetorical goals (e.g., "Contributions", "Significance").

% =========================================================================
% CONDITIONAL DEBUGGING FLAGS
% =========================================================================

\newif\ifDEBUG  % Declares a boolean switch called "DEBUG" for conditional compilation
\DEBUGtrue      % Set DEBUG to true (activates draft/todo notes)
\DEBUGfalse     % Set DEBUG to false (silences draft notes for clean compilation)
% NET EFFECT: \DEBUGfalse runs last, so debug is OFF by default. While OFF,
% \TODO{}, \CITE{}, and \NMS{} (co-author notes) render as nothing.
% To see draft markers, comment out the \DEBUGfalse line above.

% =========================================================================
% MODULAR PREAMBLE IMPORTS
% =========================================================================

\input{misc/data}        % Imports global stats, variables, and data macros
\input{misc/typesetting} % Imports formatting overrides, extra packages, and layout tweaks

% =========================================================================
% CUSTOM LOGO MACROS
% =========================================================================

% Delays definition until document start to prevent package conflicts
\AtBeginDocument{%
    % Safely defines \BibTeX logo if not already defined by a loaded package
    \providecommand\BibTeX{{%
        \normalfont B\kern-0.5em{%   % Backspaces slightly after 'B'
            \scshape                 % Switches to small caps for 'ib'
            i\kern-0.25em b}\kern-0.8em\TeX}}} 
% Micro-spaces 'ib' and pulls 'TeX' close

% =========================================================================
% RIGHTS & COPYRIGHT MANAGEMENT (Suppressed while 'nonacm' is active)
% =========================================================================

\setcopyright{acmlicensed}          % Specifies the ACM licensing agreement type
\copyrightyear{2026}                % Sets the explicit copyright year
\acmYear{2026}                      % Sets the publishing year index
\acmDOI{XXXXXXX.XXXXXXX}            % Placeholder for the Digital Object Identifier string
\acmISBN{978-1-4503-XXXX-X/18/06}   % Placeholder for the standard publication ISBN

% Defines conference acronym, full metadata description, dates, and geographic location
\acmConference[Conference acronym 'XX]{Make sure to enter the correct
  conference title from your rights confirmation email}{June 03--05,
  2018}{Woodstock, NY}

% =========================================================================
% DOCUMENT CONTENT START
% =========================================================================

\begin{document}

\title{\textbf{INSERT TITLE}} % Formats the core paper title in bold face

% -------------------------------------------------------------------------
% AUTHOR METADATA
% -------------------------------------------------------------------------

\author{Nicholas M. Synovic}
\email{nsynovic@luc.edu}
\orcid{0000-0003-0413-4594} % Unique digital author identifier
\affiliation{
  \institution{Loyola University Chicago}
  \streetaddress{1032 W Sheridan Rd}
  \city{Chicago}
  \state{Illinois}
  \country{USA}
  \postcode{60660}
}

% Replaces the full author list in page running headers with a concise version to save space
\renewcommand{\shortauthors}{Synovic}

% -------------------------------------------------------------------------
% ABSTRACT SECTION
% -------------------------------------------------------------------------

% ABSTRACT: one self-contained, three-part narrative block (6–8 sentences total).
% HARD RULES: no citations and no raw statistical tables/numbers in the abstract.
\begin{abstract}
% Part 1 — Context & Problem (Sentences 1–3):
% Establish the broad SE domain and why it matters, then pivot to the core
% friction point or consequence. Function: context → problem.
Lorem ipsum dolor sit amet consectetur adipiscing elit. 
Quisque faucibus ex sapien vitae pellentesque sem placerat. 
In id cursus mi pretium tellus duis convallis. 

% Part 2 — Approach & Scope (Sentences 4–5):
% State "In this paper, we…" and summarize the design/method at one-sentence
% granularity, including evaluation scope (data size, system types, participants).
Tempus leo eu aenean sed diam urna tempor. 
Pulvinar vivamus fringilla lacus nec metus bibendum egestas. 

% Part 3 — Impact & Insights (Sentences 6–8):
% Deliver the headline qualitative/quantitative outcome, the supporting insight,
% and one concrete implication for practice or future research.
% Keep it high-level: no detailed numbers or statistics.
Iaculis massa nisl malesuada lacinia integer nunc posuere. 
Ut hendrerit semper vel class aptent taciti sociosqu. 
Ad litora torquent per conubia nostra inceptos himenaeos.
\end{abstract}

% =========================================================================
% ACM COMPUTING CLASSIFICATION SYSTEM (https://dl.acm.org/ccs)
% =========================================================================
\begin{CCSXML}
\end{CCSXML}

% The \ccsdesc macros render XML classifications into the document footer
\ccsdesc[500]{Software and its engineering}

% =========================================================================
% KEYWORDS INDEXING
% =========================================================================

\keywords{Replace me}

\maketitle

% =========================================================================
% SECTION: INTRODUCTION
% Function: Sell the paper. Six-paragraph arc; each block has ONE rhetorical job.
% Arc (one block each): broad context/hook → narrowed gap → goal & core insight
%   → contributions (bulleted) → significance/scope → roadmap.
% Rules: 3–7 sentences per paragraph; explicit transitions between blocks
%   ("However", "As a result"); no unquantified hype ("novel", "significant").
% =========================================================================
% RULE: every figure must be cited in-text via \Cref{fig:...} (e.g. \Cref{fig:template_figure}).
\input{figures/template_figure}
\section{Introduction}
\label{sec:introduction}

% --- PARAGRAPH 1: CONTEXT AND HOOK ---
% Function: Intro paragraph 1.
% Target: 3–6 sentences.
% Content: Establish the broad software engineering domain, why it matters (reliability, cost, UX, scientific impact), and a hook that shows the domain is non‑trivial. No detailed technicalities here.
Lorem ipsum dolor sit amet consectetur adipiscing elit.
Quisque faucibus ex sapien vitae pellentesque sem placerat. 
In id cursus mi pretium tellus duis convallis.
Tempus leo eu aenean sed diam urna tempor. 
Pulvinar vivamus fringilla lacus nec metus bibendum egestas.
Iaculis massa nisl malesuada lacinia integer nunc posuere. 

% --- PARAGRAPH 2: NARROWED PROBLEM AND GAP ---
% Function: Intro paragraph 2.
% Target: 4–7 sentences.
% Content: Summarize current practice/prior work, then pivot with a clear ‘However/Yet/Despite this’ sentence that states the precise gap or friction point. End by emphasizing why this gap is consequential.
Lorem ipsum dolor sit amet consectetur adipiscing elit. 
Quisque faucibus ex sapien vitae pellentesque sem placerat.
In id cursus mi pretium tellus duis convallis. 
Tempus leo eu aenean sed diam urna tempor.
Pulvinar vivamus fringilla lacus nec metus bibendum egestas. 
Iaculis massa nisl malesuada lacinia integer nunc posuere.
Ut hendrerit semper vel class aptent taciti sociosqu. 

% --- PARAGRAPH 3: GOAL AND APPROACH (THE CORE INSIGHT) ---
% Function: Intro paragraph 3.
% Target: 3–6 sentences.
% Content: Present the overall research goal and high‑level approach (system, study, framework). Include plain‑language intuition about why your approach should address the stated gap.
Lorem ipsum dolor sit amet consectetur adipiscing elit. 
Quisque faucibus ex sapien vitae pellentesque sem placerat.
In id cursus mi pretium tellus duis convallis. 
Tempus leo eu aenean sed diam urna tempor.
Pulvinar vivamus fringilla lacus nec metus bibendum egestas. 
Iaculis massa nisl malesuada lacinia integer nunc posuere.

% --- PARAGRAPH 4: CONTRIBUTIONS ---
% Structure: Use 3–4 bullets, each describing one type of contribution (artifact, empirical study, insight, open science). 
% Formatting Rule: Keep each item one sentence plus a short clarifying clause, avoiding vague phrases (‘novel’, ‘significant’) without specifics.
\myparagraph{Contributions}
This paper makes the following core contributions:
\begin{itemize}[leftmargin=16pt, rightmargin=5pt, topsep=0pt]
    % Contribution 1: The Artifact (Tool / Framework / Dataset)
    \item Lorem ipsum dolor sit amet consectetur adipiscing elit.
    % Contribution 2: The Empirical Scale & Matrix Analysis
    \item Lorem ipsum dolor sit amet consectetur adipiscing elit.
    % Contribution 3: The Findings, Open Science & Replication Package
    \item Lorem ipsum dolor sit amet consectetur adipiscing elit.
\end{itemize}

% --- PARAGRAPH 5: SIGNIFICANCE & SCOPE ---
% Function: Intro paragraph 5 — "headline results + scope" preview (required).
% Target: 3–5 sentences.
% Content: Provide 1–2 headline quantitative or qualitative outcomes, plus a brief statement of scope/assumptions (where it works, where it might not). Do not repeat details from the Results; this is a teaser, not the findings.
\myparagraph{Significance}
Lorem ipsum dolor sit amet consectetur adipiscing elit. 
Quisque faucibus ex sapien vitae pellentesque sem placerat.
In id cursus mi pretium tellus duis convallis. 
Tempus leo eu aenean sed diam urna tempor.
Pulvinar vivamus fringilla lacus nec metus bibendum egestas. 

% --- PARAGRAPH 6: PAPER ROADMAP ---
% Function: Manuscript outline (Intro paragraph 6).
% Target: 2–4 sentences.
% Content: Walk the reader through the paper in document order so the roadmap
%   doubles as a structural map. The actual order is:
%   Background (sec:background) → Research Questions (sec:research_questions)
%   → RQ1 (sec:rq1_evaluation) → RQ2 (sec:rq2_evaluation)
%   → RQ3 (sec:rq3_evaluation) → Discussion (sec:discussion)
%   → Threats (sec:threats) → Conclusion (sec:conclusion).
%   NOTE: there is no standalone "Methodology" section — methodology is embedded
%   inside each RQ section. Use lowercase section names + \Cref macros; never
%   hardcode numbers.
\myparagraph{Paper Outline}
This paper is organized as follows: 
\Cref{sec:background} provides essential background and reviews related work, and 
\Cref{sec:research_questions} formalizes the research questions.
\Cref{sec:rq1_evaluation}, \Cref{sec:rq2_evaluation}, and \Cref{sec:rq3_evaluation} each present the methodology and empirical findings for one research question.
The paper concludes in \Cref{sec:discussion} with a discussion of implications for the software engineering community and future work, followed by an examination of potential threats to validity in \Cref{sec:threats} and closing remarks in \Cref{sec:conclusion}.

% =========================================================================
% SECTION: BACKGROUND AND RELATED WORK
% Narrative Arc: "What world are we in?" -> "What fails here?" -> "How do we fix it?"
% =========================================================================
\section{Background and Related Work}
\label{sec:background}

% --- MAJOR SIGNPOSTING & RE-CONTEXTUALIZATION ---
% Entry paragraph (3–5 sentences): re‑state the problem space with more technical depth than the Introduction and briefly signpost the three subsections (domain overview, existing approaches, conceptual foundation).
Lorem ipsum dolor sit amet consectetur adipiscing elit. 
Quisque faucibus ex sapien vitae pellentesque sem placerat.
In id cursus mi pretium tellus duis convallis. 
Tempus leo eu aenean sed diam urna tempor.
Pulvinar vivamus fringilla lacus nec metus bibendum egestas. 

% =========================================================================
% SUBSECTION 1: THE DOMAIN OVERVIEW (Domain overview)
% =========================================================================
\subsection{Subsection 1}
\label{subsec:subsection-1}

% Paragraph 1 – Domain definitions and workflows: 
% Technical definitions and standard workflows. Clarify entities (projects, pipelines, registries), typical processes, and key terminology.
Lorem ipsum dolor sit amet consectetur adipiscing elit. 
Quisque faucibus ex sapien vitae pellentesque sem placerat.
In id cursus mi pretium tellus duis convallis. 
Tempus leo eu aenean sed diam urna tempor.
Pulvinar vivamus fringilla lacus nec metus bibendum egestas. 
Iaculis massa nisl malesuada lacinia integer nunc posuere.

% Paragraph 2 – Cause‑and‑effect linkages: 
% Explain how practices and constraints in this domain lead to specific challenges (e.g., fragility, reproducibility issues). Use ‘As a result/Consequently’ to make causal links explicit.
Lorem ipsum dolor sit amet consectetur adipiscing elit.
Quisque faucibus ex sapien vitae pellentesque sem placerat. 
In id cursus mi pretium tellus duis convallis.
Tempus leo eu aenean sed diam urna tempor. 
Pulvinar vivamus fringilla lacus nec metus bibendum egestas.
Iaculis massa nisl malesuada lacinia integer nunc posuere. 

% Paragraph 3 – Transition to existing approaches: 
% Close the subsection by hinting at what practitioners or researchers currently do to cope with these challenges, setting up Subsection 2.
Lorem ipsum dolor sit amet consectetur adipiscing elit. 
Quisque faucibus ex sapien vitae pellentesque sem placerat.
In id cursus mi pretium tellus duis convallis. 
Tempus leo eu aenean sed diam urna tempor.
Pulvinar vivamus fringilla lacus nec metus bibendum egestas. 
Iaculis massa nisl malesuada lacinia integer nunc posuere.

% =========================================================================
% SUBSECTION 2: EXISTING APPROACHES & PRIOR WORK
% =========================================================================
\subsection{Subsection 2}
\label{subsec:subsection-2}

% Paragraph 1 – Subsection entry and literature grouping: 
% Target: 3–5 sentences. 
% Content: Explicitly tie back to Subsection 1 and describe high‑level categories of existing approaches (e.g., tooling, guidelines, empirical studies). Group literature by theme, not chronologically.
Lorem ipsum dolor sit amet consectetur adipiscing elit. 
Quisque faucibus ex sapien vitae pellentesque sem placerat.
In id cursus mi pretium tellus duis convallis. 
Tempus leo eu aenean sed diam urna tempor.
Pulvinar vivamus fringilla lacus nec metus bibendum egestas. 
Iaculis massa nisl malesuada lacinia integer nunc posuere.

% Paragraph 2 – Contrastive analysis and limitations: 
% Content: Describe what these approaches achieve and where they fall short. Use contrast markers (‘However’, ‘In contrast’) and connect limitations to your gap.
Lorem ipsum dolor sit amet consectetur adipiscing elit.
Quisque faucibus ex sapien vitae pellentesque sem placerat. 
In id cursus mi pretium tellus duis convallis.
Tempus leo eu aenean sed diam urna tempor. 
Pulvinar vivamus fringilla lacus nec metus bibendum egestas.
Iaculis massa nisl malesuada lacinia integer nunc posuere. 

% Paragraph 3 – Exit transition to conceptual lens: 
% Content: Summarize the overall shortcomings of existing work and point forward to the conceptual foundation you will use to reason about them (Subsection 3).
Lorem ipsum dolor sit amet consectetur adipiscing elit. 
Quisque faucibus ex sapien vitae pellentesque sem placerat.
In id cursus mi pretium tellus duis convallis. 
Tempus leo eu aenean sed diam urna tempor.
Pulvinar vivamus fringilla lacus nec metus bibendum egestas. 
Iaculis massa nisl malesuada lacinia integer nunc posuere.

% =========================================================================
% SUBSECTION 3: CONCEPTUAL FOUNDATION
% =========================================================================
\subsection{Subsection 3}
\label{subsec:subsection-3}

% Paragraph 1 – Subsection entry and theoretical framing: 
% Content: Introduce the conceptual/theoretical lens (e.g., socio‑technical, diffusion, reliability) that underpins your study. Motivate why this lens is appropriate given the domain and prior work.
Lorem ipsum dolor sit amet consectetur adipiscing elit. 
Quisque faucibus ex sapien vitae pellentesque sem placerat.
In id cursus mi pretium tellus duis convallis. 
Tempus leo eu aenean sed diam urna tempor.
Pulvinar vivamus fringilla lacus nec metus bibendum egestas. 
Iaculis massa nisl malesuada lacinia integer nunc posuere.

% Paragraph 2 – Technical hook into methodology: 
% Content: Tie the chosen lens directly to your upcoming methodology, hinting at key constructs, relationships, or hypotheses that will be operationalized in the Methods and RQ sections.
Lorem ipsum dolor sit amet consectetur adipiscing elit. 
Quisque faucibus ex sapien vitae pellentesque sem placerat.
In id cursus mi pretium tellus duis convallis. 
Tempus leo eu aenean sed diam urna tempor.
Pulvinar vivamus fringilla lacus nec metus bibendum egestas. 
Iaculis massa nisl malesuada lacinia integer nunc posuere.

% =========================================================================
% SECTION: RESEARCH QUESTIONS (The Structural Bridge)
% =========================================================================
\section{Research Questions}
\label{sec:research_questions}

% --- INTRODUCTORY SIGNPOSTING & GAP ANCHORING ---
% Target: 1–3 sentences. 
% Content: Remind the reader of the key gap from Background and state that you now formalize the evaluation goals as research questions.
Lorem ipsum dolor sit amet consectetur adipiscing elit. 
Quisque faucibus ex sapien vitae pellentesque sem placerat.
In id cursus mi pretium tellus duis convallis. 

% --- THE FORMAL RQ LIST ---
% Function: The formal, enumerated research questions (the structural bridge).
% Content: List 3–4 concise RQs, each isolating ONE dimension (e.g., prevalence,
%   patterns, impact). Keep each RQ to a single sentence prefaced by a bolded
%   dimensional anchor (e.g., [Prevalence]).
% Format: per the style guide, the colon sits OUTSIDE the bold: \textbf{RQ\arabic*}:
\noindent We address the following research questions:
\begin{enumerate}[label=\textbf{RQ\arabic*}:, leftmargin=*]
    \item \textbf{[Dimension 1, e.g., Prevalence].} Lorem ipsum dolor sit amet consectetur adipiscing elit?
    \item \textbf{[Dimension 2, e.g., Patterns].} Lorem ipsum dolor sit amet consectetur adipiscing elit?
    \item \textbf{[Dimension 3, e.g., Impact].} Lorem ipsum dolor sit amet consectetur adipiscing elit?
\end{enumerate}

% --- MAPPING PARAGRAPH ---
% Target: 2–3 sentences. 
% Content: Map each RQ to its evaluation section (RQ1, RQ2, RQ3) and briefly mention which data and methods answer which question. This proves the downstream structure aligns with stated goals.
Lorem ipsum dolor sit amet consectetur adipiscing elit. 
Quisque faucibus ex sapien vitae pellentesque sem placerat.
In id cursus mi pretium tellus duis convallis. 

% =========================================================================
% SECTION: EVALUATION OF RESEARCH QUESTION 1 (RQ1)
% RQ1 block arc: overview & motivation → RQ1‑specific methodology → RQ1 findings → formal answer box → pointer to Discussion.
% =========================================================================
\section{Evaluation of RQ1: [Dimension Title, e.g., Prevalence]}
\label{sec:rq1_evaluation}

% --- RQ1 OVERVIEW PARAGRAPH ---
% Target: 6–8 sentences. 
% Content: Restate RQ1 verbatim, explain its role relative to other RQs (baseline, patterns, impact), and describe why this dimension matters. End with a sentence that introduces the upcoming methodology subsection.
Lorem ipsum dolor sit amet consectetur adipiscing elit. 
Quisque faucibus ex sapien vitae pellentesque sem placerat. 
In id cursus mi pretium tellus duis convallis.
Tempus leo eu aenean sed diam urna tempor. 
Pulvinar vivamus fringilla lacus nec metus bibendum egestas.
Iaculis massa nisl malesuada lacinia integer nunc posuere. 
Ut hendrerit semper vel class aptent taciti sociosqu. 
Ad litora torquent per conubia nostra inceptos himenaeos.

% --- SUBSECTION: METHODOLOGY FOR RQ1 ---
\subsection{RQ1 Methodology and Operationalization}
\label{subsec:rq1_methodology}
% RULE: cite every figure in-text via \Cref{fig:...}; never rely on placement alone.
\input{figures/template_figure}

% Paragraph 1 – Design & rationale: 
% Explain the study design (experiment, case study, mining, survey) and why it is appropriate for RQ1.
Lorem ipsum dolor sit amet consectetur adipiscing elit. 
Quisque faucibus ex sapien vitae pellentesque sem placerat. 
In id cursus mi pretium tellus duis convallis. 
Tempus leo eu aenean sed diam urna tempor.

% Paragraph 2 – Data sources & sampling: 
% Describe datasets, projects, participants, and inclusion/exclusion rules. Provide enough detail for replication.
Lorem ipsum dolor sit amet consectetur adipiscing elit. 
Quisque faucibus ex sapien vitae pellentesque sem placerat. 
In id cursus mi pretium tellus duis convallis. 
Tempus leo eu aenean sed diam urna tempor.

% Paragraph 3 – Procedure: 
% Outline the step‑by‑step execution (pipelines run, tasks assigned, scripts used), emphasizing chronology and reproducibility.
Lorem ipsum dolor sit amet consectetur adipiscing elit. 
Quisque faucibus ex sapien vitae pellentesque sem placerat. 
In id cursus mi pretium tellus duis convallis. 
Tempus leo eu aenean sed diam urna tempor.

% Paragraph 4 – Measures & analysis plan: 
% Enumerate metrics and constructs, describe how they are calculated, and specify statistical/analytical methods used to answer RQ1.
Lorem ipsum dolor sit amet consectetur adipiscing elit. 
Quisque faucibus ex sapien vitae pellentesque sem placerat. 
In id cursus mi pretium tellus duis convallis. 
Tempus leo eu aenean sed diam urna tempor.

% Threats to Validity for RQ1:
% Short paragraph. Note key RQ1‑specific threats (e.g., confounding, sampling, measurement) and point to the global Threats section for deeper analysis.
\myparagraph{Threats to Validity for RQ1}
Lorem ipsum dolor sit amet consectetur adipiscing elit. 
Quisque faucibus ex sapien vitae pellentesque sem placerat. 
In id cursus mi pretium tellus duis convallis. 

% --- SUBSECTION: EMPIRICAL FINDINGS FOR RQ1 ---
\subsection{RQ1 Empirical Findings}
\label{subsec:rq1_findings}

% Paragraph 1 – Overview and preprocessing: 
% Describe the subset of data used for RQ1, any preprocessing, exclusions, and basic descriptive stats (e.g., counts, ranges).
Lorem ipsum dolor sit amet consectetur adipiscing elit. 
Quisque faucibus ex sapien vitae pellentesque sem placerat. 
In id cursus mi pretium tellus duis convallis. 
Tempus leo eu aenean sed diam urna tempor.

% Paragraphs 2–3 – Main findings with figures/tables: 
% Present the core results in logical order (e.g., primary metrics, distributions, comparative stats). Every figure/table referenced must be explicitly cited in the text (‘As shown in Figure X…’). No interpretation beyond factual patterns.
Lorem ipsum dolor sit amet consectetur adipiscing elit. 
Quisque faucibus ex sapien vitae pellentesque sem placerat.
In id cursus mi pretium tellus duis convallis. 
Tempus leo eu aenean sed diam urna tempor.
Pulvinar vivamus fringilla lacus nec metus bibendum egestas. 
Iaculis massa nisl malesuada lacinia integer nunc posuere.
Ut hendrerit semper vel class aptent taciti sociosqu. 

% Paragraph 4 – Trends, patterns, and null results: 
% Summarize observed trends, highlight interesting edge cases, and explicitly mention negative or null results. Stay descriptive; save interpretation for the Discussion.
Lorem ipsum dolor sit amet consectetur adipiscing elit. 
Quisque faucibus ex sapien vitae pellentesque sem placerat. 
In id cursus mi pretium tellus duis convallis. 
Tempus leo eu aenean sed diam urna tempor.

% Formal ‘Answer to RQ1’ box: 
% 2–3 sentences of neutral synthesis directly answering the RQ text, with no speculation.
\input{figures/template_finding}

% Pointer sentence to Discussion: 
% Single sentence that forwards the reader to the Discussion for qualitative interpretation and implications.
Lorem ipsum dolor sit amet consectetur adipiscing elit.

% =========================================================================
% SECTION: EVALUATION OF RESEARCH QUESTION 2 (RQ2)
% RQ2 block arc: overview & motivation → RQ2‑specific methodology → RQ2 findings → formal answer box → pointer to Discussion.
% =========================================================================
\section{Evaluation of RQ2: [Dimension Title, e.g., Patterns]}
\label{sec:rq2_evaluation}

% --- RQ2 OVERVIEW PARAGRAPH ---
% Target: 6–8 sentences.
% Content: Restate RQ2 verbatim, explain its role relative to the other RQs (it
%   builds on RQ1's baseline by characterizing patterns/mechanics), and explain
%   why this dimension matters. End with a sentence introducing the methodology.
Lorem ipsum dolor sit amet consectetur adipiscing elit.
Quisque faucibus ex sapien vitae pellentesque sem placerat. 
In id cursus mi pretium tellus duis convallis.
Tempus leo eu aenean sed diam urna tempor. 
Pulvinar vivamus fringilla lacus nec metus bibendum egestas.
Iaculis massa nisl malesuada lacinia integer nunc posuere. 
Ut hendrerit semper vel class aptent taciti sociosqu. 
Ad litora torquent per conubia nostra inceptos himenaeos.

% --- SUBSECTION: METHODOLOGY FOR RQ2 ---
\subsection{RQ2 Methodology and Categorization Pipeline}
\label{subsec:rq2_methodology}
% RULE: cite every figure in-text via \Cref{fig:...}; never rely on placement alone.
\input{figures/template_figure}

% Paragraph 1 – Design & rationale:
% Explain the study design (experiment, case study, mining, survey) and why it
% is appropriate for RQ2.
Lorem ipsum dolor sit amet consectetur adipiscing elit.
Quisque faucibus ex sapien vitae pellentesque sem placerat. 
In id cursus mi pretium tellus duis convallis. 
Tempus leo eu aenean sed diam urna tempor.

% Paragraph 2 – Data sources & sampling:
% Describe datasets, projects, participants, and inclusion/exclusion rules.
% Provide enough detail for replication.
Lorem ipsum dolor sit amet consectetur adipiscing elit.
Quisque faucibus ex sapien vitae pellentesque sem placerat. 
In id cursus mi pretium tellus duis convallis. 
Tempus leo eu aenean sed diam urna tempor.

% Paragraph 3 – Procedure:
% Outline the step-by-step execution (pipelines run, tasks assigned, scripts
% used), emphasizing chronology and reproducibility.
Lorem ipsum dolor sit amet consectetur adipiscing elit.
Quisque faucibus ex sapien vitae pellentesque sem placerat. 
In id cursus mi pretium tellus duis convallis. 
Tempus leo eu aenean sed diam urna tempor.

% Paragraph 4 – Measures & analysis plan:
% Enumerate metrics and constructs, describe how they are calculated, and specify
% the statistical/analytical methods used to answer RQ2.
Lorem ipsum dolor sit amet consectetur adipiscing elit.
Quisque faucibus ex sapien vitae pellentesque sem placerat. 
In id cursus mi pretium tellus duis convallis. 
Tempus leo eu aenean sed diam urna tempor.

% Threats to Validity for RQ2:
% Short paragraph (2–3 sentences). Note key RQ2-specific threats (e.g.,
% confounding, sampling, measurement) and point to the global Threats section.
\myparagraph{Threats to Validity for RQ2}
Lorem ipsum dolor sit amet consectetur adipiscing elit.
Quisque faucibus ex sapien vitae pellentesque sem placerat. 
In id cursus mi pretium tellus duis convallis. 

% --- SUBSECTION: EMPIRICAL FINDINGS FOR RQ2 ---
\subsection{RQ2 Empirical Findings}
\label{subsec:rq2_findings}

% Paragraph 1 – Overview and preprocessing:
% Describe the subset of data used for RQ2, any preprocessing, exclusions, and
% basic descriptive stats (e.g., counts, ranges).
Lorem ipsum dolor sit amet consectetur adipiscing elit.
Quisque faucibus ex sapien vitae pellentesque sem placerat. 
In id cursus mi pretium tellus duis convallis. 
Tempus leo eu aenean sed diam urna tempor.

% Paragraphs 2–3 – Main findings with figures/tables:
% Present the core results in logical order. Every figure/table referenced must
% be explicitly cited in the text ("As shown in \Cref{fig:...}"). Stay strictly
% descriptive; no interpretation beyond factual patterns.
Lorem ipsum dolor sit amet consectetur adipiscing elit. 
Quisque faucibus ex sapien vitae pellentesque sem placerat.
In id cursus mi pretium tellus duis convallis. 
Tempus leo eu aenean sed diam urna tempor.
Pulvinar vivamus fringilla lacus nec metus bibendum egestas. 
Iaculis massa nisl malesuada lacinia integer nunc posuere.
Ut hendrerit semper vel class aptent taciti sociosqu. 

% Paragraph 4 – Trends, patterns, and null results:
% Summarize observed trends, highlight interesting edge cases, and explicitly
% mention negative or null results. Stay descriptive; save interpretation for the
% Discussion.
Lorem ipsum dolor sit amet consectetur adipiscing elit.
Quisque faucibus ex sapien vitae pellentesque sem placerat. 
In id cursus mi pretium tellus duis convallis. 
Tempus leo eu aenean sed diam urna tempor. 

% Formal 'Answer to RQ2' box:
% 2–3 sentences of neutral synthesis directly answering the RQ text, no speculation.
\input{figures/template_finding}

% Pointer sentence to Discussion:
% Single sentence forwarding the reader to the Discussion for qualitative
% interpretation and implications.
Lorem ipsum dolor sit amet consectetur adipiscing elit.

% =========================================================================
% SECTION: EVALUATION OF RESEARCH QUESTION 3 (RQ3)
% RQ3 block arc: overview & motivation → RQ3‑specific methodology → RQ3 findings → formal answer box → pointer to Discussion.
% =========================================================================
\section{Evaluation of RQ3: [Dimension Title, e.g., Impact]}
\label{sec:rq3_evaluation}

% --- RQ3 OVERVIEW PARAGRAPH ---
% Target: 6–8 sentences.
% Content: Restate RQ3 verbatim, explain its role relative to the other RQs (it
%   typically caps the arc by assessing impact/consequence), and explain why this
%   dimension matters. End with a sentence introducing the methodology subsection.
Lorem ipsum dolor sit amet consectetur adipiscing elit.
Quisque faucibus ex sapien vitae pellentesque sem placerat. 
In id cursus mi pretium tellus duis convallis.
Tempus leo eu aenean sed diam urna tempor. 
Pulvinar vivamus fringilla lacus nec metus bibendum egestas.
Iaculis massa nisl malesuada lacinia integer nunc posuere. 
Ut hendrerit semper vel class aptent taciti sociosqu. 
Ad litora torquent per conubia nostra inceptos himenaeos.

% --- SUBSECTION: METHODOLOGY FOR RQ3 ---
\subsection{RQ3 Methodology and Analytical Matrix}
\label{subsec:rq3_methodology}
% RULE: cite every figure in-text via \Cref{fig:...}; never rely on placement alone.
\input{figures/template_figure}

% Paragraph 1 – Design & rationale:
% Explain the study design (experiment, case study, mining, survey) and why it
% is appropriate for RQ3.
Lorem ipsum dolor sit amet consectetur adipiscing elit.
Quisque faucibus ex sapien vitae pellentesque sem placerat. 
In id cursus mi pretium tellus duis convallis. 
Tempus leo eu aenean sed diam urna tempor.

% Paragraph 2 – Data sources & sampling:
% Describe datasets, projects, participants, and inclusion/exclusion rules.
% Provide enough detail for replication.
Lorem ipsum dolor sit amet consectetur adipiscing elit.
Quisque faucibus ex sapien vitae pellentesque sem placerat. 
In id cursus mi pretium tellus duis convallis. 
Tempus leo eu aenean sed diam urna tempor. 

% Paragraph 3 – Procedure:
% Outline the step-by-step execution (pipelines run, tasks assigned, scripts
% used), emphasizing chronology and reproducibility.
Lorem ipsum dolor sit amet consectetur adipiscing elit.
Quisque faucibus ex sapien vitae pellentesque sem placerat. 
In id cursus mi pretium tellus duis convallis. 
Tempus leo eu aenean sed diam urna tempor.

% Paragraph 4 – Measures & analysis plan:
% Enumerate metrics and constructs, describe how they are calculated, and specify
% the statistical/analytical methods used to answer RQ3.
Lorem ipsum dolor sit amet consectetur adipiscing elit.
Quisque faucibus ex sapien vitae pellentesque sem placerat. 
In id cursus mi pretium tellus duis convallis. 
Tempus leo eu aenean sed diam urna tempor.

% Threats to Validity for RQ3:
% Short paragraph (2–3 sentences). Note key RQ3-specific threats (e.g.,
% confounding, sampling, measurement) and point to the global Threats section.
\myparagraph{Threats to Validity for RQ3}
Lorem ipsum dolor sit amet consectetur adipiscing elit.
Quisque faucibus ex sapien vitae pellentesque sem placerat. 
In id cursus mi pretium tellus duis convallis. 

% --- SUBSECTION: EMPIRICAL FINDINGS FOR RQ3 ---
\subsection{RQ3 Empirical Findings}
\label{subsec:rq3_findings}

% Paragraph 1 – Overview and preprocessing:
% Describe the subset of data used for RQ3, any preprocessing, exclusions, and
% basic descriptive stats (e.g., counts, ranges).
Lorem ipsum dolor sit amet consectetur adipiscing elit.
Quisque faucibus ex sapien vitae pellentesque sem placerat. 
In id cursus mi pretium tellus duis convallis. 
Tempus leo eu aenean sed diam urna tempor.

% Paragraphs 2–3 – Main findings with figures/tables:
% Present the core results in logical order. Every figure/table referenced must
% be explicitly cited in the text ("As shown in \Cref{fig:...}"). Stay strictly
% descriptive; no interpretation beyond factual patterns.
Lorem ipsum dolor sit amet consectetur adipiscing elit. 
Quisque faucibus ex sapien vitae pellentesque sem placerat.
In id cursus mi pretium tellus duis convallis. 
Tempus leo eu aenean sed diam urna tempor.
Pulvinar vivamus fringilla lacus nec metus bibendum egestas. 
Iaculis massa nisl malesuada lacinia integer nunc posuere.
Ut hendrerit semper vel class aptent taciti sociosqu. 

% Paragraph 4 – Trends, patterns, and null results:
% Summarize observed trends, highlight interesting edge cases, and explicitly
% mention negative or null results. Stay descriptive; save interpretation for the
% Discussion.
Lorem ipsum dolor sit amet consectetur adipiscing elit.
Quisque faucibus ex sapien vitae pellentesque sem placerat. 
In id cursus mi pretium tellus duis convallis. 
Tempus leo eu aenean sed diam urna tempor. 

% Formal 'Answer to RQ3' box:
% 2–3 sentences of neutral synthesis directly answering the RQ text, no speculation.
\input{figures/template_finding}

% Pointer sentence to Discussion:
% Single sentence forwarding the reader to the Discussion for qualitative
% interpretation and implications.
Lorem ipsum dolor sit amet consectetur adipiscing elit.
Quisque faucibus ex sapien vitae pellentesque sem placerat. 
In id cursus mi pretium tellus duis convallis. 
Tempus leo eu aenean sed diam urna tempor. 

% =========================================================================
% SECTION: DISCUSSION
% Narrative Flow: Narrow (Interpretations) -> Broad (Implications) 
%                -> Defensive (Limitations) -> Forward-Looking (Future Work)
% =========================================================================
\section{Discussion}
\label{sec:discussion}

% --- PARAGRAPH 1: DISCUSSION OPENING ---
% Target: 7–9 sentences. 
% Content: Restate the high‑level takeaway of the study, tying back to the Introduction’s big gap and summarizing the overall pattern across RQ1–RQ3. This paragraph sets the tone for the rest of the Discussion.
Lorem ipsum dolor sit amet consectetur adipiscing elit. 
Quisque faucibus ex sapien vitae pellentesque sem placerat. 
In id cursus mi pretium tellus duis convallis. 
Tempus leo eu aenean sed diam urna tempor. 
Pulvinar vivamus fringilla lacus nec metus bibendum egestas. 
Iaculis massa nisl malesuada lacinia integer nunc posuere. 
Ut hendrerit semper vel class aptent taciti sociosqu. 

% --- SUBSECTION: INTERPRETATION OF RQS ---
% Per‑RQ mini‑discussions: each \myparagraph is 8–12 sentences total, structured as interpretation → mechanism → relation to prior work → scoped implication/limit.
\subsection{Interpretation of Research Questions}
\label{subsec:interpretations}

\myparagraph{RQ$_1$ -- Baseline Prevalence Metrics}
% Interpretation: 2–3 sentences summarizing how the findings answer RQ1.
Lorem ipsum dolor sit amet consectetur adipiscing elit. Quisque faucibus ex sapien vitae pellentesque sem placerat. In id cursus mi pretium tellus duis convallis.
% Mechanism/Why: 2–3 sentences explaining why these patterns likely occur, grounded in theory or domain knowledge.
Lorem ipsum dolor sit amet consectetur adipiscing elit. Quisque faucibus ex sapien vitae pellentesque sem placerat. In id cursus mi pretium tellus duis convallis.
% Relation to prior work: 2–3 sentences comparing RQ1 findings to specific prior studies; note agreements and contradictions.
Lorem ipsum dolor sit amet consectetur adipiscing elit. Quisque faucibus ex sapien vitae pellentesque sem placerat. In id cursus mi pretium tellus duis convallis.
% Scoped actionable/limit: 2–3 sentences stating implications or recommendations for practice, plus one key limitation that tempers those recommendations.
Lorem ipsum dolor sit amet consectetur adipiscing elit. Quisque faucibus ex sapien vitae pellentesque sem placerat. In id cursus mi pretium tellus duis convallis. Tempus leo eu aenean sed diam urna tempor.

\myparagraph{RQ$_2$ -- The Mechanics of Fragile Reuse}
% Interpretation: 2–3 sentences summarizing how the findings answer RQ2.
Lorem ipsum dolor sit amet consectetur adipiscing elit.
% Mechanism/Why: 2–3 sentences explaining why these patterns likely occur, grounded in theory or domain knowledge.
Lorem ipsum dolor sit amet consectetur adipiscing elit. Quisque faucibus ex sapien vitae pellentesque sem placerat. In id cursus mi pretium tellus duis convallis.
% Relation to prior work: 2–3 sentences comparing RQ2 findings to specific prior studies.
Lorem ipsum dolor sit amet consectetur adipiscing elit. Quisque faucibus ex sapien vitae pellentesque sem placerat. In id cursus mi pretium tellus duis convallis.
% Scoped actionable/limit: 2–3 sentences stating implications or recommendations for practice.
Lorem ipsum dolor sit amet consectetur adipiscing elit. Quisque faucibus ex sapien vitae pellentesque sem placerat. In id cursus mi pretium tellus duis convallis. Tempus leo eu aenean sed diam urna tempor.

\myparagraph{RQ$_3$ -- Quantification of Pipeline Half-Life}
% Interpretation: 2–3 sentences summarizing how the findings answer RQ3.
Lorem ipsum dolor sit amet consectetur adipiscing elit.
% Mechanism/Why: 2–3 sentences explaining why these patterns likely occur, grounded in theory or domain knowledge.
Lorem ipsum dolor sit amet consectetur adipiscing elit. Quisque faucibus ex sapien vitae pellentesque sem placerat. In id cursus mi pretium tellus duis convallis.
% Relation to prior work: 2–3 sentences comparing RQ3 findings to specific prior studies.
Lorem ipsum dolor sit amet consectetur adipiscing elit. Quisque faucibus ex sapien vitae pellentesque sem placerat. In id cursus mi pretium tellus duis convallis.
% Scoped actionable/limit: 2–3 sentences stating implications or recommendations for practice.
Lorem ipsum dolor sit amet consectetur adipiscing elit. Quisque faucibus ex sapien vitae pellentesque sem placerat. In id cursus mi pretium tellus duis convallis. Tempus leo eu aenean sed diam urna tempor.

% --- SUBSECTION: WIDER IMPLICATIONS ---
\subsection{Wider Implications for Science and Engineering}
\label{subsec:implications}

% Paragraph 1 – Theoretical implications: 
% Target: 5–7 sentences. 
% Content: Synthesize across RQs to state how the study refines or challenges existing theories/models. Focus on general principles.
Lorem ipsum dolor sit amet consectetur adipiscing elit. 
Quisque faucibus ex sapien vitae pellentesque sem placerat. 
In id cursus mi pretium tellus duis convallis. 
Tempus leo eu aenean sed diam urna tempor. 
Pulvinar vivamus fringilla lacus nec metus bibendum egestas. 
Iaculis massa nisl malesuada lacinia integer nunc posuere. 
Ut hendrerit semper vel class aptent taciti sociosqu.

% Paragraph 2 – Practical/industrial implications: 
% Target: 5–7 sentences. 
% Content: Translate findings into concrete recommendations for teams, tools, and processes, mentioning feasibility and constraints.
Lorem ipsum dolor sit amet consectetur adipiscing elit. 
Quisque faucibus ex sapien vitae pellentesque sem placerat. 
In id cursus mi pretium tellus duis convallis. 
Tempus leo eu aenean sed diam urna tempor. 
Pulvinar vivamus fringilla lacus nec metus bibendum egestas. 
Iaculis massa nisl malesuada lacinia integer nunc posuere. 
Ut hendrerit semper vel class aptent taciti sociosqu.

% Optional 3rd Paragraph – Policy / Educational Frameworks:
% NOTE: Optional, beyond the two standard implication types (Theoretical and
%   Practical/Industrial) defined by the style guide. Delete if not needed.
% Target: 5–7 sentences on implications for standards, curricula, or broader scientific practice.
Lorem ipsum dolor sit amet consectetur adipiscing elit. 
Quisque faucibus ex sapien vitae pellentesque sem placerat. 
In id cursus mi pretium tellus duis convallis. 
Tempus leo eu aenean sed diam urna tempor. 
Pulvinar vivamus fringilla lacus nec metus bibendum egestas. 

% --- SUBSECTION: LIMITATIONS AND BIASES ---
\subsection{Limitations and Alternative Explanations}
\label{subsec:limitations}

% Paragraph 1 – Main limitations: 
% Target: 5–7 sentences. 
% Content: Summarize key limitations across design, data, and measures; explain how each might bias results or limit generality.
Lorem ipsum dolor sit amet consectetur adipiscing elit. 
Quisque faucibus ex sapien vitae pellentesque sem placerat. 
In id cursus mi pretium tellus duis convallis. 
Tempus leo eu aenean sed diam urna tempor. 
Pulvinar vivamus fringilla lacus nec metus bibendum egestas. 

% Paragraph 2 – Alternative explanations: 
% Target: 5–7 sentences. 
% Content: Present plausible alternative interpretations of major findings and briefly note what further data/methods would be needed to test them.
Lorem ipsum dolor sit amet consectetur adipiscing elit. 
Quisque faucibus ex sapien vitae pellentesque sem placerat. 
In id cursus mi pretium tellus duis convallis. 
Tempus leo eu aenean sed diam urna tempor. 
Pulvinar vivamus fringilla lacus nec metus bibendum egestas. 

% --- SUBSECTION: FUTURE DIRECTIONS ---
\subsection{Future Research Trajectories}
\label{subsec:future_directions}

% Paragraph 1 – Next logical steps: 
% Target: 4–6 sentences. 
% Content: Propose immediate follow‑up studies (replications, scaling, boundary conditions), each explicitly linked to a specific result or limitation.
Lorem ipsum dolor sit amet consectetur adipiscing elit. 
Quisque faucibus ex sapien vitae pellentesque sem placerat. 
In id cursus mi pretium tellus duis convallis. 
Tempus leo eu aenean sed diam urna tempor. 

% Paragraph 2 – Methodological refinements: 
% Target: 3–5 sentences. 
% Content: Suggest improvements to design, metrics, or data collection that would strengthen future evidence.
Lorem ipsum dolor sit amet consectetur adipiscing elit. 
Quisque faucibus ex sapien vitae pellentesque sem placerat. 
In id cursus mi pretium tellus duis convallis. 

% Paragraph 3 – New questions: 
% Target: 4–6 sentences. 
% Content: Outline new research questions opened by unexpected findings or implications; explain why each is important for empirical SE.
Lorem ipsum dolor sit amet consectetur adipiscing elit. 
Quisque faucibus ex sapien vitae pellentesque sem placerat. 
In id cursus mi pretium tellus duis convallis. 
Tempus leo eu aenean sed diam urna tempor. 

% Final closing block – Take-Home Message:
% Target: 4–6 sentences. 
% Content: Provide a synthesized, confident statement of what the study ultimately shows and why it matters. No new ideas or data; this should echo the core contribution and high‑level impact.
\myparagraph{Take-Home Message}
Lorem ipsum dolor sit amet consectetur adipiscing elit. 
Quisque faucibus ex sapien vitae pellentesque sem placerat. 
In id cursus mi pretium tellus duis convallis. 
Tempus leo eu aenean sed diam urna tempor. 
Pulvinar vivamus fringilla lacus nec metus bibendum egestas. 
Iaculis massa nisl malesuada lacinia integer nunc posuere. 

% =========================================================================
% SECTION: THREATS TO VALIDITY
% Function: Systematically classify study vulnerabilities into three subsections
%   (Internal, External, Construct validity). Be candid; each block names a threat,
%   its potential effect on results, and the mitigation taken.
% =========================================================================
\section{Threats to Validity}
\label{sec:threats}

% Section opening:
% Target: 3–4 sentences. 
% Content: Explain the purpose of this section (systematically analyzing internal, external, and construct validity threats) and note that detailed RQ‑specific threats were mentioned in each RQ methodology subsection.
Lorem ipsum dolor sit amet consectetur adipiscing elit. 
Quisque faucibus ex sapien vitae pellentesque sem placerat. 
In id cursus mi pretium tellus duis convallis. 
Our analysis follows the standard empirical software engineering classification distinguishing between internal, external, and construct validity.

% --- SUBSECTION: INTERNAL VALIDITY ---
\subsection{Internal Validity}
\label{subsec:internal_validity}

% First block (Confounding Variables):
% Target: 3–5 sentences.
% Content: Explain potential confounding variables and how they may distort causal
%   interpretations; name the mitigation strategies used. Use causal signposts
%   ("As a result", "Consequently").
Lorem ipsum dolor sit amet consectetur adipiscing elit. 
Quisque faucibus ex sapien vitae pellentesque sem placerat. 
In id cursus mi pretium tellus duis convallis. 
Tempus leo eu aenean sed diam urna tempor. 
Pulvinar vivamus fringilla lacus nec metus bibendum egestas. 

% Second block (History, Maturation, and Instrumentation):
% Target: 3–4 sentences.
% Content: Discuss history, maturation, and instrumentation effects and their
%   potential influence on outcomes; note how design choices attempt to control them.
Lorem ipsum dolor sit amet consectetur adipiscing elit. 
Quisque faucibus ex sapien vitae pellentesque sem placerat. 
In id cursus mi pretium tellus duis convallis. 
Tempus leo eu aenean sed diam urna tempor. 

% Third block (Attrition and Missing Data):
% Target: 2–3 sentences.
% Content: Describe attrition/missing data patterns and how they were handled
%   (e.g., exclusion, imputation, sensitivity analysis).
Lorem ipsum dolor sit amet consectetur adipiscing elit. 
Quisque faucibus ex sapien vitae pellentesque sem placerat. 
In id cursus mi pretium tellus duis convallis. 

% --- SUBSECTION: EXTERNAL VALIDITY ---
\subsection{External Validity}
\label{subsec:external_validity}

% First block (Context & Generalizability):
% Target: 3–4 sentences.
% Content: Clarify the study setting and domains, and explicitly state the boundary
%   conditions — where the results do and do not generalize.
Lorem ipsum dolor sit amet consectetur adipiscing elit. 
Quisque faucibus ex sapien vitae pellentesque sem placerat. 
In id cursus mi pretium tellus duis convallis. 
Tempus leo eu aenean sed diam urna tempor. 

% Second block (Sampling Bias):
% Target: 2–3 sentences.
% Content: Discuss how selection of projects/participants may skew results and how
%   that limits external validity.
Lorem ipsum dolor sit amet consectetur adipiscing elit. 
Quisque faucibus ex sapien vitae pellentesque sem placerat. 
In id cursus mi pretium tellus duis convallis. 

% --- SUBSECTION: CONSTRUCT VALIDITY ---
\subsection{Construct Validity}
\label{subsec:construct_validity}

% First block (Indirect Metrics and Operationalization):
% Target: 4–6 sentences.
% Content: Justify that the chosen metrics capture the intended constructs,
%   referencing prior work and triangulation where possible.
Lorem ipsum dolor sit amet consectetur adipiscing elit. 
Quisque faucibus ex sapien vitae pellentesque sem placerat. 
In id cursus mi pretium tellus duis convallis. 
Tempus leo eu aenean sed diam urna tempor. 
Pulvinar vivamus fringilla lacus nec metus bibendum egestas. 
Iaculis massa nisl malesuada lacinia integer nunc posuere. 
Ut hendrerit semper vel class aptent taciti sociosqu. 
Ad litora torquent per conubia nostra inceptos himenaeos.

% Second block (Measurement Error and Technical Bias):
% Target: 3–5 sentences.
% Content: Address logging inaccuracies, self-report biases, and tool limitations;
%   explain how they might affect interpretations.
Lorem ipsum dolor sit amet consectetur adipiscing elit. 
Quisque faucibus ex sapien vitae pellentesque sem placerat. 
In id cursus mi pretium tellus duis convallis. 
Tempus leo eu aenean sed diam urna tempor. 
Pulvinar vivamus fringilla lacus nec metus bibendum egestas. 
Iaculis massa nisl malesuada lacinia integer nunc posuere. 

% =========================================================================
% SECTION: CONCLUSION
% =========================================================================
\section{Conclusion}
\label{sec:conclusion}

% Paragraph 1 – Aims and findings synthesis: 
% Target: 2–4 sentences. 
% Content: Restate the overall problem and aims in fresh language, then synthesize the main findings across RQs (no detailed numbers).
Lorem ipsum dolor sit amet consectetur adipiscing elit. 
Quisque faucibus ex sapien vitae pellentesque sem placerat. 
In id cursus mi pretium tellus duis convallis. 
Tempus leo eu aenean sed diam urna tempor. 
Pulvinar vivamus fringilla lacus nec metus bibendum egestas. 
Iaculis massa nisl malesuada lacinia integer nunc posuere. 
Ut hendrerit semper vel class aptent taciti sociosqu.

% Paragraph 2 – Contributions and practical implications: 
% Target: 2–4 sentences. 
% Content: Explicitly name the key contributions and what they enable or change for researchers and practitioners.
Lorem ipsum dolor sit amet consectetur adipiscing elit. 
Quisque faucibus ex sapien vitae pellentesque sem placerat. 
In id cursus mi pretium tellus duis convallis. 
Tempus leo eu aenean sed diam urna tempor. 
Pulvinar vivamus fringilla lacus nec metus bibendum egestas. 
Iaculis massa nisl malesuada lacinia integer nunc posuere. 
Ut hendrerit semper vel class aptent taciti sociosqu.

% Paragraph 3 – Limitations, future work, final takeaway: 
% Target: 3–4 sentences. 
% Style Rule: The final sentence must be a confident, synthesized statement—not an apology.
% Content: Echo the most important limitation/future direction at a high level and end with a strong single sentence that captures the study’s core message.
Lorem ipsum dolor sit amet consectetur adipiscing elit. 
Quisque faucibus ex sapien vitae pellentesque sem placerat. 
In id cursus mi pretium tellus duis convallis. 
Tempus leo eu aenean sed diam urna tempor. 
Pulvinar vivamus fringilla lacus nec metus bibendum egestas. 
Iaculis massa nisl malesuada lacinia integer nunc posuere. 
Ut hendrerit semper vel class aptent taciti sociosqu. 

% =========================================================================
% ACKNOWLEDGMENTS
% =========================================================================
\section*{Acknowledgments}
% Target: 1–3 sentences. 
% Content: Thank funding sources, collaborators, and infrastructure. No technical content.
Lorem ipsum dolor sit amet consectetur adipiscing elit. 
Quisque faucibus ex sapien vitae pellentesque sem placerat. 
In id cursus mi pretium tellus duis convallis. 

\bibliographystyle{ACM-Reference-Format}
\bibliography{citations}    % For all other submissions
%\input{main.bbl}          % For arXiv submission

\end{document}          % For arXiv submission

\end{document}          % For arXiv submission

\end{document}          % For arXiv submission

\end{document}